% Le condensateur : charge, décharge et énergie — Physique Tle D — Cours signé OnBuch+
% Programme officiel MINESEC. Compiler avec : tectonic P10-condensateur.tex
\newcommand{\DOCMATIERE}{Physique}
\newcommand{\DOCNIVEAU}{Tle D}
\newcommand{\DOCMODULE}{Module 3 : Électricité — évolutions temporelles des circuits électriques et électroniques}
\newcommand{\DOCLECON}{P10}
\newcommand{\DOCTITRE}{Le condensateur : charge, décharge et énergie}
% OnBuch+ — préambule « cours signés » ENRICHI (compilé avec tectonic / XeLaTeX)
% Système visuel « Pop » d'OnBuch+ : fond crème (jamais blanc), bordures encre
% épaisses, ombres « dures » décalées, titres Archivo Black en capitales.
% Version MATHÉMATIQUES : graphes (pgfplots/tikz), boîtes pédagogiques
% (définition, propriété, méthode, attention, le savais-tu, exercice résolu,
% exercice, TP, bilan), page de garde et signature OnBuch+ ancrée en bas.
%
% Macros attendues dans main.tex AVANT \input{preamble} :
%   \DOCMATIERE  \DOCNIVEAU  \DOCMODULE  \DOCLECON (numéro)  \DOCTITRE
\documentclass[11pt]{article}

\usepackage{fontspec}
\usepackage[french]{babel}
\usepackage[a4paper,margin=2cm,top=3.4cm,bottom=2.8cm,headheight=1.4cm,footskip=1.3cm]{geometry}
\usepackage{amsmath,amssymb,amsfonts}
\usepackage{mathtools} % \xmapsto, \coloneqq... souvent utilisés par les modèles
\usepackage{cancel} % simplifications barrées (bilans de forces)
% \abs{z} (module, valeur absolue) et \norm{u} (norme), utilisés par les modèles
\providecommand{\abs}{}\let\abs\relax\DeclarePairedDelimiter{\abs}{\lvert}{\rvert}
\providecommand{\norm}{}\let\norm\relax\DeclarePairedDelimiter{\norm}{\lVert}{\rVert}
\usepackage[table]{xcolor} % [table] : fournit aussi \rowcolors (alternance de lignes)
\usepackage{pagecolor}
\usepackage{tikz}
\usetikzlibrary{arrows.meta,positioning,calc,shapes.geometric,shapes.misc,shapes.arrows,decorations.pathmorphing,decorations.markings,patterns,shadows,mindmap,trees,fit,backgrounds,matrix,intersections}
\usepackage{pgfplots}
\usepackage[version=4]{mhchem} % équations nucléaires \ce{^{235}_{92}U}
\usepackage[european,siunitx]{circuitikz} % schémas électriques
\pgfplotsset{compat=1.18}
\usepgfplotslibrary{fillbetween} % aires entre courbes (calcul intégral)
\usepackage{enumitem}
\usepackage{fancyhdr}
% [most] charge toutes les bibliothèques tcolorbox (listings, minted, raster,
% documentation...) et gonfle inutilement la mémoire TeX sur un document long
% (~300 boîtes) : on ne charge que ce qui est réellement utilisé.
\usepackage[skins,breakable]{tcolorbox}
\usepackage{graphicx}
\usepackage{colortbl}
\usepackage{booktabs}
\usepackage{tabularx}
\usepackage{diagbox} % case d'en-tête diagonale des tableaux à double entrée
% Colonnes extensibles centrée / gauche / droite (C, L, R), souvent utilisées par les modèles
\newcolumntype{C}{>{\centering\arraybackslash}X}
\newcolumntype{L}{>{\raggedright\arraybackslash}X}
\newcolumntype{R}{>{\raggedleft\arraybackslash}X}
\usepackage{array}
\usepackage{multicol}
\usepackage{siunitx}
% parse-numbers=false : accepte aussi \SI{2,5 \times 10^{-3}}{...}
\sisetup{locale=FR,output-decimal-marker={,},group-minimum-digits=4,parse-numbers=false}
\usepackage{qrcode}
% \degree : commande fréquemment utilisée par les modèles pour un angle ou une
% température en dehors de \SI{}{} (non fournie par les paquets chargés).
\providecommand{\degree}{^{\circ}}
% \qed (fin de démonstration), utilisé par les modèles sans amsthm
\providecommand{\qed}{\hfill$\square$}
% \barre : double barre des valeurs interdites dans un tableau de variations (tkz-tab)
\providecommand{\barre}{\ensuremath{\|}}
% environnement proof (amsthm non chargé) utilisé par les modèles
\ifdefined\proof\else\newenvironment{proof}[1][Démonstration]{\par\noindent\textit{#1.}\ }{\qed\par}\fi
\usepackage{lastpage}
\usepackage{hyperref}
\hypersetup{hidelinks}

% ============================================================
% Palette « Pop » OnBuch+ — mêmes hex que la classe P (plus_ui.dart)
% ============================================================
\definecolor{poporange}{HTML}{F59321}
\definecolor{popdark}{HTML}{E07A0C}
\definecolor{popcream}{HTML}{FFF6E8}
\definecolor{popink}{HTML}{1C1714}
\definecolor{poppurple}{HTML}{7A5AE0}
\definecolor{popgreen}{HTML}{2E9E6A}
\definecolor{poppink}{HTML}{E0457B}
\definecolor{popblue}{HTML}{2F7DE1}
\definecolor{popgold}{HTML}{C98A12}
\definecolor{popmuted}{HTML}{8A7F76}
\definecolor{popline}{HTML}{EADFD2}
\definecolor{popgray}{HTML}{8A7F76}
\colorlet{popgrey}{popgray}
% Alias tolérants (le modèle nomme parfois les couleurs différemment)
\colorlet{popwhite}{white}
\colorlet{popblack}{popink}
\colorlet{popred}{poppink}
\colorlet{popyellow}{popgold}
% Teintes claires pour fonds de tableaux / zones
\colorlet{popblueL}{popblue!10!white}
\colorlet{poporangeL}{poporange!14!white}
\colorlet{popgreenL}{popgreen!12!white}
\colorlet{poppurpleL}{poppurple!12!white}
\colorlet{poppinkL}{poppink!12!white}
\colorlet{popgoldL}{popgold!14!white}

\pagecolor{popcream}

% ============================================================
% Typographie
% ============================================================
\defaultfontfeatures{Ligatures=TeX}
\setmainfont{PlusJakartaSans-Medium.ttf}[Path=../../../fonts/,BoldFont=PlusJakartaSans-Bold.ttf]
% Maths en sans-serif (Fira Math) avec chiffres et lettres droites de Plus
% Jakarta, pour que les formules se fondent dans le texte.
\usepackage[math-style=ISO,bold-style=ISO]{unicode-math}
\setmathfont{FiraMath-Regular.otf}[Path=../../../fonts/]
\setmathfont{PlusJakartaSans-Medium.ttf}[Path=../../../fonts/,range={up/num,up/{Latin,latin},\mathup}]
\usepackage{icomma} % virgule décimale sans espace en mode math
% Caractères mathématiques Unicode absents des polices de texte (Plus Jakarta,
% Archivo Black) : rendus par leur équivalent mathématique, en texte comme en
% formule — sinon ils s'affichent en carré vide (ex. « Arithmétique dans ℤ »).
\usepackage{newunicodechar}
\newunicodechar{ℝ}{\ensuremath{\mathbb{R}}}
\newunicodechar{ℤ}{\ensuremath{\mathbb{Z}}}
\newunicodechar{ℂ}{\ensuremath{\mathbb{C}}}
\newunicodechar{ℕ}{\ensuremath{\mathbb{N}}}
\newunicodechar{ℚ}{\ensuremath{\mathbb{Q}}}
\newunicodechar{𝔻}{\ensuremath{\mathbb{D}}}
\newunicodechar{α}{\ensuremath{\alpha}}
\newunicodechar{β}{\ensuremath{\beta}}
\newunicodechar{γ}{\ensuremath{\gamma}}
\newunicodechar{δ}{\ensuremath{\delta}}
\newunicodechar{ε}{\ensuremath{\varepsilon}}
\newunicodechar{θ}{\ensuremath{\theta}}
\newunicodechar{λ}{\ensuremath{\lambda}}
\newunicodechar{μ}{\ensuremath{\mu}}
\newunicodechar{σ}{\ensuremath{\sigma}}
\newunicodechar{τ}{\ensuremath{\tau}}
\newunicodechar{φ}{\ensuremath{\varphi}}
\newunicodechar{ψ}{\ensuremath{\psi}}
\newunicodechar{ω}{\ensuremath{\omega}}
\newunicodechar{Σ}{\ensuremath{\Sigma}}
% majuscules produites par \MakeUppercase dans les titres (α-aminés -> Α)
\newunicodechar{Α}{\ensuremath{\alpha}}
\newunicodechar{Β}{\ensuremath{\beta}}
\newunicodechar{Γ}{\ensuremath{\Gamma}}
\newunicodechar{Δ}{\ensuremath{\Delta}}
\newunicodechar{Ε}{\ensuremath{\varepsilon}}
\newunicodechar{Θ}{\ensuremath{\Theta}}
\newunicodechar{Λ}{\ensuremath{\Lambda}}
\newunicodechar{Μ}{\ensuremath{\mu}}
\newunicodechar{Τ}{\ensuremath{\tau}}
\newunicodechar{Φ}{\ensuremath{\Phi}}
\newunicodechar{Ψ}{\ensuremath{\Psi}}
\newunicodechar{Ω}{\ensuremath{\Omega}}
\newunicodechar{↦}{\ensuremath{\mapsto}}
\newunicodechar{∈}{\ensuremath{\in}}
\newunicodechar{∉}{\ensuremath{\notin}}
\newunicodechar{∧}{\ensuremath{\wedge}}
\newunicodechar{⇔}{\ensuremath{\Leftrightarrow}}
\newunicodechar{⟺}{\ensuremath{\Longleftrightarrow}}
\newunicodechar{⇒}{\ensuremath{\Rightarrow}}
\newunicodechar{⟹}{\ensuremath{\Longrightarrow}}
\newunicodechar{∘}{\ensuremath{\circ}}
\newunicodechar{∪}{\ensuremath{\cup}}
\newunicodechar{∩}{\ensuremath{\cap}}
\newunicodechar{≡}{\ensuremath{\equiv}}
\newunicodechar{⊂}{\ensuremath{\subset}}
\newunicodechar{⊥}{\ensuremath{\perp}}
\newunicodechar{∃}{\ensuremath{\exists}}
\newunicodechar{∀}{\ensuremath{\forall}}
\newunicodechar{∛}{\ensuremath{\sqrt[3]{\,}}}
\newunicodechar{∜}{\ensuremath{\sqrt[4]{\,}}}
\newunicodechar{⃗}{\ensuremath{{}^{\to}}}
\newunicodechar{ⁿ}{\ifmmode^{n}\else\textsuperscript{\ensuremath{n}}\fi}
\newunicodechar{ˣ}{\ifmmode^{x}\else\textsuperscript{\ensuremath{x}}\fi}
\newunicodechar{ᵇ}{\ifmmode^{b}\else\textsuperscript{\ensuremath{b}}\fi}
\newunicodechar{ʳ}{\ifmmode^{r}\else\textsuperscript{\ensuremath{r}}\fi}
\newunicodechar{ᵘ}{\ifmmode^{u}\else\textsuperscript{\ensuremath{u}}\fi}
\newunicodechar{ʸ}{\ifmmode^{y}\else\textsuperscript{\ensuremath{y}}\fi}
\newunicodechar{ᵃ}{\ifmmode^{a}\else\textsuperscript{\ensuremath{a}}\fi}
\newunicodechar{ᵉ}{\ifmmode^{e}\else\textsuperscript{\ensuremath{e}}\fi}
\newunicodechar{ᵏ}{\ifmmode^{k}\else\textsuperscript{\ensuremath{k}}\fi}
\newunicodechar{ᵖ}{\ifmmode^{p}\else\textsuperscript{\ensuremath{p}}\fi}
\newunicodechar{ᴺ}{\ifmmode^{N}\else\textsuperscript{\ensuremath{N}}\fi}
\newunicodechar{ᶜ}{\ifmmode^{c}\else\textsuperscript{\ensuremath{c}}\fi}
\newunicodechar{ʰ}{\ifmmode^{h}\else\textsuperscript{\ensuremath{h}}\fi}
\newunicodechar{ᵠ}{\ifmmode^{q}\else\textsuperscript{\ensuremath{q}}\fi}
\newunicodechar{ₙ}{\ifmmode_{n}\else\textsubscript{\ensuremath{n}}\fi}
\newunicodechar{ₐ}{\ifmmode_{a}\else\textsubscript{\ensuremath{a}}\fi}
\newunicodechar{ᵢ}{\ifmmode_{i}\else\textsubscript{\ensuremath{i}}\fi}
\newunicodechar{ᵣ}{\ifmmode_{r}\else\textsubscript{\ensuremath{r}}\fi}
\newunicodechar{ₖ}{\ifmmode_{k}\else\textsubscript{\ensuremath{k}}\fi}
\newunicodechar{ᵦ}{\ifmmode_{\beta}\else\textsubscript{\ensuremath{\beta}}\fi}
\newunicodechar{ₓ}{\ifmmode_{x}\else\textsubscript{\ensuremath{x}}\fi}
\newfontfamily\popdisplay{ArchivoBlack-Regular.ttf}[Path=../../../fonts/]
\newfontfamily\poplogo{SpaceGrotesk-Bold.ttf}[Path=../../../fonts/]
\newfontfamily\popsemi{PlusJakartaSans-SemiBold.ttf}[Path=../../../fonts/]
\newfontfamily\popxbold{PlusJakartaSans-ExtraBold.ttf}[Path=../../../fonts/]
\newfontfamily\popxboldls{PlusJakartaSans-ExtraBold.ttf}[Path=../../../fonts/,LetterSpace=3.0]

\setlist[enumerate]{leftmargin=1.7em,itemsep=5pt,topsep=4pt}
\setlist[itemize]{leftmargin=1.7em,itemsep=4pt,topsep=4pt,label={\color{poporange}\textbullet}}
\setlength{\parindent}{0pt}
\setlength{\parskip}{5pt}
\renewcommand{\arraystretch}{1.25}

% pgfplots : style Pop par défaut
\pgfplotsset{
  every axis/.append style={axis line style={popink,line width=1pt},tick style={popink},
    grid style={popline,line width=0.5pt},label style={font=\small},tick label style={font=\footnotesize}},
  popcurve/.style={poporange,line width=1.8pt,no markers,smooth},
}
% Même style disponible dans un \draw TikZ simple (hors pgfplots)
\tikzset{popcurve/.style={poporange,line width=1.8pt}}

% ============================================================
% Pill / squiggle
% ============================================================
\newcommand{\poppill}[3]{%
  \tikz[baseline=-0.55ex]\node[draw=popink,line width=1.2pt,rounded corners=2.6pt,%
    fill=#2,inner xsep=6.5pt,inner ysep=3pt]{%
    \popxboldls\fontsize{7}{7}\selectfont\color{#3}\MakeUppercase{#1}};%
}
\newcommand{\popsquiggle}[1]{%
  \tikz[baseline=-0.4ex]{
    \draw[#1,line width=2.6pt,line cap=round]
      plot [smooth] coordinates {(0,0.09) (0.34,0.01) (0.68,0.21) (1.02,0.01) (1.36,0.10)};
  }%
}
% Mot-clé surligné (marqueur orange sous le texte)
\newcommand{\cle}[1]{{\popxbold\color{popdark}#1}}

% ============================================================
% En-tête / pied
% ============================================================
\pagestyle{fancy}
\fancyhf{}
\renewcommand{\headrulewidth}{0pt}
\renewcommand{\footrulewidth}{0pt}
\lhead{\raisebox{-0.15ex}{\poplogo\fontsize{13}{13}\selectfont\color{popink}OnBuch\color{poporange}+}}
\rhead{\poppill{\DOCMATIERE\ \textbullet\ \DOCNIVEAU\ \textbullet\ Leçon \DOCLECON}{white}{popink}}
\lfoot{\popxbold\fontsize{7.6}{7.6}\selectfont\color{popmuted}\MakeUppercase{Cours signé OnBuch+}\ \textbullet\ {\popsemi\DOCTITRE}}
\rfoot{\tikz[baseline=-0.6ex]\node[rounded corners=8.5pt,fill=popink,minimum height=17pt,minimum width=17pt,inner xsep=5pt,inner sep=0]{\popxbold\fontsize{8}{8}\selectfont\color{popcream}\thepage\,/\,\pageref*{LastPage}};}

% ============================================================
% Titres
% ============================================================
\newcommand{\chaptitre}[2]{%
  \begin{tcolorbox}[enhanced,colback=poporange,colframe=popink,boxrule=2.6pt,arc=11pt,
      drop shadow=popink,left=15pt,right=15pt,top=12pt,bottom=11pt,before skip=3pt,after skip=18pt]
    \poppill{#2}{popink}{popcream}\par\vspace{9pt}
    {\raggedright\hyphenpenalty=10000\exhyphenpenalty=10000\popdisplay\fontsize{22}{24}\selectfont\color{popcream}\MakeUppercase{#1}\par}
  \end{tcolorbox}
}
% Section numérotée : pastille orange avec le numéro + titre Archivo
\newcounter{popsec}
\newcommand{\coursec}[1]{%
  \stepcounter{popsec}\setcounter{popsub}{0}%
  \par\vspace{16pt}\noindent
  \begin{minipage}{\linewidth}
  \tikz[baseline=-0.7ex]\node[draw=popink,line width=1.6pt,fill=poporange,rounded corners=4pt,minimum size=22pt,inner sep=0]{\popdisplay\fontsize{12}{12}\selectfont\color{popcream}\thepopsec};%
  \hspace{8pt}{\popdisplay\fontsize{15}{17}\selectfont\color{popink}\MakeUppercase{#1}}\par
  \vspace{4pt}\noindent{\color{poporange}\rule{36pt}{3.6pt}}
  \end{minipage}\par\nopagebreak\vspace{8pt}
}
\newcounter{popsub}
\newcommand{\courssub}[1]{%
  \stepcounter{popsub}%
  \par\vspace{9pt}\noindent{\popxbold\fontsize{11.5}{13}\selectfont\color{popink}{\color{poporange}\thepopsec.\thepopsub}\hspace{6pt}#1}\par\nopagebreak\vspace{4pt}
}

% ============================================================
% Boîtes pédagogiques — même recette : fond blanc, bordure épaisse colorée,
% ombre dure encre, badge plein. Titre optionnel [#1].
% ============================================================
\tcbset{popbase/.style={breakable,enhanced,colback=white,boxrule=2.3pt,arc=9pt,
  shadow={3pt}{-3pt}{0pt}{popink},left=14pt,right=14pt,top=11pt,bottom=11pt,before skip=11pt,after skip=15pt,
  fontupper=\popsemi\fontsize{10.6}{14.5}\selectfont\color{popink},
  fontlower=\popsemi\fontsize{10.6}{14.5}\selectfont\color{popink}}}
% Coche dessinée (le glyphe ✓ n'existe pas dans les polices du document)
\newcommand{\popcheck}[1]{\tikz[baseline=-0.5ex]\draw[#1,line width=1.6pt,line cap=round,line join=round](-0.13,0.0)--(-0.04,-0.09)--(0.13,0.1);}
\providecommand{\coche}{\popcheck{popgreen}}
\newcommand{\popboxhead}[3]{\poppill{#1}{#2}{white}\ifx\relax#3\relax\else\hspace{6pt}{\popxbold\fontsize{10.6}{12}\selectfont\color{popink}#3}\fi\par\vspace{7pt}}

\newtcolorbox{aretenir}[1][]{popbase,colframe=popgreen,before upper={\popboxhead{À retenir}{popgreen}{#1}}}
\newtcolorbox{exemplebox}[1][]{popbase,colframe=poporange,before upper={\popboxhead{Exemple}{poporange}{#1}}}
\newtcolorbox{definition}[1][]{popbase,colframe=popblue,before upper={\popboxhead{Définition}{popblue}{#1}}}
\newtcolorbox{propriete}[1][]{popbase,colframe=poppurple,before upper={\popboxhead{Propriété}{poppurple}{#1}}}
\newtcolorbox{methode}[1][]{popbase,colframe=popgold,before upper={\popboxhead{Méthode}{popgold}{#1}}}
\newtcolorbox{attention}[1][]{popbase,colframe=poppink,before upper={\popboxhead{Attention}{poppink}{#1}}}
\newtcolorbox{experience}[1][]{popbase,colframe=popblue,colback=popblueL,before upper={\popboxhead{Expérience}{popblue}{#1}}}
% « Le savais-tu ? » : carte violette pleine (contexte, culture, Cameroun)
\newtcolorbox{savaistu}[1][]{popbase,colback=poppurple,colframe=popink,
  fontupper=\popsemi\fontsize{10.4}{14.2}\selectfont\color{white},
  before upper={\poppill{Le savais-tu\,?}{white}{poppurple}\ifx\relax#1\relax\else\hspace{6pt}{\popxbold\color{white}#1}\fi\par\vspace{7pt}}}
% Exercice résolu : énoncé en haut, \tcblower puis la solution sur fond crème
\newcounter{popexo}\newcounter{popexr}
\newtcolorbox{exoresolu}[1][]{popbase,colframe=poporange,colbacklower=popcream,
  segmentation style={popink,line width=1.2pt,dashed},
  before upper={\stepcounter{popexr}\popboxhead{Exercice résolu \thepopexr}{poporange}{#1}},
  before lower={\poppill{Solution}{popink}{popcream}\par\vspace{6pt}}}
% Exercice (non résolu en place) — la correction est dans la partie Corrigés
\newtcolorbox{exercice}[1][]{popbase,colframe=popink,
  before upper={\stepcounter{popexo}\popboxhead{Exercice \thepopexo}{popink}{#1}}}
% Corrigé
\newtcolorbox{corrige}[1][]{popbase,colframe=popgreen,colback=popgreenL,
  before upper={\popboxhead{Corrigé}{popgreen}{#1}}}
% Objectifs / prérequis / bilan
\newtcolorbox{objectifs}{popbase,colframe=popink,colback=poporangeL,
  before upper={\popboxhead{Objectifs de la leçon}{popink}{}}}
\newtcolorbox{prerequis}{popbase,colframe=popink,colback=popblueL,
  before upper={\popboxhead{Prérequis}{popblue}{}}}
\newtcolorbox{bilan}{popbase,colframe=popink,colback=white,boxrule=2.8pt,
  before upper={\popboxhead{Fiche bilan}{poporange}{L'essentiel à connaître pour l'examen}}}
% Figure encadrée avec légende
\newtcolorbox{popfigure}[1][]{enhanced,breakable=false,colback=white,colframe=popline,boxrule=1.4pt,arc=8pt,
  left=10pt,right=10pt,top=10pt,bottom=8pt,before skip=10pt,after skip=12pt,halign=center,
  lower separated=false,halign lower=center,
  fontlower=\popsemi\fontsize{9}{11.5}\selectfont\color{popmuted},
  #1}
\newcommand{\legende}[1]{\par\vspace{6pt}{\popsemi\fontsize{9}{11.5}\selectfont\color{popmuted}#1}}

% ============================================================
% Signature OnBuch+
% ============================================================
\newcommand{\poplogoinline}[1]{{\poplogo\fontsize{#1}{#1}\selectfont\color{popink}OnBuch\color{poporange}+}}
\newcommand{\signatureonbuch}{%
  \begin{tcolorbox}[enhanced,colback=popink,colframe=popink,boxrule=2.2pt,arc=10pt,
      drop shadow=poporange,left=14pt,right=14pt,top=10pt,bottom=10pt,before skip=0pt,after skip=0pt]
    \tikz[baseline=-0.7ex]\node[circle,fill=popgreen,draw=popcream,line width=1.3pt,minimum size=22pt,inner sep=0]{\popcheck{white}};%
    \hspace{9pt}\begin{minipage}[c]{0.62\linewidth}
      {\popxbold\fontsize{12}{13}\selectfont\color{popcream}Signé\ \textbullet\ {\poplogo OnBuch\color{poporange}+}}\par\vspace{2pt}
      {\raggedright\popsemi\fontsize{8}{10.5}\selectfont\color{popline}\MakeUppercase{Équipe pédagogique OnBuch+}\\\MakeUppercase{Conforme au programme officiel MINESEC}\par}
    \end{minipage}\hfill
    {\poplogo\fontsize{22}{22}\selectfont\color{popcream}OnBuch\color{poporange}+}
  \end{tcolorbox}%
}

% ============================================================
% Page de garde
% #1 titre · #2 matière · #3 niveau · #4 module · #5 n° leçon · #6 durée ·
% #7 description / objectifs (texte ou itemize)
% ============================================================
\newcommand{\pagegarde}[7]{%
  \thispagestyle{empty}
  \newgeometry{a4paper,margin=1.7cm,top=1.6cm,bottom=1.4cm}%
  \begin{tikzpicture}[remember picture,overlay]
    \fill[poporange!18] ([xshift=-3.2cm,yshift=-3.6cm]current page.north east) circle (4.6cm);
    \fill[poppurple!14] ([xshift=2.4cm,yshift=7.6cm]current page.south west) circle (3.8cm);
  \end{tikzpicture}%
  {\poplogo\fontsize{30}{30}\selectfont\color{popink}OnBuch\color{poporange}+}\hfill
  \raisebox{0.6ex}{\poppill{Cours signé \textbullet\ Premium}{popink}{popcream}}\par
  \vspace{1pt}\popsquiggle{poppurple}\par
  \vspace{8pt}
  \begin{tcolorbox}[enhanced,colback=poporange,colframe=popink,boxrule=2.8pt,arc=13pt,
      drop shadow=popink,left=18pt,right=18pt,top=12pt,bottom=13pt,after skip=10pt]
    \poppill{#2 \textbullet\ #3}{popink}{popcream}\hspace{5pt}\poppill{#4}{white}{popink}\par\vspace{8pt}
    {\popdisplay\fontsize{14}{14}\selectfont\color{popink}LEÇON #5}\par\vspace{4pt}
    {\raggedright\hyphenpenalty=10000\exhyphenpenalty=10000\popdisplay\fontsize{25}{27}\selectfont\color{popcream}\MakeUppercase{#1}\par}
    \vspace{8pt}
    {\popxbold\fontsize{9.5}{9.5}\selectfont\color{popink}\MakeUppercase{Durée conseillée : #6}}
  \end{tcolorbox}
  \begin{tcolorbox}[enhanced,colback=white,colframe=popink,boxrule=2.2pt,arc=10pt,
      drop shadow=popink,left=16pt,right=16pt,top=10pt,bottom=10pt,after skip=8pt]
    \poppill{Dans cette leçon}{poppurple}{white}\par\vspace{5pt}
    \popsemi\fontsize{9.3}{12.6}\selectfont\color{popink}#7
  \end{tcolorbox}
  \noindent\begin{minipage}[t]{0.487\linewidth}
    \begin{tcolorbox}[enhanced,colback=popgreen,colframe=popink,boxrule=2.2pt,arc=10pt,
        drop shadow=popink,left=11pt,right=11pt,top=10pt,bottom=10pt,height=3.1cm,valign=center]
      \tcbox[colback=white,colframe=popink,boxrule=1.3pt,arc=5pt,boxsep=0pt,nobeforeafter,
        left=3pt,right=3pt,top=3pt,bottom=3pt,valign=center]{\qrcode[height=1.9cm]{https://play.google.com/store/apps/details?id=com.luvvix.onbuch&hl=fr}}%
      \hspace{8pt}\begin{minipage}[c]{0.5\linewidth}
        {\popdisplay\fontsize{11}{12.5}\selectfont\color{white}\MakeUppercase{Télécharge OnBuch}}\par\vspace{4pt}
        {\raggedright\popsemi\fontsize{8.4}{11}\selectfont\color{white}Gratuit \textbullet\ Google Play\\Tuteur IA, annales, concours, résultats\par}
      \end{minipage}
    \end{tcolorbox}
  \end{minipage}\hfill
  \begin{minipage}[t]{0.487\linewidth}
    \begin{tcolorbox}[enhanced,colback=poppurple,colframe=popink,boxrule=2.2pt,arc=10pt,
        drop shadow=popink,left=11pt,right=11pt,top=10pt,bottom=10pt,height=3.1cm,valign=center]
      \tikz[baseline=-0.7ex]\node[circle,fill=white,draw=popink,line width=1.3pt,minimum size=1.6cm,inner sep=0]{\popdisplay\fontsize{24}{24}\selectfont\color{poppurple}+};%
      \hspace{8pt}\begin{minipage}[c]{0.6\linewidth}
        {\popdisplay\fontsize{11}{12.5}\selectfont\color{white}\MakeUppercase{Passe à OnBuch+}}\par\vspace{4pt}
        {\raggedright\popsemi\fontsize{8.4}{11}\selectfont\color{white}Tous les cours signés, Tuteur IA illimité, fiches PDF — onglet OnBuch+ de l'app\par}
      \end{minipage}
    \end{tcolorbox}
  \end{minipage}
  \vspace{8pt}
  \popcontenu
  \vfill
  \signatureonbuch
  \restoregeometry
  \newpage
}

% Rangée « Ce cours contient » (page de garde)
\newcommand{\poptile}[3]{% #1 couleur #2 gros texte #3 légende
  \begin{tcolorbox}[enhanced,colback=#1,colframe=popink,boxrule=2pt,arc=9pt,shadow={2.5pt}{-2.5pt}{0pt}{popink},
      left=6pt,right=6pt,top=9pt,bottom=9pt,halign=center,valign=center,height=1.75cm,width=\linewidth,nobeforeafter]
    {\popdisplay\fontsize{12.5}{13}\selectfont\color{white}\MakeUppercase{#2}}\par\vspace{3pt}
    {\centering\popsemi\fontsize{7.8}{9.5}\selectfont\color{white}#3\par}
  \end{tcolorbox}}
\newcommand{\popcontenu}{%
  \par\noindent{\popxboldls\fontsize{7.5}{7.5}\selectfont\color{popmuted}CE COURS SIGNÉ CONTIENT}\par\vspace{7pt}
  \noindent\begin{minipage}[t]{0.235\linewidth}\poptile{popblue}{Cours}{complet, illustré, pas à pas}\end{minipage}\hfill
  \begin{minipage}[t]{0.235\linewidth}\poptile{poporange}{Méthodes}{et exercices résolus}\end{minipage}\hfill
  \begin{minipage}[t]{0.235\linewidth}\poptile{poppink}{Exercices}{type examen + corrigés}\end{minipage}\hfill
  \begin{minipage}[t]{0.235\linewidth}\poptile{popgreen}{Fiche bilan}{l'essentiel à retenir}\end{minipage}\par}

% Sommaire
\newtcolorbox{sommaire}{enhanced,colback=white,colframe=popink,boxrule=2.2pt,arc=10pt,
  drop shadow=popink,left=18pt,right=18pt,top=13pt,bottom=13pt,before skip=6pt,after skip=20pt,
  before upper={\poppill{Sommaire}{poppurple}{white}\par\vspace{9pt}\popsemi\fontsize{10.6}{14.5}\selectfont\color{popink}}}

% Signature de fin de cours — poussée tout en bas de la dernière page
\newcommand{\findecours}{%
  \par\vfill\nopagebreak
  \begin{center}\popsquiggle{poporange}\end{center}\vspace{4pt}
  \signatureonbuch
}

%% FIGFIX-BEGIN v5 — correctifs globaux des figures (docs/onbuch-generation/tools/figfix)
\usepackage{adjustbox}\usepackage{etoolbox}\tcbuselibrary{raster}
% toute figure TikZ/pgfplots plus large que la ligne est réduite à la largeur disponible
\BeforeBeginEnvironment{tikzpicture}{\begin{adjustbox}{max width=\linewidth}}
\AfterEndEnvironment{tikzpicture}{\end{adjustbox}}
% légendes pgfplots lisibles par-dessus les courbes
\pgfplotsset{every axis/.append style={legend style={fill=white,fill opacity=0.92,text opacity=1,draw=black!15,inner sep=3pt}}}
% étiquettes d'arêtes (syntaxe edge["..."]) avec fond blanc
\tikzset{every edge quotes/.append style={fill=white,inner sep=1.2pt}}
%% FIGFIX-END
\begin{document}

\pagegarde{Le condensateur : charge, décharge et énergie}{Physique}{Tle D}{Module 3 : Électricité — évolutions temporelles des circuits électriques et électroniques}{P10}{3 h}{Cette leçon présente le condensateur, composant capable de stocker des charges électriques et de l'énergie. On y étudie sa constitution, sa capacité, ses associations, puis l'évolution temporelle de sa charge et de sa décharge dans un dipôle RC, enfin l'énergie qu'il emmagasine.
\vspace{4pt}
\begin{multicols}{2}\small
\begin{itemize}
\item À la fin de cette leçon, je sais décrire la constitution d'un condensateur et reconnaître son symbole normalisé
\item À la fin de cette leçon, je sais exploiter la relation q = C u entre charge et tension
\item À la fin de cette leçon, je sais calculer la capacité d'un condensateur plan à l'aide de C = εS/e
\item À la fin de cette leçon, je sais déterminer la capacité équivalente d'une association de condensateurs en série ou en parallèle
\item À la fin de cette leçon, je sais établir l'équation différentielle de la charge ou de la décharge d'un dipôle RC
\item À la fin de cette leçon, je sais exploiter la solution de cette équation et la constante de temps τ = RC
\end{itemize}\end{multicols}}

\begin{sommaire}
\begin{multicols}{2}
\begin{enumerate}
\item Le condensateur : constitution, charge et capacité
\item Capacité du condensateur plan
\item Associations de condensateurs
\item Charge du condensateur dans un dipôle RC
\item Décharge du condensateur dans un résistor
\item Énergie emmagasinée par un condensateur
\item Activité d'intégration
\item Méthodes et exercices résolus
\item Exercices
\item Corrigés des exercices
\item Fiche bilan
\end{enumerate}
\end{multicols}
\end{sommaire}

\begin{objectifs}
\begin{itemize}
\item À la fin de cette leçon, je sais décrire la constitution d'un condensateur et reconnaître son symbole normalisé
\item À la fin de cette leçon, je sais exploiter la relation q = C u entre charge et tension
\item À la fin de cette leçon, je sais calculer la capacité d'un condensateur plan à l'aide de C = εS/e
\item À la fin de cette leçon, je sais déterminer la capacité équivalente d'une association de condensateurs en série ou en parallèle
\item À la fin de cette leçon, je sais établir l'équation différentielle de la charge ou de la décharge d'un dipôle RC
\item À la fin de cette leçon, je sais exploiter la solution de cette équation et la constante de temps τ = RC
\item À la fin de cette leçon, je sais déterminer graphiquement la constante de temps τ sur une courbe expérimentale
\item À la fin de cette leçon, je sais calculer l'énergie emmagasinée par un condensateur sous ses trois formes équivalentes
\item À la fin de cette leçon, je sais réaliser le schéma d'un circuit comportant un condensateur en respectant les conventions
\end{itemize}
\end{objectifs}

\begin{prerequis}
\begin{itemize}
\item Loi d'Ohm et lois des mailles et des nœuds (Première)
\item Intensité du courant comme débit de charges : i = dq/dt
\item Notion de champ électrostatique uniforme et relation E = U/d (Première)
\item Dérivation et équations différentielles du premier ordre (mathématiques, Tle)
\item Fonction exponentielle et ses propriétés graphiques
\end{itemize}
\end{prerequis}

\newpage

\coursec{Le condensateur : constitution, charge et capacité}

\courssub{Situation d'accroche : le flash qui survit à la coupure d'Eneo}

C'est le soir d'un mariage à Yaoundé. Soudain, coupure générale d'Eneo : la salle plonge dans le noir, les enceintes se taisent. Pourtant, le photographe continue de travailler : à chaque cliché, son flash projette une lumière éblouissante, encore et encore, pendant plusieurs secondes. D'où vient cette énergie, alors que tout est éteint ? Quelque part, dans un petit composant cylindrique caché dans l'appareil, elle était stockée \emph{avant} la coupure.

Le technicien qui répare les téléviseurs au quartier, lui, connaît bien la règle : même débranché du secteur, un poste de télévision reste dangereux pendant quelques minutes. Il attend, ou il « vide » prudemment un composant avant de mettre les mains dedans. Ce composant « qui reste chargé », c'est le \cle{condensateur}. On le trouve partout : dans les radios, les chargeurs de téléphone, les flashs d'appareils photo, les démarreurs de moteurs, les cartes électroniques des décodeurs.

\textbf{Problématique.} Comment ce petit composant stocke-t-il des charges électriques et de l'énergie ? Quelle relation relie la charge accumulée à la tension à ses bornes ? Et comment le courant se comporte-t-il dans sa branche ? C'est tout l'objet de cette première section.

\courssub{Constitution et symbole du condensateur}

\begin{experience}[Observer pour comprendre : deux conducteurs qui se font face]
Prenons deux plaques métalliques planes, placées face à face, séparées par une fine feuille de plastique, et reliées chacune à une borne d'un générateur. Aucun contact électrique direct n'existe entre les plaques : le plastique est un isolant. Pourtant, un galvanomètre placé dans le circuit détecte un bref passage de courant à la fermeture du circuit : des charges se sont déplacées et se sont \emph{accumulées} sur les plaques. Le dispositif a « emmagasiné » de la charge. Nous venons de fabriquer un condensateur.
\end{experience}

\begin{definition}[Condensateur]
Un \cle{condensateur} est un dipôle constitué de deux conducteurs, appelés \cle{armatures}, séparés par un isolant appelé \cle{diélectrique}. Un condensateur peut emmagasiner des charges électriques sur ses armatures et, par suite, de l'énergie électrique.
\end{definition}

Le diélectrique joue un double rôle : il empêche le contact direct entre les armatures (pas de court-circuit) et, comme nous le verrons dans la section suivante, il augmente la capacité du composant. Les diélectriques usuels sont :

\begin{itemize}
\item l'\textbf{air} (condensateurs variables des vieux postes radio) ;
\item le \textbf{mica} et la \textbf{céramique} (condensateurs de petite taille, très stables, utilisés dans les circuits de téléphonie) ;
\item les \textbf{films plastiques} (polyester, polypropylène) ;
\item l'\textbf{oxyde d'aluminium ou de tantale} formé par électrolyse : c'est le principe des condensateurs dits \emph{électrochimiques}, de forte capacité, reconnaissables à leur forme cylindrique et à leur polarité marquée.
\end{itemize}

\begin{attention}[Condensateurs polarisés]
Les condensateurs électrochimiques sont \cle{polarisés} : ils possèdent une borne $+$ et une borne $-$ qu'il faut impérativement respecter lors du branchement. Branché à l'envers, un tel condensateur chauffe et peut exploser. Au laboratoire, repère toujours la bande marquée sur le boîtier avant de fermer le circuit.
\end{attention}

Le symbole normalisé du condensateur est constitué de deux traits parallèles qui évoquent les deux armatures face à face :

\begin{center}
\begin{circuitikz}[european]
\draw (0,0) to[C, l=$C$] (2.5,0);
\end{circuitikz}
\end{center}

Dans un schéma de circuit, on le représente ainsi, en précisant la tension $u$ à ses bornes et le courant $i$ de sa branche.

\courssub{Le phénomène de charge : des armatures qui portent $+q$ et $-q$}

Pourquoi des charges s'accumulent-elles sur les armatures alors que le diélectrique est isolant ? C'est le phénomène d'\cle{influence électrostatique}.

Lorsqu'on relie le condensateur à un générateur, celui-ci agit comme une « pompe à électrons » : il arrache des électrons à l'armature reliée à sa borne $+$ et en dépose sur l'armature reliée à sa borne $-$. L'armature privée d'électrons porte alors la charge $+q$ ; l'autre, enrichie en électrons, porte la charge $-q$. Les charges $+q$ et $-q$, situées en regard l'une de l'autre de part et d'autre du diélectrique, s'attirent : elles restent « collées » sur les armatures, sans pouvoir se rejoindre à travers l'isolant.

\begin{propriete}[Conservation de la charge]
À tout instant, les deux armatures d'un condensateur portent des charges opposées : $+q$ sur l'une, $-q$ sur l'autre. La charge totale du condensateur est nulle ; on dit qu'il est \emph{chargé} lorsque $q \neq 0$. Par convention, la \cle{charge du condensateur} est la charge $q$ de l'armature vers laquelle arrive le courant $i$ (armature d'entrée du courant en convention récepteur).
\end{propriete}

Entre les armatures, les charges $+q$ et $-q$ créent un champ électrique $\vec{E}$ uniforme (pour un condensateur plan), dirigé de l'armature positive vers l'armature négative. C'est dans ce champ que réside l'énergie stockée.

\begin{popfigure}
\begin{center}
\begin{tikzpicture}[scale=1.0]
% Armatures
\fill[popblue!80] (-0.15,-2) rectangle (0.15,2);
\fill[poporange!80] (2.85,-2) rectangle (3.15,2);
% Charges
\foreach \y in {-1.5,-0.75,0,0.75,1.5}
  \node[popdark] at (0.28,\y) {$+$};
\foreach \y in {-1.5,-0.75,0,0.75,1.5}
  \node[popdark] at (2.72,\y) {$-$};
% Champ E
\foreach \y in {-1.2,-0.4,0.4,1.2}
  \draw[-{Stealth[length=2.5mm]},popgreen!70!black,thick] (0.6,\y) -- (2.4,\y);
\node[popgreen!70!black] at (1.5,1.65) {$\vec{E}$};
% Labels charges
\node[popblue!90!black] at (0,2.35) {armature $A$ : $+q$};
\node[poporange!90!black] at (3,2.35) {armature $B$ : $-q$};
% Surface S
\draw[{Stealth}-{Stealth},popdark] (-0.6,-1.8) -- (-0.6,1.8) node[midway,left] {$S$};
% Distance e
\draw[{Stealth}-{Stealth},popdark] (0,-2.5) -- (3,-2.5) node[midway,below] {$e$};
% Dielectrique
\node[popmuted] at (1.5,-1.75) {diélectrique (isolant)};
\end{tikzpicture}
\end{center}
\legende{Condensateur plan : deux armatures de surface $S$, distantes de $e$, portant les charges $+q$ et $-q$. Le champ électrique $\vec{E}$ est uniforme entre les armatures.}
\end{popfigure}

\courssub{Expérience fondamentale : la charge est proportionnelle à la tension}

\begin{experience}[La charge proportionnelle à la tension]
\textbf{Matériel :} un générateur de tension continue réglable, un condensateur (par exemple $C = 100~\mu$F), un interrupteur, un galvanomètre balistique (ou un dispositif de mesure de charge), des fils.

\textbf{Protocole :}
\begin{enumerate}
\item On règle le générateur sur une tension $u$ connue et on charge complètement le condensateur.
\item On bascule l'interrupteur pour décharger le condensateur à travers le galvanomètre balistique, qui mesure la charge $q$ accumulée.
\item On recommence pour plusieurs valeurs de $u$ et on reporte $q$ en fonction de $u$.
\end{enumerate}

\textbf{Observations :} les points $(u, q)$ sont alignés sur une droite passant par l'origine :

\begin{center}
\begin{tabular}{|c|c|c|c|c|}
\hline
\rowcolor{popblueL} $u$ (V) & 3,0 & 6,0 & 9,0 & 12,0 \\
\hline
$q$ ($10^{-4}$ C) & 3,0 & 6,0 & 9,0 & 12,0 \\
\hline
$q/u$ ($10^{-4}$ C/V) & 1,0 & 1,0 & 1,0 & 1,0 \\
\hline
\end{tabular}
\end{center}

\textbf{Interprétation :} le rapport $q/u$ est constant pour un condensateur donné. Cette constante caractérise le condensateur : c'est sa \emph{capacité}.
\end{experience}

\begin{propriete}[Relation charge--tension (admise à partir de l'expérience)]
La charge $q$ portée par l'armature positive d'un condensateur est proportionnelle à la tension $u$ entre ses bornes :
\[ q = C\,u \]
où $C$ est la \cle{capacité} du condensateur, grandeur positive qui ne dépend que de sa géométrie et de son diélectrique.
\begin{itemize}
\item $q$ en coulombs (C) ;
\item $u$ en volts (V) ;
\item $C$ en \cle{farads} (F).
\end{itemize}
Vérification dimensionnelle : $1~\text{F} = 1~\text{C}\cdot\text{V}^{-1}$.
\end{propriete}

\begin{aretenir}[Le farad et ses sous-multiples]
Le farad est une unité \emph{très grande}. Les condensateurs courants ont des capacités comprises entre le picofarad et le millifarad :
\[ 1~\mu\text{F} = 10^{-6}~\text{F} \qquad 1~\text{nF} = 10^{-9}~\text{F} \qquad 1~\text{pF} = 10^{-12}~\text{F} \]
Ordres de grandeur : céramique de découplage $100$ nF ; condensateur de flash $100~\mu$F à $1$ mF ; condensateur de démarrage de moteur quelques dizaines de $\mu$F.
\end{aretenir}

\begin{exemplebox}[Charge d'un condensateur de flash]
Le condensateur du flash de l'appareil photo du mariage a une capacité $C = 100~\mu\text{F} = 1,0\times 10^{-4}$ F et est chargé sous la tension $u = 12$ V. Sa charge vaut :
\[ q = C\,u = 1,0\times 10^{-4} \times 12 = 1,2\times 10^{-3}~\text{C} \]
Soit $q = 1,2$ mC. C'est cette charge, brutalement libérée dans la lampe du flash, qui produit l'éclair lumineux.
\end{exemplebox}

\begin{savaistu}[Des picofarads aux supercondensateurs]
Le farad est une unité énorme : un condensateur de $1$ F de la taille d'un composant classique était impensable il y a encore trente ans. Les condensateurs courants vont du picofarad ($10^{-12}$ F, circuits radio des téléphones) au millifarad. Mais les \cle{supercondensateurs}, développés pour l'électronique de puissance, atteignent plusieurs centaines de farads ! Ils se rechargent en quelques secondes et servent à récupérer l'énergie de freinage des bus électriques ou à assurer la continuité d'alimentation des équipements télécoms — une technologie précieuse là où le réseau est instable.
\end{savaistu}

\courssub{Intensité du courant dans la branche : $i = C\,\dfrac{du}{dt}$}

Par définition, l'intensité du courant est le débit de charge électrique : c'est la quantité de charge qui traverse une section du conducteur par unité de temps. Pour la branche contenant le condensateur, la charge qui arrive sur l'armature d'entrée est précisément $q$, donc :
\[ i = \frac{dq}{dt} \]
Comme $q = C\,u$ et que la capacité $C$ est une constante du composant, on obtient la relation fondamentale :

\begin{propriete}[Relation courant--tension du condensateur]
En \cle{convention récepteur} (les flèches de $u$ et de $i$ sont de sens opposés sur le dipôle, le courant $i$ arrivant sur l'armature qui porte la charge $+q$) :
\[ i = \frac{dq}{dt} = C\,\frac{du}{dt} \]
avec $i$ en ampères (A), $C$ en farads (F), $u$ en volts (V) et $t$ en secondes (s).

Vérification dimensionnelle : $[\,C\,]\times[\,u\,]/[\,t\,] = \text{F}\cdot\text{V}\cdot\text{s}^{-1} = \text{C}\cdot\text{s}^{-1} = \text{A}$, ce qui est bien homogène à une intensité.
\end{propriete}

\begin{methode}[Avant d'écrire $i = C\,\dfrac{du}{dt}$ : orienter correctement]
\begin{enumerate}
\item Sur le schéma, choisir un sens pour le courant $i$ dans la branche du condensateur et le flécher.
\item Flécher la tension $u$ aux bornes du condensateur \emph{en sens opposé} à $i$ (convention récepteur).
\item Vérifier que la charge $q$ est bien celle de l'armature sur laquelle arrive $i$.
\item Alors seulement, écrire $i = \dfrac{dq}{dt} = C\,\dfrac{du}{dt}$. Si la convention est inversée, un signe moins apparaît : c'est la source d'erreur la plus fréquente dans les équations différentielles des dipôles RC.
\end{enumerate}
\end{methode}

\begin{popfigure}
\begin{center}
\begin{circuitikz}[european, scale=1.0, transform shape]
\draw (0,0) to[battery1, l=$E$] (0,3)
      to[nos, l=K] (2.5,3)
      to[R, l=$R$, i=$i$] (5,3)
      to[C, l=$C$, v=$u$] (5,0)
      -- (0,0);
\end{circuitikz}
\end{center}
\legende{Circuit de charge d'un condensateur : générateur de force électromotrice $E$, interrupteur K, résistor $R$, condensateur $C$. Le courant $i$ et la tension $u$ sont fléchés en convention récepteur pour le condensateur.}
\end{popfigure}

\courssub{Conséquence : le condensateur en régime permanent continu}

Supposons la tension $u$ aux bornes du condensateur \emph{constante} (régime permanent continu). Alors :
\[ \frac{du}{dt} = 0 \quad\Longrightarrow\quad i = C\,\frac{du}{dt} = 0 \]

\begin{propriete}[Condensateur chargé en régime continu]
En régime permanent continu, un condensateur chargé se comporte comme un \cle{interrupteur ouvert} : aucun courant ne circule dans sa branche. C'est pourquoi, une fois le flash chargé, la pile de l'appareil photo ne se décharge plus à travers lui : le condensateur « garde » sa charge et son énergie, prêtes à l'emploi.
\end{propriete}

\begin{attention}[Le courant ne « traverse » pas le condensateur !]
Erreur fréquente : imaginer que les électrons traversent le condensateur de part en part. C'est \textbf{faux} : le diélectrique est un isolant, aucune charge ne le franchit. Le courant $i$ existe dans les \emph{fils} de la branche : il traduit l'arrivée des charges sur une armature et le départ simultané des charges de l'autre armature. Tout se passe comme si le courant traversait le dipôle, mais physiquement ce sont des \emph{accumulations} de charges $+q$ et $-q$ sur les armatures. C'est exactement pour cela que le technicien peut recevoir une décharge en touchant les deux bornes d'un condensateur chargé, même appareil débranché : les charges stockées cherchent un chemin pour se recombiner.
\end{attention}

\courssub{Tension maximale d'usage : le claquage du diélectrique}

Un condensateur ne supporte pas n'importe quelle tension. Si $u$ devient trop grande, le champ électrique $\vec{E}$ dans le diélectrique devient si intense qu'il arrache des électrons aux atomes de l'isolant : celui-ci devient brutalement conducteur, une étincelle jaillit entre les armatures. C'est le \cle{claquage} du diélectrique, qui détruit généralement le composant.

\begin{definition}[Tension nominale]
La \cle{tension nominale} (ou tension maximale d'usage), inscrite sur le boîtier du condensateur (par exemple « $100~\mu$F / $25$ V »), est la tension maximale que le composant peut supporter durablement sans risque de claquage. Dans tout montage, on veille à ce que la tension appliquée reste inférieure à cette valeur, avec une marge de sécurité.
\end{definition}

Le diélectrique joue donc un rôle essentiel : il permet de rapprocher les armatures à très petite distance $e$ sans contact, ce qui augmente la capacité, tout en fixant la limite de tension supportable. C'est tout l'art du fabricant de condensateurs : choisir un isolant à la fois fin, très résistant au claquage et favorable à une grande capacité — nous y reviendrons dans la section suivante avec la formule $C = \varepsilon S/e$.

\begin{exoresolu}[Charge, capacité et tension de claquage]
Un condensateur porte les indications « $47~\mu$F — $16$ V ».

1. Calculer la charge maximale $q_{\max}$ qu'il peut emmagasiner.

2. On mesure dans un montage une charge $q = 3,76\times 10^{-4}$ C sur son armature positive. Quelle est la tension $u$ à ses bornes ? Le composant est-il en sécurité ?

3. Dans ce montage, la tension devient constante. Que vaut alors l'intensité du courant dans la branche ?

\tcblower
\textbf{Solution détaillée.}

1. La charge maximale correspond à la tension nominale :
\[ q_{\max} = C\,u_{\max} = 47\times 10^{-6} \times 16 = 7,52\times 10^{-4}~\text{C} \]
Soit $q_{\max} \approx 7,5\times 10^{-4}$ C $= 0,75$ mC.

2. De $q = C\,u$, on tire :
\[ u = \frac{q}{C} = \frac{3,76\times 10^{-4}}{47\times 10^{-6}} = 8,0~\text{V} \]
Comme $u = 8,0~\text{V} < 16~\text{V}$, la tension reste inférieure à la tension nominale : le condensateur fonctionne en toute sécurité, avec une marge confortable.

3. Si la tension est constante, $\dfrac{du}{dt} = 0$, donc :
\[ i = C\,\frac{du}{dt} = 0~\text{A} \]
Le condensateur chargé se comporte comme un interrupteur ouvert : plus aucun courant ne circule dans sa branche.
\end{exoresolu}

\begin{aretenir}[L'essentiel de la section]
\begin{itemize}
\item Un condensateur est un dipôle formé de deux armatures conductrices séparées par un diélectrique isolant ; il stocke des charges $+q$ et $-q$ par influence électrostatique.
\item La charge est proportionnelle à la tension : $q = C\,u$, avec $C$ la capacité en farads (F). Sous-multiples usuels : $\mu$F, nF, pF.
\item En convention récepteur : $i = \dfrac{dq}{dt} = C\,\dfrac{du}{dt}$.
\item En régime permanent continu, le condensateur chargé équivaut à un interrupteur ouvert ($i = 0$).
\item Aucune charge ne traverse le diélectrique ; au-delà de la tension nominale, le diélectrique claque et le composant est détruit.
\end{itemize}
\end{aretenir}

\coursec{Capacité du condensateur plan}

Dans la section précédente, nous avons découvert le condensateur : deux armatures conductrices séparées par un isolant (le \cle{diélectrique}), capables d'emmagasiner une charge $q$ proportionnelle à la tension $u$ appliquée, selon la loi $q = C\,u$. Nous avons aussi compris que la capacité $C$ dépend de la \emph{géométrie} du composant et de la \emph{nature} de l'isolant. Une question naturelle se pose alors : comment le fabricant d'un condensateur choisit-il la taille des armatures, leur écartement et le matériau isolant pour obtenir une capacité donnée ? Autrement dit, de quels paramètres dépend $C$, et selon quelle loi ? C'est tout l'objet de cette section, consacrée au cas fondamental du \cle{condensateur plan}, modèle de référence pour comprendre tous les autres.

\courssub{Géométrie du condensateur plan}

\begin{definition}[Condensateur plan]
Un \cle{condensateur plan} est constitué de deux armatures conductrices planes, parallèles, de surface $S$ en regard l'une de l'autre, séparées par une distance $e$ constante. L'espace entre les armatures est rempli par un diélectrique (isolant) de permittivité $\varepsilon$.
\end{definition}

Trois paramètres géométriques et électriques caractérisent donc complètement le condensateur plan :
\begin{itemize}
\item la \cle{surface en regard} $S$ des armatures (en mètres carrés, m$^2$) : c'est l'aire de la zone où les deux plaques se font face ;
\item la \cle{distance} $e$ entre les armatures (en mètres, m), souvent appelée épaisseur du diélectrique ;
\item la \cle{permittivité} $\varepsilon$ du diélectrique placé entre les armatures (en farads par mètre, F$\cdot$m$^{-1}$).
\end{itemize}

\begin{popfigure}
\begin{tikzpicture}[scale=1.0, >=Stealth]
% Armatures
\fill[popblue!70] (-2.2,-1.6) rectangle (-1.9,1.6);
\fill[poporange!80] (1.9,-1.6) rectangle (2.2,1.6);
\node[above] at (-2.05,1.65) {\small Armature A ($+q$)};
\node[above] at (2.05,1.65) {\small Armature B ($-q$)};
% Charges
\foreach \y in {-1.2,-0.6,0,0.6,1.2} {\node[white] at (-2.05,\y) {\small $+$}; \node[white] at (2.05,\y) {\small $-$};}
% Lignes de champ
\foreach \y in {-1.2,-0.6,0,0.6,1.2} {\draw[->, popgreen!80!black, thick] (-1.75,\y) -- (1.75,\y);}
\node[popgreen!60!black] at (0,1.45) {\small $\vec{E}$ uniforme};
% Cote e
\draw[<->, popdark, thick] (-1.9,-2.0) -- (1.9,-2.0) node[midway, below] {\small $e$};
% Cote S (perspective)
\draw[popdark] (2.2,-1.6) -- (2.9,-1.1);
\draw[popdark] (2.2,1.6) -- (2.9,2.1);
\draw[popdark] (2.9,-1.1) -- (2.9,2.1);
\draw[<->, popdark] (3.05,-1.05) -- (3.05,2.05) node[midway, right] {\small $S$};
% Diélectrique
\node[popmuted] at (0,-1.45) {\small diélectrique $\varepsilon = \varepsilon_0 \varepsilon_r$};
\end{tikzpicture}
\legende{Condensateur plan : deux armatures de surface $S$ en regard, séparées de $e$. Entre elles règne un champ électrique $\vec{E}$ uniforme, dirigé de l'armature positive vers l'armature négative.}
\end{popfigure}

Lorsque le condensateur est chargé, l'armature A porte la charge $+q$ et l'armature B la charge $-q$. Il règne entre elles un champ électrique $\vec{E}$ \emph{uniforme} (même direction, même sens, même norme en tout point, tant qu'on reste loin des bords), dirigé de la plaque positive vers la plaque négative. C'est ce champ qui « relie » la tension à la géométrie, comme nous le verrons en fin de section.

\courssub{Expression de la capacité}

\textbf{Intuition d'abord.} Imagine un entrepôt de stockage : plus sa surface au sol est grande, plus il peut accueillir de marchandises ; plus les rayonnages sont proches les uns des autres, plus l'ensemble est compact et efficace. Pour le condensateur, l'intuition est analogue :
\begin{itemize}
\item si l'on \emph{augmente la surface} $S$ des armatures, on offre plus de place aux charges : à tension égale, le condensateur stocke davantage de charge, donc $C$ augmente ;
\item si l'on \emph{rapproche les armatures} (on diminue $e$), les charges $+$ et $-$ en vis-à-vis s'attirent plus fortement à travers l'isolant : l'armature négative « aide » la positive à accueillir encore plus de charges, donc $C$ augmente ;
\item si l'on remplace l'air par un \emph{meilleur diélectrique} (permittivité $\varepsilon$ plus grande), le matériau se polarise et affaiblit le champ interne, ce qui permet de loger plus de charge à tension égale : $C$ augmente.
\end{itemize}

L'expérience (mesure de $C$ en faisant varier séparément $S$, $e$ et le diélectrique) confirme cette intuition et conduit à un résultat remarquablement simple.

\begin{propriete}[Capacité du condensateur plan (admise)]
La capacité d'un condensateur plan est donnée par :
\[ C = \dfrac{\varepsilon\, S}{e} \qquad \text{avec} \qquad \varepsilon = \varepsilon_0\,\varepsilon_r \]
où $S$ est la surface des armatures en regard (m$^2$), $e$ la distance entre elles (m), $\varepsilon_0 \approx 8{,}85\times 10^{-12}$~F$\cdot$m$^{-1}$ la permittivité du vide et $\varepsilon_r$ la permittivité relative du diélectrique (sans dimension, $\varepsilon_r \geq 1$).
\end{propriete}

\begin{aretenir}[Vérification dimensionnelle]
Vérifions l'homogénéité de $C = \varepsilon S / e$ :
\[ [\varepsilon] \times \dfrac{[S]}{[e]} = \mathrm{F\cdot m^{-1}} \times \dfrac{\mathrm{m^2}}{\mathrm{m}} = \mathrm{F}. \]
On obtient bien des farads : la formule est dimensionnellement cohérente. Retiens ce réflexe : une analyse dimensionnelle rapide permet de détecter immédiatement une formule mal recopiée (par exemple $\varepsilon e/S$, qui donnerait des F$\cdot$m$^{-2}$, est forcément fausse).
\end{aretenir}

\courssub{Permittivité du diélectrique : \texorpdfstring{$\varepsilon = \varepsilon_0 \varepsilon_r$}{epsilon = epsilon0 epsilonr}}

\begin{definition}[Permittivités absolue et relative]
La \cle{permittivité du vide} $\varepsilon_0$ est une constante fondamentale de la nature :
\[ \varepsilon_0 \approx 8{,}85\times 10^{-12}\ \mathrm{F\cdot m^{-1}}. \]
Chaque matériau isolant est caractérisé par sa \cle{permittivité relative} $\varepsilon_r$, nombre sans dimension supérieur ou égal à 1. La permittivité absolue du matériau est alors $\varepsilon = \varepsilon_0\,\varepsilon_r$.
\end{definition}

Physiquement, $\varepsilon_r$ mesure l'efficacité du diélectrique à se polariser sous l'effet du champ électrique : ses molécules s'orientent et créent un champ opposé qui affaiblit le champ total. À tension imposée, le condensateur peut alors stocker $\varepsilon_r$ fois plus de charge que s'il était vide. Voici quelques valeurs usuelles :

\begin{center}
\begin{tabularx}{\linewidth}{|l|X|}\hline
\rowcolor{popblueL} \textbf{Diélectrique} & \textbf{Permittivité relative $\varepsilon_r$ (ordre de grandeur)} \\ \hline
Vide & $1$ (exactement) \\ \hline
Air sec & $\approx 1{,}0006 \simeq 1$ \\ \hline
Papier paraffiné & $2$ à $3$ \\ \hline
Verre & $\approx 5$ \\ \hline
Mica & $\approx 7$ \\ \hline
Céramiques spéciales (titanates) & de $100$ à plusieurs milliers \\ \hline
\end{tabularx}
\end{center}

\begin{attention}[Ne pas confondre $\varepsilon_0$ et $\varepsilon_r$]
Erreur très fréquente au Baccalauréat : écrire $C = \varepsilon_r S / e$ ou $C = \varepsilon_0 S / e$ en oubliant l'un des deux facteurs. Retiens bien :
\begin{itemize}
\item $\varepsilon_0 = 8{,}85\times 10^{-12}$~F$\cdot$m$^{-1}$ : constante universelle, \emph{avec unité} ;
\item $\varepsilon_r$ : nombre \emph{sans dimension}, caractéristique du matériau ;
\item la formule exige le \textbf{produit} : $C = \dfrac{\varepsilon_0\,\varepsilon_r\, S}{e}$.
\end{itemize}
Pour un condensateur à air ($\varepsilon_r \approx 1$), on a simplement $C = \varepsilon_0 S/e$.
\end{attention}

\courssub{Influence des paramètres : proportionnalité}

La relation $C = \varepsilon S / e$ se lit comme suit :
\begin{itemize}
\item $C$ est \textbf{proportionnelle à $S$} : doubler la surface double la capacité ;
\item $C$ est \textbf{inversement proportionnelle à $e$} : diviser l'écartement par deux double la capacité ;
\item $C$ est \textbf{proportionnelle à $\varepsilon_r$} : remplacer l'air ($\varepsilon_r \approx 1$) par du mica ($\varepsilon_r \approx 7$) multiplie la capacité par 7.
\end{itemize}

Une façon élégante de vérifier expérimentalement la dépendance en $e$ consiste à tracer $C$ en fonction de $1/e$ : la relation $C = \varepsilon S \times \dfrac{1}{e}$ est celle d'une fonction linéaire de coefficient directeur $\varepsilon S$. On obtient une droite passant par l'origine.

\begin{popfigure}
\begin{tikzpicture}
\begin{axis}[width=0.85\linewidth, height=6.5cm,
xlabel={$1/e$ \ (m$^{-1}$)}, ylabel={$C$ \ (nF)},
xmin=0, xmax=12000, ymin=0, ymax=12,
xtick={0,2000,4000,6000,8000,10000,12000},
ytick={0,2,4,6,8,10,12},
grid=major, grid style={popline!50},
axis lines=left, tick label style={font=\small}, label style={font=\small}]
\addplot[popcurve, domain=0:12000, samples=50]{0.001*x};
\addplot[only marks, mark=*, mark size=1.8pt, color=poporange] coordinates {(2000,2.0) (4000,4.0) (6000,6.0) (8000,8.0) (10000,10.0)};
\node[popblue, font=\small, anchor=south west] at (axis cs:5500,4.2) {$C = \varepsilon S \times \dfrac{1}{e}$};
\node[popblue, font=\small, anchor=north west] at (axis cs:5500,3.6) {pente $= \varepsilon S$};
\end{axis}
\end{tikzpicture}
\legende{Capacité $C$ en fonction de $1/e$ pour un condensateur plan de surface $S$ et de diélectrique $\varepsilon$ fixés : droite passant par l'origine, de pente $\varepsilon S$. Ici $\varepsilon S = 10^{-3}$~nF$\cdot$m (condensateur à air de $S \approx 113$~cm$^2$).}
\end{popfigure}

\begin{methode}[Calculer une capacité sans se tromper]
Pour appliquer numériquement $C = \varepsilon_0 \varepsilon_r S / e$ :
\begin{enumerate}
\item \textbf{Convertir} la surface en mètres carrés : $1~\mathrm{cm^2} = 10^{-4}~\mathrm{m^2}$ ;
\item \textbf{Convertir} l'épaisseur en mètres : $1~\mathrm{mm} = 10^{-3}$~m, $1~\mu\mathrm{m} = 10^{-6}$~m ;
\item \textbf{Identifier} le diélectrique et prendre $\varepsilon = \varepsilon_0 \varepsilon_r$ (pour l'air, $\varepsilon_r \approx 1$) ;
\item \textbf{Appliquer} la formule et exprimer le résultat avec un préfixe adapté (pF, nF, $\mu$F) ;
\item \textbf{Vérifier} l'ordre de grandeur : les condensateurs plans de laboratoire donnent typiquement des pF à des nF ; les $\mu$F et plus relèvent des condensateurs chimiques ou céramiques multicouches.
\end{enumerate}
\end{methode}

\begin{exemplebox}[Calcul de la capacité d'un condensateur plan]
Un condensateur plan a des armatures de surface $S = 200$~cm$^2$ séparées par une distance $e = 0{,}1$~mm. Le diélectrique a une permittivité relative $\varepsilon_r = 6$.

Conversions : $S = 200\times 10^{-4}~\mathrm{m^2} = 2{,}0\times 10^{-2}$~m$^2$ et $e = 1{,}0\times 10^{-4}$~m.

Application numérique :
\begin{align*}
C &= \dfrac{\varepsilon_0\,\varepsilon_r\, S}{e} = \dfrac{8{,}85\times 10^{-12}\times 6\times 2{,}0\times 10^{-2}}{1{,}0\times 10^{-4}} \\
&= 1{,}06\times 10^{-8}\ \mathrm{F} \approx 10{,}6\ \mathrm{nF}.
\end{align*}
La capacité vaut environ $10{,}6$~nF. Avec de l'air à la place du diélectrique, on n'aurait obtenu que $C \approx 1{,}77$~nF : le diélectrique a bien multiplié la capacité par $\varepsilon_r = 6$.
\end{exemplebox}

\courssub{Complément utile au Bac : lien avec le champ électrique}

\begin{propriete}[Champ entre les armatures (complément)]
Entre les armatures d'un condensateur plan soumis à la tension $u$, le champ électrique est uniforme et sa norme vaut :
\[ E = \dfrac{u}{e}. \]
En combinant avec $q = Cu$ et $C = \varepsilon S/e$, on obtient la charge portée par chaque armature :
\[ q = \varepsilon\, S\, E. \]
Ces relations ne sont pas exigibles en elles-mêmes, mais elles éclairent les formules du cours et apparaissent régulièrement dans les sujets.
\end{propriete}

Vérifions rapidement la cohérence : $q = Cu = \dfrac{\varepsilon S}{e}\times u = \varepsilon S \times \dfrac{u}{e} = \varepsilon S E$. Dimensionnellement, $[\varepsilon S E] = \mathrm{F\cdot m^{-1}}\times \mathrm{m^2}\times \mathrm{V\cdot m^{-1}} = \mathrm{F\cdot V} = \mathrm{C}$ (coulombs), puisque $1~\mathrm{F} = 1~\mathrm{C\cdot V^{-1}}$. Tout est homogène.

Ce lien explique aussi une contrainte technologique majeure : on ne peut pas rapprocher indéfiniment les armatures pour augmenter $C$, car à tension fixée, diminuer $e$ augmente le champ $E = u/e$. Au-delà d'un champ critique (le \emph{champ disruptif} du diélectrique, environ $3\times 10^{6}$~V$\cdot$m$^{-1}$ pour l'air sec), l'isolant claque : une étincelle traverse le condensateur et le détruit. Tout condensateur commercial porte donc une \emph{tension maximale de service} gravée sur son boîtier.

\courssub{Applications technologiques}

\textbf{Condensateurs céramiques.} En utilisant des céramiques de très forte permittivité ($\varepsilon_r$ de quelques centaines à plusieurs milliers) et en empilant de fines couches d'armatures métalliques (technologie « multicouche »), les fabricants obtiennent des capacités de l'ordre du $\mu$F dans des boîtiers de quelques millimètres. Ce sont les condensateurs les plus répandus dans les téléphones portables et les cartes électroniques que l'on répare à Douala ou Yaoundé.

\textbf{Condensateurs variables.} Dans les postes radio analogiques, le réglage de la station (la « syntonisation ») se faisait avec un condensateur variable à air : un jeu de lames mobiles tourne entre des lames fixes, ce qui fait varier la \emph{surface en regard} $S$, donc la capacité $C = \varepsilon_0 S/e$, sans modifier l'écartement. Associé à une bobine, ce condensateur règle la fréquence de résonance du circuit d'accord — tu en reparleras dans l'étude des oscillations électriques.

\begin{savaistu}[Le secret des condensateurs électrolytiques]
Comment les condensateurs « chimiques » (électrolytiques) atteignent-ils des capacités de plusieurs milliers de $\mu$F dans un cylindre de la taille d'un dé ? Leur diélectrique est une couche d'oxyde d'aluminium d'épaisseur microscopique, de l'ordre du micromètre voire moins, formée électrochimiquement à la surface d'une feuille d'aluminium gravée (ce qui augmente énormément $S$). Comme $C = \varepsilon S/e$, un $e$ minuscule et un $S$ gigantesque donnent une capacité énorme. Contrepartie : ces condensateurs sont \emph{polarisés} — il faut respecter le sens de branchement, sous peine de destruction, parfois explosive ! C'est pourquoi on ne les utilise jamais en courant alternatif.
\end{savaistu}

\begin{exoresolu}[Exploitation complète de la formule du condensateur plan]
Énoncé. Un condensateur plan à air ($\varepsilon_r = 1$) est formé de deux disques de rayon $R = 10$~cm, distants de $e = 2{,}0$~mm. On donne $\varepsilon_0 = 8{,}85\times 10^{-12}$~F$\cdot$m$^{-1}$.

1. Calculer sa capacité $C$.

2. On intercale entre les armatures une plaque de verre ($\varepsilon_r = 5$) de même épaisseur. Que devient la capacité ?

3. Le condensateur au verre est soumis à une tension $u = 12$~V. Calculer la charge $q$ emmagasinée et la norme $E$ du champ entre les armatures.

4. Quelle tension maximale peut-on appliquer si le verre claque pour un champ de $10^{7}$~V$\cdot$m$^{-1}$ ?
\tcblower
Solution détaillée.

1. Surface des armatures : $S = \pi R^2 = \pi \times (0{,}10)^2 = 3{,}14\times 10^{-2}$~m$^2$. D'où :
\[ C_{\text{air}} = \dfrac{\varepsilon_0 S}{e} = \dfrac{8{,}85\times 10^{-12}\times 3{,}14\times 10^{-2}}{2{,}0\times 10^{-3}} = 1{,}39\times 10^{-10}\ \mathrm{F} \approx 139\ \mathrm{pF}. \]

2. Le verre multiplie la capacité par $\varepsilon_r = 5$ :
\[ C_{\text{verre}} = \varepsilon_r\, C_{\text{air}} = 5\times 139 \approx 695\ \mathrm{pF} \approx 0{,}70\ \mathrm{nF}. \]

3. Charge : $q = C_{\text{verre}}\, u = 6{,}95\times 10^{-10}\times 12 = 8{,}3\times 10^{-9}$~C, soit $q \approx 8{,}3$~nC. Champ :
\[ E = \dfrac{u}{e} = \dfrac{12}{2{,}0\times 10^{-3}} = 6{,}0\times 10^{3}\ \mathrm{V\cdot m^{-1}}. \]
Vérification par la relation complémentaire : $q = \varepsilon_0 \varepsilon_r S E = 8{,}85\times 10^{-12}\times 5\times 3{,}14\times 10^{-2}\times 6{,}0\times 10^{3} \approx 8{,}3\times 10^{-9}$~C. Cohérent.

4. Tension de claquage : $u_{\max} = E_{\text{disruptif}}\times e = 10^{7}\times 2{,}0\times 10^{-3} = 2{,}0\times 10^{4}$~V, soit $20$~kV.
\end{exoresolu}

\begin{attention}[Pièges classiques sur le condensateur plan]
\begin{itemize}
\item \textbf{Conversions} : $S$ en m$^2$ et $e$ en m, toujours ! Une surface laissée en cm$^2$ fausse le résultat d'un facteur $10^{4}$.
\item \textbf{Le produit $\varepsilon_0 \varepsilon_r$} : n'oublie aucun des deux facteurs ; $\varepsilon_r$ seul n'a pas d'unité et ne peut pas figurer seul dans une formule homogène.
\item \textbf{Surface en regard} : si les armatures ne se recouvrent que partiellement (condensateur variable), c'est l'aire de la zone commune qui compte, pas l'aire totale des plaques.
\item \textbf{Champ disruptif} : augmenter $C$ en réduisant $e$ se paie par une tension maximale d'utilisation plus faible.
\end{itemize}
\end{attention}

\begin{exercice}[Pour t'entraîner]
Un condensateur plan à air a une capacité $C_0 = 47$~pF. Ses armatures carrées ont $5{,}0$~cm de côté.

1. Calculer la distance $e$ entre les armatures.

2. On glisse du mica ($\varepsilon_r = 7$) entre les armatures et on réduit l'écartement à $e/2$. Quelle est la nouvelle capacité ?

3. Ce condensateur est chargé sous $u = 24$~V. Quelle charge porte chaque armature ?
\end{exercice}

\begin{corrige}[Exercice « Pour t'entraîner »]
1. $S = (5{,}0\times 10^{-2})^2 = 2{,}5\times 10^{-3}$~m$^2$. De $C_0 = \varepsilon_0 S/e$, on tire :
\[ e = \dfrac{\varepsilon_0 S}{C_0} = \dfrac{8{,}85\times 10^{-12}\times 2{,}5\times 10^{-3}}{47\times 10^{-12}} = 4{,}7\times 10^{-4}\ \mathrm{m} \approx 0{,}47\ \mathrm{mm}. \]

2. Les deux modifications multiplient $C$ par $\varepsilon_r = 7$ (diélectrique) et par $2$ (écartement divisé par deux) :
\[ C = 7\times 2\times 47 = 658\ \mathrm{pF} \approx 0{,}66\ \mathrm{nF}. \]

3. $q = Cu = 6{,}58\times 10^{-10}\times 24 \approx 1{,}6\times 10^{-8}$~C $= 16$~nC. L'armature reliée au pôle $+$ porte $+q$, l'autre $-q$.
\end{corrige}

\begin{aretenir}[L'essentiel de la section]
\begin{itemize}
\item Condensateur plan : deux armatures de surface $S$ en regard, séparées de $e$, diélectrique de permittivité $\varepsilon$.
\item Capacité : $C = \dfrac{\varepsilon_0\,\varepsilon_r\, S}{e}$, avec $\varepsilon_0 = 8{,}85\times 10^{-12}$~F$\cdot$m$^{-1}$ et $\varepsilon_r \geq 1$ sans dimension.
\item $C$ croît avec $S$ et $\varepsilon_r$, décroît avec $e$ ; le graphe $C = f(1/e)$ est une droite de pente $\varepsilon S$.
\item Complément : $E = u/e$ et $q = \varepsilon S E$ ; le champ disruptif limite la tension d'usage.
\item Avant tout calcul : convertir $S$ en m$^2$ et $e$ en m.
\end{itemize}
\end{aretenir}

Maintenant que tu sais calculer la capacité d'un condensateur plan à partir de sa géométrie, une question pratique se pose : que se passe-t-il lorsqu'on \emph{associe} plusieurs condensateurs dans un circuit, en série ou en parallèle ? C'est l'objet de la section suivante.

\coursec{Associations de condensateurs}

Dans la section précédente, nous avons vu que la capacité d'un condensateur plan est fixée par sa géométrie : $C = \varepsilon\,S/e$. En pratique, on ne fabrique pas un condensateur sur mesure pour chaque besoin : les condensateurs du commerce existent en valeurs normalisées (1 nF, 10 nF, 100 nF, 1 µF, 4,7 µF, 10 µF, 100 µF...). Comment obtenir alors la capacité exacte dont on a besoin dans un circuit ? La réponse est simple et puissante : on \cle{associe} plusieurs condensateurs, soit en \cle{parallèle}, soit en \cle{série}. C'est exactement ce que font les techniciens qui réparent les téléviseurs et les amplificateurs à Douala ou à Yaoundé : ils combinent des condensateurs de leur stock pour remplacer un composant défectueux. Toute cette section repose sur une seule relation fondamentale, déjà rencontrée : $q = C\,u$, que l'on écrira aussi $u = q/C$.

\courssub{Association en parallèle}

\textbf{Observation et intuition.} Deux condensateurs sont dits associés en \cle{parallèle} (ou en dérivation) lorsque leurs armatures sont reliées deux à deux : les deux armatures « gauches » sont connectées au même nœud, les deux armatures « droites » à l'autre nœud. Intuitivement, tout se passe comme si l'on avait réuni les armatures de même signe : on a fabriqué un condensateur dont la surface des armatures est la somme des surfaces. Or, pour un condensateur plan, $C = \varepsilon S/e$ augmente avec $S$. On s'attend donc à ce que la capacité équivalente soit \emph{plus grande} que chacune des capacités. C'est exactement ce que nous allons démontrer.

\begin{popfigure}
\begin{center}
\begin{circuitikz}[european]
\draw (0,0) to[short, i=$i$] (1,0);
\draw (1,0) -- (1,1.2) to[C, l_=$C_1$, i>_=$i_1$] (4,1.2);
\draw (1,0) -- (1,-1.2) to[C, l_=$C_2$, i>_=$i_2$] (4,-1.2);
\draw (4,1.2) -- (4,0) -- (4,-1.2);
\draw (4,0) -- (5.2,0);
\draw (0,0) -- (0,-2.6) to[sV=$u$] (5.2,-2.6) -- (5.2,0);
\node[popblue] at (2.5,2.0) {$q_1 = C_1 u$};
\node[poporange] at (2.5,-2.0) {$q_2 = C_2 u$};
\end{circuitikz}
\end{center}
\legende{Deux condensateurs en parallèle : la même tension $u$ règne aux bornes de $C_1$ et de $C_2$ ; les charges $q_1$ et $q_2$ s'ajoutent.}
\end{popfigure}

\begin{propriete}[Association en parallèle : démonstration exigible]
Deux condensateurs $C_1$ et $C_2$ montés en parallèle sont soumis à la \textbf{même tension} $u$ (propriété des dérivations). Le générateur débite une charge totale $q$ qui se répartit sur les deux condensateurs : par \textbf{conservation de la charge électrique},
\[ q = q_1 + q_2. \]
Or $q_1 = C_1\,u$ et $q_2 = C_2\,u$, donc :
\[ q = C_1 u + C_2 u = (C_1 + C_2)\,u. \]
En comparant avec la définition du condensateur équivalent, $q = C_{\text{éq}}\,u$, on obtient :
\[ \boxed{C_{\text{éq}} = C_1 + C_2} \]
Généralisation à $n$ condensateurs en parallèle : $C_{\text{éq}} = \displaystyle\sum_{i=1}^{n} C_i$.
\end{propriete}

\begin{aretenir}[En parallèle]
En parallèle : \textbf{même tension}, les \textbf{charges s'ajoutent}, les \textbf{capacités s'ajoutent}. La capacité équivalente est toujours \emph{supérieure} à la plus grande des capacités. On utilise le montage parallèle pour \cle{augmenter la capacité} d'un circuit.
\end{aretenir}

\courssub{Association en série}

\textbf{Observation et intuition.} Deux condensateurs sont associés en \cle{série} lorsqu'ils sont branchés l'un à la suite de l'autre, sans nœud intermédiaire. Imaginons que le générateur dépose une charge $+q$ sur l'armature d'entrée du premier condensateur. Par influence électrostatique, son autre armature porte $-q$. Mais cette armature est reliée par un fil à l'armature d'entrée du second condensateur : cet ensemble fil + deux armatures forme un \cle{îlot conducteur isolé}, initialement neutre. La charge totale de l'îlot doit rester nulle : l'armature d'entrée du second condensateur porte donc nécessairement $+q$, et sa seconde armature $-q$. Conclusion remarquable : \textbf{les deux condensateurs portent la même charge} $q$, quelles que soient leurs capacités.

\begin{popfigure}
\begin{center}
\begin{circuitikz}[european]
\draw (0,0) to[short, i=$i$] (1,0) to[C, l_=$C_1$, v=$u_1$] (3.4,0) to[C, l_=$C_2$, v=$u_2$] (5.8,0);
\draw (0,0) -- (0,-2.4) to[sV=$u$] (5.8,-2.4) -- (5.8,0);
\node[popblue] at (2.2,0.9) {$+q$};
\node[popblue] at (3.1,-0.55) {$-q$};
\node[poporange] at (4.4,0.9) {$+q$};
\node[poporange] at (5.3,-0.55) {$-q$};
\draw[dashed, popmuted, rounded corners] (2.9,-0.9) rectangle (4.9,0.75);
\node[popmuted] at (3.9,-1.25) {\footnotesize îlot isolé : charge totale nulle};
\end{circuitikz}
\end{center}
\legende{Deux condensateurs en série : la conservation de la charge sur l'îlot d'armatures (en pointillés) impose la même charge $q$ sur chaque condensateur ; les tensions $u_1$ et $u_2$ s'ajoutent.}
\end{popfigure}

\begin{propriete}[Association en série : démonstration exigible]
Deux condensateurs $C_1$ et $C_2$ montés en série portent la \textbf{même charge} $q$ (conservation de la charge sur l'îlot d'armatures reliées). D'après la loi d'additivité des tensions, la tension totale est :
\[ u = u_1 + u_2. \]
Or $u_1 = \dfrac{q}{C_1}$ et $u_2 = \dfrac{q}{C_2}$, donc :
\[ u = \frac{q}{C_1} + \frac{q}{C_2} = q\left(\frac{1}{C_1} + \frac{1}{C_2}\right). \]
En comparant avec $u = \dfrac{q}{C_{\text{éq}}}$, on obtient :
\[ \boxed{\frac{1}{C_{\text{éq}}} = \frac{1}{C_1} + \frac{1}{C_2}} \]
Généralisation à $n$ condensateurs en série : $\dfrac{1}{C_{\text{éq}}} = \displaystyle\sum_{i=1}^{n} \dfrac{1}{C_i}$.
\end{propriete}

\begin{propriete}[Cas particulier : deux condensateurs en série]
Pour deux condensateurs en série, on peut écrire directement la formule « produit sur somme » :
\[ C_{\text{éq}} = \frac{C_1\,C_2}{C_1 + C_2}. \]
\emph{Démonstration :} $\dfrac{1}{C_{\text{éq}}} = \dfrac{1}{C_1} + \dfrac{1}{C_2} = \dfrac{C_2 + C_1}{C_1 C_2}$, d'où le résultat en prenant l'inverse.
\end{propriete}

\begin{attention}[Erreur fréquente : ne pas confondre avec les résistances]
C'est le piège classique du Baccalauréat ! Les règles sont \textbf{inversées} par rapport aux résistances :
\begin{itemize}
\item Résistances : en série $R_{\text{éq}} = R_1 + R_2$ ; en parallèle $\dfrac{1}{R_{\text{éq}}} = \dfrac{1}{R_1} + \dfrac{1}{R_2}$.
\item Condensateurs : en \textbf{parallèle} $C_{\text{éq}} = C_1 + C_2$ ; en \textbf{série} $\dfrac{1}{C_{\text{éq}}} = \dfrac{1}{C_1} + \dfrac{1}{C_2}$.
\end{itemize}
Ainsi, la formule « produit/somme » $C_{\text{éq}} = \dfrac{C_1 C_2}{C_1 + C_2}$ concerne la \textbf{série} pour les condensateurs (alors qu'elle concerne le parallèle pour les résistances). Vérifie toujours la cohérence : en série, $C_{\text{éq}}$ doit être \emph{plus petite} que la plus petite des capacités ; en parallèle, \emph{plus grande} que la plus grande.
\end{attention}

\begin{exoresolu}[Parallèle ou série : deux condensateurs, deux résultats]
On dispose de deux condensateurs de capacités $C_1 = \SI{4}{\micro\farad}$ et $C_2 = \SI{12}{\micro\farad}$, soumis à une tension $u = \SI{24}{\volt}$.
\begin{enumerate}
\item On les associe en parallèle. Calculer $C_{\text{éq}}$, puis la charge de chaque condensateur et la charge totale.
\item On les associe en série. Calculer $C_{\text{éq}}$, la charge commune $q$, puis les tensions $u_1$ et $u_2$.
\end{enumerate}
\tcblower
\textbf{1. Association en parallèle.} Les capacités s'ajoutent :
\[ C_{\text{éq}} = C_1 + C_2 = 4 + 12 = \SI{16}{\micro\farad}. \]
Chaque condensateur est soumis à $u = \SI{24}{\volt}$ :
\[ q_1 = C_1 u = 4\times 10^{-6} \times 24 = \SI{96}{\micro\coulomb}, \qquad q_2 = C_2 u = 12\times 10^{-6} \times 24 = \SI{288}{\micro\coulomb}. \]
Vérification : $q = q_1 + q_2 = 96 + 288 = \SI{384}{\micro\coulomb}$ et $C_{\text{éq}}\,u = 16\times 10^{-6}\times 24 = \SI{384}{\micro\coulomb}$. Cohérent.

\textbf{2. Association en série.} Formule produit/somme :
\[ C_{\text{éq}} = \frac{C_1 C_2}{C_1 + C_2} = \frac{4\times 12}{4 + 12} = \frac{48}{16} = \SI{3}{\micro\farad}. \]
On remarque que $C_{\text{éq}} = \SI{3}{\micro\farad} < C_1 = \SI{4}{\micro\farad}$ : la capacité équivalente est bien plus petite que la plus petite des capacités. La charge commune est :
\[ q = C_{\text{éq}}\,u = 3\times 10^{-6}\times 24 = \SI{72}{\micro\coulomb}. \]
Les tensions se répartissent ainsi :
\[ u_1 = \frac{q}{C_1} = \frac{72\times 10^{-6}}{4\times 10^{-6}} = \SI{18}{\volt}, \qquad u_2 = \frac{q}{C_2} = \frac{72\times 10^{-6}}{12\times 10^{-6}} = \SI{6}{\volt}. \]
Vérification : $u_1 + u_2 = 18 + 6 = \SI{24}{\volt} = u$. Cohérent. On observe que le \textbf{plus petit condensateur} ($C_1$) supporte la \textbf{plus grande tension} ($\SI{18}{\volt}$).
\end{exoresolu}

\courssub{Répartition de la tension en série et tension de claquage}

\begin{propriete}[Répartition de la tension en série]
En série, la charge $q$ est commune, donc $u_i = \dfrac{q}{C_i}$ : la tension aux bornes d'un condensateur est \textbf{inversement proportionnelle} à sa capacité. Le plus petit condensateur supporte la plus grande tension. Pour deux condensateurs :
\[ \frac{u_1}{u_2} = \frac{C_2}{C_1}. \]
\end{propriete}

\begin{attention}[Tension de claquage : un danger réel]
Tout condensateur du commerce porte une indication de \cle{tension maximale} (ou tension de claquage), par exemple « 10 µF / 25 V ». Au-delà, le diélectrique claque et le condensateur est détruit. En série, ne te contente pas de vérifier que la tension totale est acceptable : \textbf{calcule la tension supportée par chaque condensateur}. Exemple : deux condensateurs de $\SI{10}{\micro\farad}$ / $\SI{25}{\volt}$ en série sous $\SI{48}{\volt}$ supportent chacun $\SI{24}{\volt}$ : c'est juste, mais admissible. Si l'un vaut $\SI{4}{\micro\farad}$ et l'autre $\SI{12}{\micro\farad}$, le premier supporte $\dfrac{3}{4}\times 48 = \SI{36}{\volt}$ : il claque ! L'intérêt du montage en série est précisément de \cle{supporter une tension plus élevée} qu'un condensateur seul, à condition de choisir des capacités adaptées.
\end{attention}

\begin{savaistu}[À quoi servent ces associations ?]
Dans les alimentations des postes radio et des téléviseurs, on place de gros condensateurs en parallèle pour obtenir de fortes capacités de filtrage (plusieurs milliers de µF), capables de lisser la tension redressée du réseau ENEO. À l'inverse, dans les circuits haute tension (flashs d'appareils photo, défibrillateurs des hôpitaux, lignes de transmission SONATREL), on monte des condensateurs en série pour répartir des tensions de plusieurs centaines de volts sans dépasser la tension de claquage de chacun.
\end{savaistu}

\courssub{Associations mixtes : méthode de réduction}

\begin{methode}[Réduire un circuit de condensateurs par étapes]
Pour déterminer la capacité équivalente d'une association mixte (série + parallèle) :
\begin{enumerate}
\item \textbf{Identifier} les groupements les plus simples : deux condensateurs sont en série s'ils sont parcourus par la même charge (pas de nœud entre eux) ; en parallèle s'ils sont soumis à la même tension (reliés aux deux mêmes nœuds).
\item \textbf{Remplacer} chaque groupement par son condensateur équivalent, en redessinant le schéma à chaque étape.
\item \textbf{Recommencer} jusqu'à n'obtenir qu'un seul condensateur équivalent.
\item \textbf{Remonter} ensuite le calcul pour déterminer charges et tensions de chaque condensateur, en appliquant : même tension en parallèle, même charge en série.
\end{enumerate}
\end{methode}

\begin{exoresolu}[Association mixte]
Trois condensateurs $C_1 = \SI{6}{\micro\farad}$, $C_2 = \SI{3}{\micro\farad}$ et $C_3 = \SI{2}{\micro\farad}$ sont montés ainsi : $C_2$ et $C_3$ en parallèle, l'ensemble en série avec $C_1$. Le tout est soumis à $u = \SI{30}{\volt}$. Déterminer la capacité équivalente, la charge de $C_1$ et la tension aux bornes du groupement parallèle.
\tcblower
\textbf{Étape 1 :} le groupement parallèle $C_2 \parallel C_3$ a pour capacité :
\[ C_{23} = C_2 + C_3 = 3 + 2 = \SI{5}{\micro\farad}. \]
\textbf{Étape 2 :} $C_{23}$ est en série avec $C_1$ :
\[ C_{\text{éq}} = \frac{C_1\,C_{23}}{C_1 + C_{23}} = \frac{6\times 5}{6 + 5} = \frac{30}{11} \approx \SI{2,73}{\micro\farad}. \]
\textbf{Étape 3 :} la charge débitée par le générateur est :
\[ q = C_{\text{éq}}\,u = \frac{30}{11}\times 10^{-6}\times 30 \approx \SI{81,8}{\micro\coulomb}. \]
$C_1$ est en série avec le reste : il porte toute cette charge, $q_1 = q \approx \SI{81,8}{\micro\coulomb}$.
\textbf{Étape 4 :} la tension aux bornes du groupement parallèle est :
\[ u_{23} = \frac{q}{C_{23}} = \frac{81{,}8\times 10^{-6}}{5\times 10^{-6}} \approx \SI{16,4}{\volt}. \]
Vérification : $u_1 = q/C_1 \approx \SI{13,6}{\volt}$ et $u_1 + u_{23} \approx 13{,}6 + 16{,}4 = \SI{30}{\volt} = u$. Cohérent.
\end{exoresolu}

\begin{aretenir}[Aide-mémoire : l'essentiel de la section]
\begin{itemize}
\item \textbf{Parallèle} : même tension $u$ ; les charges s'ajoutent ; $C_{\text{éq}} = C_1 + C_2 + \cdots$ (capacité augmentée).
\item \textbf{Série} : même charge $q$ (îlot d'armatures) ; les tensions s'ajoutent ; $\dfrac{1}{C_{\text{éq}}} = \dfrac{1}{C_1} + \dfrac{1}{C_2} + \cdots$ ; pour deux condensateurs, $C_{\text{éq}} = \dfrac{C_1 C_2}{C_1 + C_2}$.
\item En série, le plus petit condensateur supporte la plus grande tension : vérifier la tension de claquage de chacun.
\item Les règles sont \textbf{inverses} de celles des résistances.
\end{itemize}
\end{aretenir}

\begin{exercice}[À toi de jouer]
Un technicien dispose de trois condensateurs identiques de capacité $C = \SI{9}{\micro\farad}$ et de tension de claquage $\SI{50}{\volt}$ chacun.
\begin{enumerate}
\item Quelles capacités peut-il obtenir en les associant tous les trois ? (Envisager tous les montages possibles.)
\item Il souhaite réaliser un dipôle de capacité $\SI{6}{\micro\farad}$ pouvant fonctionner sous $\SI{80}{\volt}$. Proposer un montage et vérifier les tensions supportées.
\end{enumerate}
\end{exercice}

\begin{corrige}[Exercice : trois condensateurs identiques]
\textbf{1.} Quatre montages possibles :
\begin{itemize}
\item Trois en parallèle : $C_{\text{éq}} = 3C = \SI{27}{\micro\farad}$.
\item Trois en série : $\dfrac{1}{C_{\text{éq}}} = \dfrac{3}{C}$, soit $C_{\text{éq}} = \dfrac{C}{3} = \SI{3}{\micro\farad}$.
\item Deux en parallèle, en série avec le troisième : $C_{\text{éq}} = \dfrac{2C \times C}{2C + C} = \dfrac{2C}{3} = \SI{6}{\micro\farad}$.
\item Deux en série, en parallèle avec le troisième : $C_{\text{éq}} = \dfrac{C}{2} + C = \dfrac{3C}{2} = \SI{13,5}{\micro\farad}$.
\end{itemize}
\textbf{2.} Le montage « deux en parallèle ($\SI{18}{\micro\farad}$) en série avec le troisième ($\SI{9}{\micro\farad}$) » donne $C_{\text{éq}} = \SI{6}{\micro\farad}$. Sous $u = \SI{80}{\volt}$ : $q = C_{\text{éq}} u = 6\times 10^{-6}\times 80 = \SI{480}{\micro\coulomb}$. Le condensateur seul supporte $u_1 = q/C = 480/9 \approx \SI{53,3}{\volt} > \SI{50}{\volt}$ : il claque ! Le montage ne convient pas. En revanche, \textbf{deux en série} (chacun supporte $\SI{40}{\volt} < \SI{50}{\volt}$) \textbf{en parallèle avec le troisième} (qui supporte $\SI{80}{\volt}$ : claquage aussi !). Finalement, seul le montage \textbf{trois en série} supporte $\SI{80}{\volt}$ ($\approx \SI{26,7}{\volt}$ chacun), mais il ne donne que $\SI{3}{\micro\farad}$. \textbf{Avec seulement trois condensateurs, il est impossible d'obtenir $\SI{6}{\micro\farad}$ sous $\SI{80}{\volt}$}. Le montage $6\,\mu\text{F}$ (deux en parallèle, en série avec le troisième) ne supporte au maximum que $U_{\text{max}} = 50 \times 3/2 = \SI{75}{\volt}$ (le condensateur seul voit $2/3$ de la tension totale). Pour atteindre $\SI{6}{\micro\farad}$ sous $\SI{80}{\volt}$ en toute sécurité, il faudrait six condensateurs (deux branches de trois en série, mises en parallèle).
\end{corrige}

\coursec{Charge du condensateur dans un dipôle RC}

Dans la section précédente, nous avons appris à remplacer une association de condensateurs par un condensateur unique de capacité équivalente. Nous savons donc désormais calculer la charge $q = Cu$ d'un condensateur soumis à une tension donnée. Mais une question fondamentale reste ouverte : \emph{comment} un condensateur atteint-il cette charge ? Instantanément ? Progressivement ? En combien de temps ?

C'est toute la richesse du \cle{dipôle RC} : en intercalant un résistor entre le générateur et le condensateur, on contrôle la \emph{vitesse} de la charge. Ce circuit, d'une simplicité déconcertante (trois composants !), est pourtant au cœur des temporisateurs, des clignotants, des filtres et de la quasi-totalité des circuits électroniques qui nous entourent. Suivons la démarche complète : observation, modélisation, équation différentielle, solution, exploitation.

\courssub{Le circuit étudié : observation expérimentale}

\begin{experience}[Charge d'un condensateur : ce que l'on observe]
\textbf{Matériel :} un générateur de tension continue de fem $E$ (par exemple $E = \SI{6,0}{V}$), un résistor de résistance $R$, un condensateur de capacité $C$ initialement \textbf{déchargé} ($u = 0$), un interrupteur $K$, un voltmètre aux bornes du condensateur (ou une interface d'acquisition type ExAO).

\textbf{Protocole :} on ferme l'interrupteur à l'instant $t = 0$ et on enregistre la tension $u(t)$ aux bornes du condensateur.

\textbf{Observations :}
\begin{itemize}
\item à $t = 0$, le voltmètre indique $u = 0$ : le condensateur est déchargé ;
\item la tension $u$ \textbf{croît progressivement}, rapidement au début puis de plus en plus lentement ;
\item au bout d'un certain temps, $u$ se stabilise à la valeur $E$ : le condensateur est \textbf{chargé} ;
\item si l'on augmente $R$ ou $C$, la charge est \textbf{plus lente}.
\end{itemize}

\textbf{Interprétation :} la charge n'est \emph{pas} instantanée. Le résistor limite l'intensité du courant, donc le débit de porteurs de charge qui viennent s'accumuler sur les armatures. Le condensateur se remplit comme un réservoir d'eau alimenté par un robinet à débit limité : plus le tuyau est étroit (grand $R$) et plus le réservoir est grand (grand $C$), plus le remplissage est long.
\end{experience}

\begin{popfigure}
\begin{center}
\begin{circuitikz}[european, scale=1.0, transform shape]
\draw (0,0) to[V=$E$] (0,3)
      to[nos, l=$K$] (2,3)
      to[R=$R$, i=$i$] (5,3)
      to[C=$C$, v=$u$] (5,0)
      -- (0,0);
\end{circuitikz}
\end{center}
\legende{Circuit de charge d'un condensateur : générateur de fem $E$, interrupteur $K$, résistor $R$ et condensateur $C$ en série. La tension $u$ aux bornes du condensateur est fléchée en convention récepteur (flèche opposée au sens du courant $i$).}
\end{popfigure}

\begin{attention}[Convention d'orientation : le point de départ de tout le raisonnement]
Sur le schéma, le condensateur est orienté en \textbf{convention récepteur} : la flèche de la tension $u$ est \emph{opposée} au sens choisi pour le courant $i$, et l'armature reliée à l'entrée du courant porte la charge $+q$ avec $q = Cu$. C'est \emph{uniquement} dans cette convention que la relation $i = \dfrac{dq}{dt} = C\dfrac{du}{dt}$ est vraie \emph{sans signe moins}. Une erreur d'orientation fausse toute l'équation différentielle !
\end{attention}

\courssub{Établissement de l'équation différentielle}

Appliquons la \cle{loi des mailles} (loi d'additivité des tensions) au circuit, parcouru dans le sens du courant $i$ :
\[ E = u_R + u \]
où $u_R$ est la tension aux bornes du résistor et $u$ celle aux bornes du condensateur.

Deux relations constitutives entrent alors en jeu, et il est capital de ne pas les confondre :
\begin{itemize}
\item pour le \textbf{résistor}, la loi d'Ohm : $u_R = Ri$ ;
\item pour le \textbf{condensateur}, la relation courant-charge : $i = \dfrac{dq}{dt}$ avec $q = Cu$, d'où, $C$ étant une constante :
\[ i = C\,\dfrac{du}{dt} \]
\end{itemize}

En substituant dans la loi des mailles :
\[ E = R\,C\,\dfrac{du}{dt} + u \]

\begin{propriete}[Équation différentielle de la charge du condensateur]
La tension $u(t)$ aux bornes d'un condensateur de capacité $C$, chargé à travers un résistor de résistance $R$ par un générateur de fem $E$, vérifie l'équation différentielle linéaire du premier ordre à coefficients constants :
\[ \boxed{\;RC\,\dfrac{du}{dt} + u = E\;} \]
avec la condition initiale $u(0) = 0$ (condensateur initialement déchargé). \emph{Savoir établir cette équation à partir de la loi des mailles est exigible au Baccalauréat.}
\end{propriete}

\begin{attention}[L'erreur classique : écrire $i = u/R$ pour le condensateur]
Beaucoup d'élèves, par automatisme, appliquent la loi d'Ohm au condensateur. C'est \textbf{faux} : la loi d'Ohm $u_R = Ri$ ne s'applique qu'au \emph{résistor}. Le condensateur n'obéit pas à la loi d'Ohm : la relation qui lie courant et tension à ses bornes est $i = C\,\dfrac{du}{dt}$, une relation \emph{différentielle}. Retiens : le courant dans la branche est proportionnel à la \textbf{dérivée} de la tension du condensateur, pas à la tension elle-même. D'ailleurs, une fois la charge terminée, $u = E$ est constante, donc $\dfrac{du}{dt} = 0$ et $i = 0$ : le condensateur chargé se comporte comme un \textbf{interrupteur ouvert} en régime permanent continu.
\end{attention}

\courssub{Résolution de l'équation différentielle : la solution exponentielle}

L'équation $RC\,\dfrac{du}{dt} + u = E$ admet une solution de la forme $u(t) = A + B\,e^{-\alpha t}$. Plutôt que de l'admettre, \textbf{vérifions par substitution} que la fonction
\[ u(t) = E\left(1 - e^{-t/(RC)}\right) \]
est bien solution et satisfait la condition initiale. C'est la méthode attendue au Baccalauréat.

\begin{methode}[Vérifier qu'une fonction est solution de l'équation différentielle]
\begin{enumerate}
\item \textbf{Calculer la dérivée} de la fonction proposée :
\[ \dfrac{du}{dt} = E \times \left(-\frac{-1}{RC}\right) e^{-t/(RC)} = \frac{E}{RC}\,e^{-t/(RC)} \]
\item \textbf{Substituer} dans le membre de gauche $RC\,\dfrac{du}{dt} + u$ :
\begin{align*}
RC\,\dfrac{du}{dt} + u &= RC \times \frac{E}{RC}\,e^{-t/(RC)} + E\left(1 - e^{-t/(RC)}\right) \\
&= E\,e^{-t/(RC)} + E - E\,e^{-t/(RC)} = E
\end{align*}
On retrouve le membre de droite : la fonction est bien solution.
\item \textbf{Vérifier la condition initiale} : $u(0) = E(1 - e^{0}) = E(1-1) = 0$. Le condensateur est bien déchargé à $t = 0$.
\item \textbf{Conclure} : l'équation différentielle du premier ordre avec une condition initiale admet une solution \emph{unique} ; c'est donc \emph{la} solution du problème physique.
\end{enumerate}
\end{methode}

\begin{propriete}[Tension aux bornes du condensateur en charge]
La tension aux bornes du condensateur en charge dans un dipôle $RC$ soumis à un échelon de tension $E$ est :
\[ \boxed{\;u(t) = E\left(1 - e^{-t/\tau}\right) \quad \text{avec} \quad \tau = RC\;} \]
C'est une \textbf{croissance exponentielle} de $0$ vers la valeur asymptotique $E$.
\end{propriete}

\courssub{La constante de temps $\tau = RC$}

\begin{definition}[Constante de temps du dipôle RC]
On appelle \cle{constante de temps} du dipôle $RC$ la grandeur :
\[ \tau = RC \]
Elle caractérise la \textbf{rapidité} de la charge (ou de la décharge) : plus $\tau$ est grand, plus l'évolution est \textbf{lente} ; plus $\tau$ est petit, plus le condensateur se charge vite.
\end{definition}

\textbf{Analyse dimensionnelle (homogénéité).} Vérifions que $\tau$ est bien homogène à un temps. D'après la loi d'Ohm, $[R] = \dfrac{[U]}{[I]}$, et d'après $q = Cu$, $[C] = \dfrac{[Q]}{[U]}$. Donc :
\[ [RC] = \frac{[U]}{[I]} \times \frac{[Q]}{[U]} = \frac{[Q]}{[I]} = \frac{[I]\cdot[T]}{[I]} = [T] \]
La constante $\tau = RC$ s'exprime bien en \textbf{secondes} (s) si $R$ est en ohms ($\Omega$) et $C$ en farads (F). Cette vérification, rapide à rédiger, est très appréciée des correcteurs.

\textbf{Valeurs remarquables.} Évaluons $u$ aux instants caractéristiques :
\begin{itemize}
\item à $t = \tau$ : $u(\tau) = E\left(1 - e^{-1}\right) \approx E(1 - 0,368) \approx 0,63\,E$ : la charge atteint \textbf{63 \%} de sa valeur finale ;
\item à $t = 5\tau$ : $u(5\tau) = E\left(1 - e^{-5}\right) \approx E(1 - 0,0067) \approx 0,99\,E$ : on considère la charge \textbf{pratiquement terminée} (à 99 \%).
\end{itemize}

\begin{aretenir}[Les trois durées à connaître par cœur]
\begin{itemize}
\item $t = \tau$ : $u = 0,63\,E$ (charge à 63 \%) ;
\item $t = 5\tau$ : $u \approx 0,99\,E$ : \textbf{régime permanent atteint}, le condensateur est chargé ;
\item durée de charge totale : théoriquement infinie (asymptote), pratiquement égale à $5\tau$.
\end{itemize}
\end{aretenir}

\courssub{Étude graphique de la tension $u(t)$}

\begin{popfigure}
\begin{center}
\begin{tikzpicture}
\begin{axis}[
  width=0.95\linewidth, height=7cm,
  xlabel={$t$ (s)}, ylabel={$u$ (V)},
  xmin=0, xmax=6, ymin=0, ymax=7.2,
  xtick={1,5}, xticklabels={$\tau$,$5\tau$},
  ytick={3.78,6}, yticklabels={$0{,}63\,E$,$E$},
  grid=both, grid style={popline!40},
  axis lines=left, clip=false
]
\addplot[popcurve, domain=0:6, samples=200] {6*(1-exp(-x))};
\addplot[poporange, dashed, thick, domain=0:6] {6};
\addplot[popgreen, thick, domain=0:2.2] {6*x};
\addplot[popmuted, dotted] coordinates {(1,0) (1,3.78)};
\addplot[popmuted, dotted] coordinates {(0,3.78) (1,3.78)};
\addplot[popmuted, dotted] coordinates {(5,0) (5,5.96)};
\addplot[only marks, mark=*, poppink, mark size=2pt] coordinates {(1,3.78)};
\node[popdark, anchor=south west] at (axis cs:1.1,3.9) {\small $u(\tau) \approx 0,63\,E$};
\node[poporange, anchor=south west] at (axis cs:3.2,6.15) {\small asymptote $u = E$};
\node[popgreen, anchor=south east] at (axis cs:2.1,5.6) {\small tangente à l'origine};
\node[popblue, anchor=west] at (axis cs:3.1,2.6) {\small $u(t) = E\left(1 - e^{-t/\tau}\right)$};
\end{axis}
\end{tikzpicture}
\end{center}
\legende{Charge d'un condensateur ($E = \SI{6}{V}$, $\tau = \SI{1}{s}$) : croissance exponentielle vers l'asymptote $u = E$. La tangente à l'origine coupe l'asymptote en $t = \tau$ ; la charge est complète à 99 \% en $t = 5\tau$.}
\end{popfigure}

\textbf{Propriété géométrique remarquable : la tangente à l'origine.} Calculons la pente de la courbe en $t = 0$ :
\[ \left.\dfrac{du}{dt}\right|_{t=0} = \frac{E}{\tau}\,e^{0} = \frac{E}{\tau} \]
La tangente à l'origine a donc pour équation $u = \dfrac{E}{\tau}\,t$. Elle coupe l'asymptote $u = E$ lorsque $\dfrac{E}{\tau}\,t = E$, c'est-à-dire exactement en $t = \tau$. \emph{Si la charge se poursuivait à la vitesse initiale, elle serait terminée au bout d'une seule constante de temps.} Cette propriété fournit une méthode graphique de mesure de $\tau$.

\begin{methode}[Déterminer graphiquement la constante de temps $\tau$]
Deux méthodes équivalentes, à savoir rédiger au Bac :
\begin{enumerate}
\item \textbf{Méthode de la tangente à l'origine :} tracer la tangente à la courbe $u(t)$ au point $t = 0$ ; elle coupe l'asymptote horizontale $u = E$ en un point dont l'abscisse est $t = \tau$.
\item \textbf{Méthode des 63 \% :} calculer $0,63\,E$ ; repérer sur la courbe le point d'ordonnée $u = 0,63\,E$ ; son abscisse est $t = \tau$.
\end{enumerate}
Dans les deux cas, on en déduit ensuite $C = \dfrac{\tau}{R}$ si $R$ est connue : c'est une méthode \emph{expérimentale de mesure d'une capacité}.
\end{methode}

\courssub{Évolution de l'intensité du courant}

L'intensité se déduit immédiatement de $i = C\,\dfrac{du}{dt}$ :
\[ i(t) = C \times \frac{E}{RC}\,e^{-t/\tau} = \frac{E}{R}\,e^{-t/\tau} \]

\begin{propriete}[Intensité pendant la charge]
L'intensité du courant dans le circuit de charge décroît exponentiellement :
\[ i(t) = i_0\,e^{-t/\tau} \quad \text{avec} \quad i_0 = \frac{E}{R} \]
À $t = 0$, le condensateur déchargé se comporte comme un \textbf{fil} ($u = 0$, toute la tension $E$ est aux bornes de $R$, d'où $i_0 = E/R$) ; à $t \gg \tau$, il se comporte comme un \textbf{interrupteur ouvert} ($i \to 0$). L'intensité subit une \textbf{discontinuité} à la fermeture de $K$ (elle saute de $0$ à $i_0$), tandis que la tension $u$ aux bornes du condensateur est, elle, toujours \textbf{continue}.
\end{propriete}

\begin{popfigure}
\begin{center}
\begin{tikzpicture}
\begin{axis}[
  width=0.95\linewidth, height=6cm,
  xlabel={$t$ (s)}, ylabel={$i$ (mA)},
  xmin=0, xmax=6, ymin=0, ymax=0.75,
  xtick={1,5}, xticklabels={$\tau$,$5\tau$},
  ytick={0.221,0.6}, yticklabels={$0{,}37\,i_0$,$i_0 = \frac{E}{R}$},
  grid=both, grid style={popline!40},
  axis lines=left, clip=false
]
\addplot[popcurve, domain=0:6, samples=200] {0.6*exp(-x)};
\addplot[poporange, dashed, thick, domain=0:1.6] {0.6-0.6*x};
\addplot[popmuted, dotted] coordinates {(1,0) (1,0.221)};
\addplot[popmuted, dotted] coordinates {(0,0.221) (1,0.221)};
\addplot[only marks, mark=*, poppink, mark size=2pt] coordinates {(1,0.221)};
\node[poporange, anchor=south west] at (axis cs:1.05,0.05) {\small tangente à l'origine : coupe l'axe en $t = \tau$};
\node[popblue, anchor=west] at (axis cs:2.6,0.42) {\small $i(t) = \dfrac{E}{R}\,e^{-t/\tau}$};
\end{axis}
\end{tikzpicture}
\end{center}
\legende{Décroissance exponentielle de l'intensité pendant la charge ($i_0 = \SI{0,6}{mA}$, $\tau = \SI{1}{s}$). La tangente à l'origine coupe l'axe des temps en $t = \tau$ ; à $t = \tau$, l'intensité ne vaut plus que $0,37\,i_0$.}
\end{popfigure}

\courssub{Évolution de la charge $q(t)$ du condensateur}

Puisque $q = Cu$ à tout instant, la charge de l'armature positive suit la même loi que la tension, multipliée par $C$ :
\[ q(t) = CE\left(1 - e^{-t/\tau}\right) = q_{\max}\left(1 - e^{-t/\tau}\right) \]
où la \textbf{charge maximale} atteinte en régime permanent est :
\[ q_{\max} = CE \]
C'est cohérent avec la relation fondamentale $q = Cu$ appliquée au régime permanent où $u = E$. La courbe $q(t)$ a exactement la même allure que $u(t)$, avec $q(\tau) \approx 0,63\,q_{\max}$.

\courssub{Exemple chiffré et exercice résolu}

\begin{exemplebox}[Une constante de temps de une seconde]
Prenons $R = \SI{10}{k\Omega} = \SI{1,0e4}{\Omega}$ et $C = \SI{100}{\micro F} = \SI{1,0e-4}{F}$, avec $E = \SI{6,0}{V}$.
\[ \tau = RC = \num{1,0e4} \times \num{1,0e-4} = \SI{1,0}{s} \]
Conséquences concrètes :
\begin{itemize}
\item à $t = \SI{1,0}{s}$ : $u = 0,63 \times 6,0 \approx \SI{3,8}{V}$ ;
\item la charge est terminée (à 99 \%) à $t = 5\tau = \SI{5,0}{s}$ ;
\item l'intensité initiale vaut $i_0 = \dfrac{E}{R} = \dfrac{6,0}{\num{1,0e4}} = \SI{6,0e-4}{A} = \SI{0,60}{mA}$ ;
\item la charge maximale vaut $q_{\max} = CE = \num{1,0e-4} \times 6,0 = \SI{6,0e-4}{C} = \SI{0,60}{mC}$.
\end{itemize}
Une constante de temps de l'ordre de la seconde est parfaitement observable à l'œil nu sur un voltmètre : c'est ce qui rend ce circuit si précieux pour les temporisations.
\end{exemplebox}

\begin{exoresolu}[Exploitation complète d'une charge de condensateur]
\textbf{Énoncé.} Un condensateur de capacité $C$ inconnue, initialement déchargé, est chargé à travers un résistor $R = \SI{2,0}{k\Omega}$ par un générateur de fem $E = \SI{12}{V}$. L'enregistrement de $u(t)$ montre que la tension atteint $\SI{7,6}{V}$ à l'instant $t = \SI{4,0}{ms}$. (1) Déterminer $\tau$ puis $C$. (2) Donner l'expression de $i(t)$ et calculer $i_0$. (3) À quel instant la tension vaut-elle $\SI{11,9}{V}$ ? Conclure.
\tcblower
\textbf{Solution.} (1) On remarque que $\dfrac{7,6}{12} \approx 0,63$ : la tension atteint 63 \% de sa valeur finale à $t = \SI{4,0}{ms}$, donc $\tau = \SI{4,0}{ms}$. Alors :
\[ C = \frac{\tau}{R} = \frac{\num{4,0e-3}}{\num{2,0e3}} = \SI{2,0e-6}{F} = \SI{2,0}{\micro F} \]
(2) $i(t) = \dfrac{E}{R}\,e^{-t/\tau}$ avec $i_0 = \dfrac{12}{\num{2,0e3}} = \SI{6,0e-3}{A} = \SI{6,0}{mA}$.
(3) On résout $E\left(1 - e^{-t/\tau}\right) = \SI{11,9}{V}$ :
\[ e^{-t/\tau} = 1 - \frac{11,9}{12} = \num{8,3e-3} \quad\Rightarrow\quad t = \tau\,\ln\!\left(\frac{1}{\num{8,3e-3}}\right) \approx 4,8\,\tau \approx \SI{19}{ms} \]
On vérifie que $\dfrac{11,9}{12} \approx 0,99$ : la tension atteint 99 \% de $E$ autour de $t \approx 5\tau = \SI{20}{ms}$, en parfait accord avec le critère des $5\tau$.
\end{exoresolu}

\begin{savaistu}[Des clignotants aux feux de circulation : la temporisation RC partout]
Le clignotant d'une mobylette ou d'une mototaxi exploite le temps de charge d'un circuit $RC$ : tant que la tension aux bornes du condensateur n'a pas atteint un certain seuil, une lampe reste éteinte ; le seuil atteint, elle s'allume et le condensateur se décharge, puis le cycle recommence. La fréquence du clignotement est directement réglée par $\tau = RC$. De même, la temporisation de certains feux tricolores, l'essuie-vitre intermittent des voitures, la temporisation d'allumage des lampes de cage d'escalier et les circuits anti-rebond des boutons poussoirs de nos téléphones reposent tous sur la constante de temps d'un dipôle $RC$. Un simple résistor et un simple condensateur : et voilà le temps domestiqué !
\end{savaistu}

\begin{aretenir}[L'essentiel de la section]
\begin{itemize}
\item Loi des mailles + loi d'Ohm ($u_R = Ri$) + relation du condensateur ($i = C\,\dfrac{du}{dt}$) conduisent à l'équation différentielle $RC\,\dfrac{du}{dt} + u = E$.
\item Solution avec $u(0) = 0$ : $u(t) = E\left(1 - e^{-t/\tau}\right)$, avec la constante de temps $\tau = RC$ (en secondes).
\item Repères : $u(\tau) \approx 0,63\,E$ ; charge complète à $t \approx 5\tau$ ; tangente à l'origine coupant l'asymptote $u = E$ en $t = \tau$.
\item Intensité : $i(t) = \dfrac{E}{R}\,e^{-t/\tau}$, décroissante depuis $i_0 = E/R$ ; charge : $q(t) = CE\left(1 - e^{-t/\tau}\right)$, de maximum $q_{\max} = CE$.
\item La tension $u$ est continue ; l'intensité $i$ est discontinue à la fermeture de l'interrupteur.
\end{itemize}
\end{aretenir}

\coursec{Décharge du condensateur dans un résistor}

Dans la section précédente, nous avons chargé un condensateur à travers un résistor : branché sur un générateur de f.é.m. $E$, le condensateur a vu sa tension croître selon $u(t) = E\left(1 - e^{-t/\tau}\right)$ jusqu'à atteindre pratiquement la valeur $E$ au bout d'environ $5\tau$. Le condensateur est alors un véritable \cle{réservoir d'énergie électrique}. Une question naturelle se pose : que se passe-t-il si l'on supprime le générateur et que l'on referme le condensateur chargé sur le seul résistor ? C'est le phénomène de \cle{décharge}, que nous étudions maintenant. Tu le rencontres chaque jour sans le savoir : l'écran d'un téléphone qui s'éteint avec un léger délai, la lampe de poche qui faiblit progressivement, le flash d'appareil photo qui se « vide » entre deux photos.

\courssub{Le circuit étudié}

\begin{experience}[Circuit de charge puis de décharge]
On utilise un interrupteur à deux positions (K), un générateur de tension continue de f.é.m. $E$, un résistor de résistance $R$ et un condensateur de capacité $C$.
\begin{itemize}
\item \textbf{Position 1} : le condensateur se charge à travers $R$ sous la tension $E$. On attend assez longtemps ($t \gg 5\tau$) pour que la charge soit complète : $u = E$ et $q = CE$.
\item \textbf{Position 2} : à l'instant $t = 0$, on bascule l'interrupteur. Le générateur est hors circuit ; le condensateur chargé est refermé sur le résistor seul. On observe à l'oscilloscope (ou avec un voltmètre à mémoire) la décroissance de $u$ jusqu'à zéro.
\end{itemize}
\textbf{Observation} : la tension aux bornes du condensateur ne s'annule pas instantanément ; elle décroît progressivement, d'autant plus lentement que $R$ et $C$ sont grands. Un courant transitoire circule dans le résistor, \emph{dans le sens opposé} au courant de charge.
\end{experience}

\begin{popfigure}
\begin{center}
\begin{circuitikz}[european, scale=0.95, transform shape]
\draw (0,0) to[sV=$E$] (0,3);
\draw (0,3) -- (1,3);
\draw (1,3) -- (1.7,3.5);
\draw (1,3.5) node[left] {1};
\draw (1,2.5) node[left] {2};
\fill (1,3) circle (1.5pt);
\fill (2,3.5) circle (1.5pt);
\fill (2,2.5) circle (1.5pt);
\draw (2,3.5) -- (3,3.5) -- (3,3);
\draw (2,2.5) -- (3,2.5) -- (3,3);
\draw (1.05,3.08) -- (1.95,3.45);
\draw (1.5,3.9) node {K};
\draw (3,3) to[R=$R$, i=$i$, v=$u_R$] (6,3);
\draw (6,3) -- (6,1.5) to[C=$C$, v=$u$] (6,0);
\draw (6,0) -- (0,0);
\end{circuitikz}
\end{center}
\legende{Circuit à interrupteur deux positions : en position 1, charge du condensateur sous la tension $E$ ; en position 2, décharge à travers le résistor $R$. Le condensateur et le résistor sont orientés en convention récepteur.}
\end{popfigure}

\courssub{Équation différentielle de la décharge}

\begin{methode}[Établir l'équation différentielle de la décharge]
\begin{enumerate}
\item Flécher la tension $u$ aux bornes du condensateur et choisir le courant $i$ en \textbf{convention récepteur} pour le condensateur : la flèche de $i$ arrive sur l'armature où $u$ est comptée positivement. On a alors $i = \dfrac{dq}{dt} = C\dfrac{du}{dt}$.
\item Orienter le résistor aussi en convention récepteur : $u_R = Ri$.
\item Écrire la \textbf{loi des mailles} sur le circuit fermé condensateur--résistor : $u + u_R = 0$.
\item Remplacer $u_R$ par $Ri$ puis $i$ par $C\dfrac{du}{dt}$.
\end{enumerate}
\end{methode}

Appliquons cette méthode. Une fois l'interrupteur basculé en position 2, la maille ne contient plus que le condensateur et le résistor. La loi des mailles s'écrit :
\[ u + Ri = 0. \]
Le condensateur étant en convention récepteur, l'intensité est liée à sa tension par $i = \dfrac{dq}{dt}$ avec $q = Cu$, donc $i = C\dfrac{du}{dt}$. En reportant :
\[ u + RC\,\dfrac{du}{dt} = 0. \]

\begin{propriete}[Équation différentielle de la décharge]
Un condensateur de capacité $C$, initialement chargé sous la tension $E$, se déchargeant à travers un résistor de résistance $R$, obéit à l'équation différentielle :
\[ RC\,\dfrac{du}{dt} + u = 0 \qquad \text{avec la condition initiale } u(0) = E. \]
C'est une équation différentielle linéaire du premier ordre, à coefficients constants, sans second membre.
\end{propriete}

\begin{propriete}[Solution : loi de décharge]
La solution de l'équation précédente, vérifiant $u(0) = E$, est :
\[ u(t) = E\,e^{-t/\tau} \qquad \text{avec } \tau = RC \text{ la constante de temps (en secondes).} \]
\textbf{Vérification guidée (exigible)} : calculons la dérivée : $\dfrac{du}{dt} = -\dfrac{E}{\tau}e^{-t/\tau}$. Alors, puisque $\tau = RC$ :
\[ RC\,\dfrac{du}{dt} + u = RC\left(-\frac{E}{RC}e^{-t/\tau}\right) + E\,e^{-t/\tau} = -E\,e^{-t/\tau} + E\,e^{-t/\tau} = 0. \]
L'équation est bien vérifiée, et $u(0) = E\,e^{0} = E$ : la condition initiale est respectée.
\end{propriete}

\begin{attention}[Homogénéité de $\tau$]
Vérifie toujours que $\tau = RC$ est bien un temps : $[R] = \Omega = \mathrm{V\,A^{-1}}$ et $[C] = \mathrm{F} = \mathrm{C\,V^{-1}} = \mathrm{A\,s\,V^{-1}}$, donc $[RC] = \mathrm{V\,A^{-1}} \times \mathrm{A\,s\,V^{-1}} = \mathrm{s}$. L'exposant $-t/\tau$ est bien sans dimension, comme il se doit pour l'exponentielle.
\end{attention}

\courssub{Intensité du courant pendant la décharge}

L'intensité s'obtient en dérivant la tension :
\[ i(t) = C\,\dfrac{du}{dt} = C\left(-\frac{E}{RC}e^{-t/\tau}\right) = -\frac{E}{R}\,e^{-t/\tau}. \]

\begin{aretenir}[Intensité en décharge]
\[ i(t) = -\frac{E}{R}\,e^{-t/\tau}. \]
Le signe « $-$ » n'est pas une erreur : il signifie que le courant circule \textbf{en sens inverse} du sens choisi pendant la charge. Physiquement, les charges quittent l'armature positive du condensateur et traversent le résistor en sens contraire. À $t = 0^+$, l'intensité vaut $i(0^+) = -\dfrac{E}{R}$ en valeur algébrique, puis elle tend vers zéro.
\end{aretenir}

\begin{attention}[Erreur fréquente : le sens du courant]
Beaucoup d'élèves écrivent $i = \dfrac{E}{R}e^{-t/\tau}$ (positif) à la décharge par habitude de la charge. C'est faux \emph{avec l'orientation en convention récepteur} : le condensateur se comporte alors comme un générateur, le courant sort par son armature positive, donc $i < 0$ dans l'orientation choisie. Le signe de $i$ doit toujours être cohérent avec l'orientation fléchée sur ton schéma. En revanche, la tension $u$ aux bornes du condensateur reste positive pendant toute la décharge.
\end{attention}

\courssub{Allure des courbes et constante de temps}

\begin{popfigure}
\begin{center}
\begin{tikzpicture}
\begin{axis}[
  width=0.85\linewidth, height=6.5cm,
  axis lines=left,
  xlabel={$t$}, ylabel={$u(t)$},
  xmin=0, xmax=5.2, ymin=0, ymax=1.25,
  xtick={0,1,2,3,4,5}, xticklabels={$0$,$\tau$,$2\tau$,$3\tau$,$4\tau$,$5\tau$},
  ytick={0.368,1}, yticklabels={$0{,}37\,E$,$E$},
  grid=both, grid style={popline!40},
  legend style={at={(0.97,0.97)}, anchor=north east, draw=popline}
]
\addplot[popcurve, domain=0:5, samples=200] {exp(-x)};
\addlegendentry{$u(t) = E\,e^{-t/\tau}$}
\addplot[poporange, thick, dashed, domain=0:1.6, samples=20] {1-x};
\addlegendentry{tangente à l'origine}
\addplot[popmuted, dashed] coordinates {(1,0) (1,0.368)};
\addplot[popmuted, dashed] coordinates {(0,0.368) (1,0.368)};
\addplot[only marks, mark=*, poppink, mark size=2pt] coordinates {(1,0.368)};
\end{axis}
\end{tikzpicture}
\end{center}
\legende{Décharge du condensateur : décroissance exponentielle de $u(t)$. La tangente à l'origine coupe l'axe des temps en $t = \tau$ ; à $t = \tau$, la tension ne vaut plus que $0{,}37\,E$.}
\end{popfigure}

\begin{propriete}[Lecture graphique de $\tau$ et durée de décharge]
\begin{itemize}
\item La \textbf{tangente à l'origine} de la courbe $u(t)$ coupe l'axe des temps en $t = \tau$. En effet, la tangente en $t = 0$ a pour équation $y = E - \dfrac{E}{\tau}t$, qui s'annule en $t = \tau$.
\item À $t = \tau$ : $u(\tau) = E\,e^{-1} \approx 0{,}37\,E$ : la tension a perdu $63\,\%$ de sa valeur initiale.
\item À $t = 5\tau$ : $u(5\tau) = E\,e^{-5} \approx 0{,}007\,E$, soit moins de $1\,\%$ de $E$. On considère la décharge \textbf{pratiquement terminée au bout de $t \approx 5\tau$}.
\end{itemize}
La décharge est d'autant plus lente que $R$ (le courant de décharge est limité) et $C$ (le réservoir est grand) sont élevées.
\end{propriete}

\courssub{Continuité de la tension aux bornes du condensateur}

\begin{propriete}[Continuité de $u$ — complément utile au Bac]
La tension aux bornes d'un condensateur ne subit \textbf{jamais de discontinuité} : à l'instant du basculement de l'interrupteur, $u(0^-) = u(0^+) = E$. En effet, $u = q/C$ et la charge $q$ portée par les armatures ne peut pas varier instantanément : il faudrait un courant infini, ce qui est physiquement impossible. \emph{En revanche, l'intensité $i$ peut, elle, sauter brutalement} : ici elle passe de $0$ (fin de charge) à $-E/R$ (début de décharge). Cette propriété de continuité de $u$ est souvent utilisée au Baccalauréat pour justifier les conditions initiales.
\end{propriete}

\courssub{Interprétation énergétique}

D'où vient le courant qui traverse le résistor, alors qu'il n'y a plus de générateur ? Du \cle{réservoir d'énergie} qu'est le condensateur chargé. Pendant la décharge, l'énergie électrique emmagasinée $\mathcal{E} = \dfrac{1}{2}CE^{2}$ est progressivement restituée au circuit et intégralement dissipée par \cle{effet Joule} dans le résistor : le résistor chauffe légèrement. Quand $u \to 0$, l'énergie stockée tend vers zéro : tout a été converti en chaleur. Le condensateur joue, pendant la décharge, le rôle de \textbf{générateur temporaire} dont la « réserve » s'épuise.

\begin{exoresolu}[Décharge d'un condensateur : calculs complets]
\textbf{Énoncé.} Un condensateur de capacité $C = \SI{470}{\micro F}$ est chargé sous la tension $E = \SI{12}{V}$. À $t = 0$, on le connecte à un résistor de résistance $R = \SI{2}{k\Omega}$.
\begin{enumerate}
\item Calculer la constante de temps $\tau$ et la durée pratique de décharge.
\item Donner l'expression de $u(t)$ et calculer $u$ à $t = \tau$.
\item Donner l'expression de $i(t)$ et calculer l'intensité à $t = 0^+$.
\end{enumerate}
\tcblower
\textbf{Solution.}
\begin{enumerate}
\item Constante de temps :
\[ \tau = RC = \num{2000} \times \num{470e-6} = \SI{0,94}{s}. \]
La décharge est pratiquement terminée à $t \approx 5\tau = \SI{4,7}{s}$.
\item $u(t) = 12\,e^{-t/0{,}94}$ (en volts, $t$ en secondes). À $t = \tau = \SI{0,94}{s}$ :
\[ u(\tau) = 12\,e^{-1} \approx 12 \times 0{,}368 \approx \SI{4,4}{V}. \]
\item $i(t) = -\dfrac{E}{R}e^{-t/\tau} = -\dfrac{12}{2000}\,e^{-t/0{,}94} = -\num{6,0e-3}\,e^{-t/0{,}94}$ (en ampères). À $t = 0^+$ :
\[ i(0^+) = -\SI{6,0}{mA}. \]
Le signe négatif traduit l'inversion du sens du courant par rapport à la charge ; le résistor est traversé par un courant de $\SI{6,0}{mA}$ en valeur absolue, qui décroît ensuite exponentiellement.
\end{enumerate}
\end{exoresolu}

\begin{exoresolu}[Détermination expérimentale de la constante de temps]
\textbf{Énoncé.} On enregistre la décharge d'un condensateur initialement chargé sous $E = \SI{9,0}{V}$ à travers un résistor $R = \SI{10}{k\Omega}$. On relève $u = \SI{3,3}{V}$ à la date $t_1 = \SI{0,22}{s}$.
\begin{enumerate}
\item Montrer que $t_1$ correspond approximativement à la constante de temps $\tau$.
\item En déduire une valeur approchée de la capacité $C$.
\end{enumerate}
\tcblower
\textbf{Solution.}
\begin{enumerate}
\item Calculons le rapport : $\dfrac{u(t_1)}{E} = \dfrac{3{,}3}{9{,}0} \approx 0{,}37 \approx e^{-1}$. Or $u(\tau) = E\,e^{-1} \approx 0{,}37\,E$. La tension relevée correspond donc bien à $t_1 \approx \tau = \SI{0,22}{s}$.
\item De $\tau = RC$, on tire :
\[ C = \dfrac{\tau}{R} = \dfrac{0{,}22}{\num{10e3}} = \num{2,2e-5}\ \mathrm{F} = \SI{22}{\micro F}. \]
C'est une méthode classique de mesure d'une capacité au laboratoire.
\end{enumerate}
\end{exoresolu}

\begin{savaistu}[Pourquoi un appareil débranché peut encore être dangereux]
C'est la décharge \emph{lente} des condensateurs qui rend dangereux l'intérieur d'un téléviseur, d'un flash d'appareil photo ou d'un four à micro-ondes, \textbf{même débranché}. Les condensateurs de ces appareils sont chargés sous des tensions de plusieurs centaines, parfois plusieurs milliers de volts. S'ils ne sont reliés à aucun résistor de décharge (ou si $R$ est très grande), la constante de temps $\tau = RC$ peut atteindre des minutes, voire des heures : l'énergie $\frac{1}{2}Cu^{2}$ reste stockée longtemps après le débranchement. C'est pourquoi les techniciens déchargent volontairement les condensateurs avec un résistor avant toute intervention — et pourquoi il ne faut jamais ouvrir soi-même un tel appareil. À l'inverse, cette même propriété est mise à profit dans les flashs : on charge lentement le condensateur, puis on le décharge brutalement dans la lampe pour obtenir un éclair bref et intense. Les défibrillateurs utilisés dans les hôpitaux fonctionnent sur le même principe : un condensateur restitue en quelques millisecondes une énergie de plusieurs centaines de joules.
\end{savaistu}

\begin{aretenir}[L'essentiel de la décharge]
\begin{itemize}
\item Circuit : condensateur chargé ($u(0) = E$) refermé sur un résistor seul.
\item Équation différentielle : $RC\dfrac{du}{dt} + u = 0$.
\item Solutions : $u(t) = E\,e^{-t/\tau}$ et $i(t) = -\dfrac{E}{R}e^{-t/\tau}$, avec $\tau = RC$.
\item Repères : tangente à l'origine coupant l'axe des temps en $\tau$ ; $u(\tau) \approx 0{,}37\,E$ ; décharge terminée à $t \approx 5\tau$.
\item La tension $u$ est continue ; le courant $i$ s'inverse et peut être discontinu.
\item Bilan : l'énergie stockée $\frac{1}{2}CE^{2}$ est intégralement dissipée par effet Joule dans $R$.
\end{itemize}
\end{aretenir}

\coursec{Énergie emmagasinée par un condensateur}

Dans la section précédente, nous avons étudié la décharge d'un condensateur à travers une résistance : la tension $u_C$ décroît exponentiellement avec la constante de temps $\tau = RC$, et le courant circule dans le sens inverse du courant de charge. Mais une question fondamentale reste en suspens : d'où vient le courant qui traverse la résistance pendant la décharge, alors qu'aucun générateur n'est branché ? Quelque chose, dans le condensateur chargé, \emph{pousse} les charges et les fait circuler. Ce « quelque chose » est de l'\cle{énergie}, emmagasinée lors de la charge et restituée lors de la décharge. C'est toute la richesse du condensateur : c'est un \textbf{réservoir d'énergie électrique}. Cette dernière section quantifie cette énergie, établit ses expressions, et montre pourquoi le condensateur est devenu indispensable, du flash de ton téléphone au défibrillateur des hôpitaux de Yaoundé et de Douala.

\courssub{Mise en évidence expérimentale : le condensateur stocke de l'énergie}

\begin{experience}[Le condensateur fait briller une lampe]
\textbf{Matériel :} un condensateur de forte capacité ($C = 1\ \mathrm{F}$, tolérant $5{,}5\ \mathrm{V}$), un générateur de tension continue $E = 5\ \mathrm{V}$, une petite lampe à incandescence ou une DEL avec résistance de protection, un interrupteur à deux positions, des fils.

\textbf{Protocole :}
\begin{enumerate}
\item Placer l'interrupteur en position 1 : le condensateur se charge sous la tension $E$ du générateur.
\item Basculer l'interrupteur en position 2 : le générateur est débranché, le condensateur se retrouve seul en série avec la lampe.
\end{enumerate}

\textbf{Observations :} à la bascule, la lampe brille vivement un bref instant, puis son éclat décroît progressivement jusqu'à l'extinction, sur une durée de l'ordre de $\tau = RC$. Avec un condensateur chargé sous une tension plus élevée (montage de démonstration), on peut même observer une \textbf{étincelle} entre deux pointes reliées à ses armatures.

\textbf{Interprétation :} la lampe a émis de la lumière et de la chaleur. Or, par conservation de l'énergie, cette énergie ne peut pas apparaître ex nihilo : elle était \textbf{stockée dans le condensateur chargé}, sous forme d'énergie du champ électrique régnant entre ses armatures. La décharge n'est que la restitution de cette énergie.
\end{experience}

\begin{aretenir}[Idée clé]
Un condensateur de charge $q$ porté à la tension $u$ possède une \cle{énergie emmagasinée}, localisée dans le champ électrique entre les armatures. Cette énergie est restituée au circuit lors de la décharge.
\end{aretenir}

\courssub{Établissement de l'expression de l'énergie}

\textbf{Intuition.} Charger un condensateur, c'est transférer des charges d'une armature vers l'autre, contre le champ électrique qui s'oppose de plus en plus à ce transfert. Au début (condensateur déchargé, $u = 0$), transférer une charge ne coûte rien ; à la fin (tension $u$ élevée), chaque charge supplémentaire doit être « poussée » contre une tension déjà grande. Le travail total est donc une somme de petits travaux croissants : c'est exactement une intégrale.

\begin{propriete}[Énergie emmagasinée par un condensateur — démonstration exigible]
Un condensateur de capacité $C$, portant la charge $q$ sous la tension $u$, a emmagasiné l'énergie :
\[ \mathcal{E} = \frac{1}{2} C u^2 = \frac{1}{2} q\,u = \frac{q^2}{2C} \]
avec $\mathcal{E}$ en joules (J), $C$ en farads (F), $u$ en volts (V) et $q$ en coulombs (C).

\textbf{Démonstration.} Pendant la charge, le condensateur reçoit du reste du circuit la puissance électrique :
\[ p = u\,i \]
En convention récepteur, l'intensité qui \emph{entre} dans l'armature de charge $q$ est $i = \dfrac{dq}{dt}$, et comme $q = Cu$ avec $C$ constante, on a $i = C\dfrac{du}{dt}$. Donc :
\[ p = u \times C\,\frac{du}{dt} = C\,u\,\frac{du}{dt} \]
L'énergie reçue pendant la durée élémentaire $dt$ est $d\mathcal{E} = p\,dt = C\,u\,du$. L'énergie totale emmagasinée quand la tension passe de $0$ (condensateur déchargé) à la valeur finale $u$ s'obtient par intégration :
\[ \mathcal{E} = \int_0^{u} C\,u'\,du' = C \left[ \frac{u'^2}{2} \right]_0^{u} = \frac{1}{2} C u^2 \]
En utilisant $q = Cu$, on en déduit les deux formes équivalentes : $\mathcal{E} = \dfrac{1}{2}qu$ et $\mathcal{E} = \dfrac{q^2}{2C}$. \hfill $\blacksquare$
\end{propriete}

\textbf{Vérification par analyse dimensionnelle.} D'après $q = Cu$, on a $[C] = \mathrm{C \cdot V^{-1}}$, donc $[Cu^2] = \mathrm{C \cdot V}$. Or $1\ \mathrm{V} = 1\ \mathrm{J \cdot C^{-1}}$, d'où $[Cu^2] = \mathrm{C \cdot J \cdot C^{-1}} = \mathrm{J}$. L'expression est homogène à une énergie. 

\begin{popfigure}
\begin{center}
\begin{tikzpicture}
\begin{axis}[
  width=0.72\linewidth, height=6cm,
  axis lines=left,
  xlabel={$q$ (C)}, ylabel={$u$ (V)},
  xmin=0, xmax=1.25, ymin=0, ymax=1.25,
  xtick={1}, xticklabels={$q$},
  ytick={1}, yticklabels={$u$},
  clip=false]
  \addplot[popcurve,domain=0:1,samples=50] {x};
  \addplot[fill=popblue!25, draw=none] coordinates {(0,0) (1,0) (1,1)} \closedcycle;
  \addplot[popcurve,domain=0:1,samples=50] {x};
  \draw[dashed, popmuted] (axis cs:1,0) -- (axis cs:1,1) -- (axis cs:0,1);
  \node[popdark] at (axis cs:0.62,0.28) {$\mathcal{E} = \dfrac{1}{2} q\,u$};
  \node[popblue, rotate=38] at (axis cs:0.42,0.52) {$u = \dfrac{q}{C}$};
\end{axis}
\end{tikzpicture}
\end{center}
\legende{Interprétation graphique de l'énergie : dans le diagramme $u = f(q)$, l'énergie emmagasinée est l'aire du triangle sous la droite $u = q/C$, soit $\mathcal{E} = \tfrac{1}{2} q u$.}
\end{popfigure}

\begin{attention}[Le facteur $\tfrac{1}{2}$, piège classique du Bac]
L'erreur la plus fréquente est d'écrire $\mathcal{E} = qu$ ou $\mathcal{E} = Cu^2$ en oubliant le facteur $\dfrac{1}{2}$. Retiens l'image graphique : l'énergie est l'aire d'un \textbf{triangle} (et non d'un rectangle) car la tension \emph{croît} pendant la charge, de $0$ à $u$. Chaque charge élémentaire $dq$ n'a été transférée que sous la tension \emph{instantanée} $u'$, pas sous la tension finale. D'où la moyenne $\tfrac{1}{2}u$ et le facteur $\tfrac{1}{2}$. Vérifie aussi l'homogénéité de ton résultat : une énergie s'exprime toujours en joules.
\end{attention}

\courssub{Choisir la bonne expression selon les données}

Les trois expressions sont rigoureusement équivalentes, mais l'une est toujours plus pratique que les autres selon la grandeur qui est \emph{imposée} ou \emph{conservée} dans le problème.

\begin{methode}[Quelle expression choisir ?]
\begin{enumerate}
\item \textbf{Identifier la grandeur connue ou conservée.} Le problème donne-t-il la tension $u$, la charge $q$, ou les deux ?
\item \textbf{Si la tension est imposée} (condensateur branché en permanence sur un générateur, ou chargé sous une tension connue), utiliser :
\[ \mathcal{E} = \frac{1}{2} C u^2 \]
\item \textbf{Si la charge est conservée} (condensateur \emph{chargé puis isolé} : on débranche le générateur avant toute modification ; alors $q$ ne peut plus s'échapper), utiliser :
\[ \mathcal{E} = \frac{q^2}{2C} \]
C'est le cas typique lorsqu'on écarte les armatures d'un condensateur isolé : $q$ reste fixe, $C$ diminue, donc $\mathcal{E}$ augmente — le travail mécanique fourni pour écarter les armatures s'est transformé en énergie électrique.
\item \textbf{Si $q$ et $u$ sont toutes deux connues}, la forme $\mathcal{E} = \dfrac{1}{2}qu$ est la plus rapide.
\item \textbf{Vérifier} le résultat : unité en joules, ordre de grandeur plausible.
\end{enumerate}
\end{methode}

\begin{exemplebox}[Le flash d'un appareil photo]
Le condensateur d'un flash a pour capacité $C = 200\ \mu\mathrm{F} = \num{2,0e-4}\ \mathrm{F}$ et est chargé sous la tension $u = 300\ \mathrm{V}$. Calculons l'énergie disponible pour l'éclair :
\[ \mathcal{E} = \frac{1}{2} C u^2 = \frac{1}{2} \times \num{2,0e-4} \times (300)^2 = \frac{1}{2} \times \num{2,0e-4} \times \num{9,0e4} = 9{,}0\ \mathrm{J} \]
Ces $9\ \mathrm{J}$ sont libérés en une fraction de milliseconde dans le tube du flash : la puissance instantanée atteint plusieurs dizaines de kilowatts, ce qu'aucune pile ne pourrait fournir directement. C'est le principe même du condensateur : \textbf{accumuler lentement, restituer brutalement}.
\end{exemplebox}

\courssub{Bilan énergétique de la charge (complément utile au Bac)}

\begin{propriete}[Bilan énergétique de la charge d'un dipôle RC — complément]
Lors de la charge d'un condensateur (initialement déchargé) sous la tension constante $E$ à travers une résistance $R$ :
\begin{itemize}
\item le générateur fournit l'énergie $W_{\text{gén}} = q_{\max} E = C E^2$ ;
\item le condensateur n'en stocke que la moitié : $\mathcal{E}_C = \dfrac{1}{2} C E^2$ ;
\item l'autre moitié est dissipée par effet Joule dans la résistance : $W_R = \dfrac{1}{2} C E^2$.
\end{itemize}
Remarquablement, ce partage moitié-moitié est \textbf{indépendant de la valeur de $R$} : $R$ ne fixe que la \emph{durée} de la charge ($\tau = RC$), pas le bilan énergétique.
\end{propriete}

\textbf{Justification.} Le générateur débite un courant $i$ sous tension constante $E$ ; l'énergie fournie est $W_{\text{gén}} = \displaystyle\int_0^{+\infty} E\,i\,dt = E \int_0^{+\infty} i\,dt = E\,q_{\max} = CE^2$. Le condensateur stocke $\mathcal{E}_C = \tfrac{1}{2}CE^2$. La conservation de l'énergie impose alors $W_R = W_{\text{gén}} - \mathcal{E}_C = \tfrac{1}{2}CE^2$. (On peut le retrouver en calculant $\int_0^{+\infty} R i^2\,dt$ avec $i = \tfrac{E}{R}e^{-t/RC}$.)

\begin{popfigure}
\begin{center}
\begin{tikzpicture}[font=\small]
  \node[draw=popblue, fill=popblueL, rounded corners, minimum width=3.4cm, minimum height=1.1cm, align=center] (gen) at (0,0) {\textbf{Générateur}\\ $W_{\text{gén}} = CE^2$};
  \node[draw=popgreen, fill=popgreenL, rounded corners, minimum width=3.4cm, minimum height=1.1cm, align=center] (cond) at (5.2,1.3) {\textbf{Condensateur}\\ $\mathcal{E}_C = \dfrac{1}{2}CE^2$};
  \node[draw=poporange, fill=poporangeL, rounded corners, minimum width=3.4cm, minimum height=1.1cm, align=center] (res) at (5.2,-1.3) {\textbf{Résistance (Joule)}\\ $W_R = \dfrac{1}{2}CE^2$};
  \draw[-{Stealth[length=3mm]}, very thick, popgreen] (gen.east) -- (cond.west) node[midway, above, sloped, popdark] {stockée : $50\ \%$};
  \draw[-{Stealth[length=3mm]}, very thick, poporange] (gen.east) -- (res.west) node[midway, below, sloped, popdark] {dissipée : $50\ \%$};
\end{tikzpicture}
\end{center}
\legende{Bilan énergétique de la charge : l'énergie fournie par le générateur se partage en deux parts égales, quelle que soit la résistance $R$.}
\end{popfigure}

\courssub{Applications technologiques}

Le condensateur est le champion du \textbf{stockage temporaire} et de la \textbf{restitution rapide} d'énergie. Quelques applications emblématiques :

\begin{itemize}
\item \textbf{Flash photographique :} le condensateur se charge en quelques secondes depuis la batterie, puis se décharge en moins d'une milliseconde dans le tube à éclat.
\item \textbf{Défibrillateur cardiaque :} un condensateur restitue en quelques millisecondes un choc électrique calibré capable de resynchroniser les battements du cœur.
\item \textbf{Démarrage des moteurs monophasés} (réfrigérateurs, pompes à eau, climatiseurs très répandus dans les foyers camerounais) : un condensateur fournit l'impulsion de courant déphasé nécessaire au lancement du moteur.
\item \textbf{Alimentation de secours :} les condensateurs (et leurs cousins les supercondensateurs) maintiennent brièvement l'alimentation des mémoires et horloges des appareils électroniques lors d'une coupure — précieux quand le réseau ENEO connaît des délestages.
\end{itemize}

\begin{savaistu}[Le défibrillateur : la physique qui réanime]
Lors d'une fibrillation cardiaque, le cœur se contracte de façon chaotique et ne pompe plus le sang. Le défibrillateur, présent dans les services d'urgence des hôpitaux, charge un condensateur (de l'ordre de $100$ à $200\ \mu\mathrm{F}$) sous une tension pouvant dépasser $2000\ \mathrm{V}$, puis le décharge à travers le thorax du patient en $5$ à $10\ \mathrm{ms}$. L'énergie délivrée, typiquement de $100$ à $360\ \mathrm{J}$, traverse le cœur et « réinitialise » l'activité électrique des cellules cardiaques, permettant au stimulateur naturel de reprendre le rythme. Chaque seconde compte : la probabilité de survie chute d'environ $10\ \%$ par minute sans défibrillation. Derrière ce geste qui sauve des vies, il y a tout simplement... la formule $\mathcal{E} = \tfrac{1}{2}Cu^2$ !
\end{savaistu}

\begin{exoresolu}[Énergie d'un condensateur et choix de l'expression]
Un condensateur de capacité $C = 47\ \mu\mathrm{F}$ est chargé sous une tension $u = 12\ \mathrm{V}$.

\textbf{1.} Calculer sa charge $q$ et l'énergie emmagasinée $\mathcal{E}$.

\textbf{2.} On débranche le condensateur (il est désormais isolé), puis on modifie la géométrie de sorte que sa capacité soit divisée par deux. Que deviennent la tension et l'énergie ? Commenter.
\tcblower
\textbf{1.} Charge : $q = Cu = \num{47e-6} \times 12 = \num{5,64e-4}\ \mathrm{C}$, soit environ $0{,}56\ \mathrm{mC}$.

La tension étant connue, on utilise $\mathcal{E} = \dfrac{1}{2}Cu^2$ :
\[ \mathcal{E} = \frac{1}{2} \times \num{47e-6} \times 12^2 = \frac{1}{2} \times \num{47e-6} \times 144 \approx \num{3,4e-3}\ \mathrm{J} \]
soit $\mathcal{E} \approx 3{,}4\ \mathrm{mJ}$.

\textbf{2.} Le condensateur est \textbf{isolé} : sa charge ne peut pas s'échapper, donc $q$ est conservée. La capacité devient $C' = \dfrac{C}{2} = 23{,}5\ \mu\mathrm{F}$.

Nouvelle tension : $u' = \dfrac{q}{C'} = 2\,\dfrac{q}{C} = 2u = 24\ \mathrm{V}$.

Pour l'énergie, la grandeur conservée étant $q$, on utilise $\mathcal{E}' = \dfrac{q^2}{2C'}$ :
\[ \mathcal{E}' = \frac{q^2}{2C'} = 2\,\frac{q^2}{2C} = 2\mathcal{E} \approx 6{,}8\ \mathrm{mJ} \]

\textbf{Commentaire :} l'énergie a \emph{doublé} alors qu'aucun générateur n'est branché ! D'où vient ce supplément ? Du \textbf{travail mécanique} fourni par l'opérateur pour écarter les armatures : les charges opposées portées par les deux armatures s'attirent, il faut donc lutter contre cette attraction. Le travail fourni s'est intégralement converti en énergie électrique. La conservation de l'énergie est sauve — elle l'est toujours.
\end{exoresolu}

\begin{exercice}[Pour t'entraîner]
Un condensateur de $1000\ \mu\mathrm{F}$, initialement déchargé, est chargé sous $E = 24\ \mathrm{V}$ à travers une résistance.
\begin{enumerate}
\item Calculer l'énergie fournie par le générateur, l'énergie stockée par le condensateur et l'énergie dissipée dans la résistance.
\item Le condensateur chargé est ensuite déchargé dans une lampe. Quelle énergie maximale la lampe peut-elle recevoir ?
\end{enumerate}
\end{exercice}

\begin{corrige}[Exercice : bilan énergétique]
\textbf{1.} Énergie fournie : $W_{\text{gén}} = CE^2 = \num{1,0e-3} \times 24^2 = \num{1,0e-3} \times 576 = 0{,}576\ \mathrm{J}$.

Énergie stockée : $\mathcal{E}_C = \dfrac{1}{2}CE^2 = 0{,}288\ \mathrm{J}$.

Énergie dissipée dans la résistance : $W_R = W_{\text{gén}} - \mathcal{E}_C = 0{,}288\ \mathrm{J}$.

\textbf{2.} À la décharge, le condensateur restitue toute l'énergie qu'il a stockée : la lampe peut recevoir au maximum $\mathcal{E}_C = 0{,}288\ \mathrm{J}$ (convertie en lumière et en chaleur).
\end{corrige}

\begin{aretenir}[L'essentiel de la section]
\begin{itemize}
\item Un condensateur chargé emmagasine de l'énergie, restituée à la décharge : $\mathcal{E} = \dfrac{1}{2}Cu^2 = \dfrac{1}{2}qu = \dfrac{q^2}{2C}$ (en joules).
\item Le facteur $\dfrac{1}{2}$ vient de la croissance progressive de la tension pendant la charge (aire du triangle sous $u = q/C$).
\item Tension imposée $\Rightarrow$ $\mathcal{E} = \tfrac{1}{2}Cu^2$ ; charge conservée (condensateur isolé) $\Rightarrow$ $\mathcal{E} = \dfrac{q^2}{2C}$.
\item Lors de la charge sous tension constante $E$ : le générateur fournit $CE^2$, dont exactement la moitié est stockée et la moitié dissipée par effet Joule, quelle que soit $R$.
\item Applications : flash, défibrillateur, démarrage de moteurs, alimentation de secours.
\end{itemize}
\end{aretenir}

\newpage

\coursec{Activité d'intégration}

\courssub{Activité d'intégration : Le flash du photographe de quartier}

\textbf{Objectifs de l'activité}\par

À l'issue de cette activité d'intégration, tu seras capable de :

\begin{itemize}
\item exploiter la relation $q = C\,u$ pour déterminer la charge maximale d'un condensateur ;
\item calculer l'énergie emmagasinée par un condensateur à l'aide de la relation $E = \dfrac{1}{2}C\,u^2$ et la comparer à un besoin concret ;
\item utiliser la constante de temps $\tau = RC$ et le critère de charge complète ($t \approx 5\tau$) pour dimensionner une résistance ;
\item résoudre un problème simple d'association de condensateurs en parallèle et en série ;
\item tracer et interpréter l'allure de la courbe $u_C(t)$ pendant la charge ;
\item rédiger une synthèse argumentée sous forme de conseil technique, en mobilisant plusieurs savoir-faire de la leçon dans une situation authentique.
\end{itemize}

\textbf{Situation}\par

\textbf{Le flash du photographe de quartier.}

Monsieur Etoundi tient un petit studio photo au quartier Deïdo, à Douala. Pour ses séances de portraits et ses reportages de mariages, il utilise un flash électronique. Le principe en est simple : un \cle{condensateur} est chargé lentement, puis il se décharge très rapidement dans le tube à éclat, produisant une lumière intense pendant une fraction de seconde.

Son vieux flash vient de tomber en panne. Un technicien du marché Deïdo lui propose de fabriquer un circuit de recharge avec le matériel suivant :

\begin{itemize}
\item une pile de $\SI{9}{V}$ associée à un convertisseur élévateur qui délivre une tension continue $E = \SI{300}{V}$ ;
\item un condensateur de capacité $C = \SI{220}{\micro F}$ ;
\item un résistor de résistance $R$, à choisir, placé en série avec le condensateur.
\end{itemize}

Monsieur Etoundi impose deux contraintes de travail :

\begin{itemize}
\item \textbf{contrainte de puissance lumineuse :} le flash doit pouvoir délivrer une énergie d'au moins $E_{\min} = \SI{8}{J}$ à chaque éclair ;
\item \textbf{contrainte de rapidité :} entre deux prises de vue, le condensateur doit être rechargé à au moins $99\,\%$ de sa charge maximale en moins de $\SI{3}{s}$, pour ne pas faire attendre les clients.
\end{itemize}

On rappelle qu'un condensateur en charge dans un dipôle $RC$ est considéré comme pratiquement chargé ($u_C \geq 0{,}99\,E$) au bout d'une durée $t \approx 5\tau$, où $\tau = RC$ est la constante de temps du circuit.

Par ailleurs, le technicien ne dispose en stock que de condensateurs de capacité $C_0 = \SI{100}{\micro F}$ supportant une tension maximale de $\SI{350}{V}$. Il faudra donc vérifier s'il peut reconstituer une capacité équivalente convenable.

\textbf{Ta mission :} en tant qu'élève de Terminale C passionné de physique, tu es chargé d'étudier la faisabilité du projet et de rédiger un conseil technique à l'attention de Monsieur Etoundi.

\textbf{Consignes}\par

\begin{enumerate}
\item \textbf{Vérification énergétique.} Calculer l'énergie $E_C$ emmagasinée par le condensateur de $\SI{220}{\micro F}$ chargé sous $\SI{300}{V}$. La contrainte de puissance lumineuse est-elle satisfaite ?
\item \textbf{Charge maximale.} Calculer la charge maximale $Q_{\max}$ portée par l'armature positive du condensateur en fin de charge.
\item \textbf{Dimensionnement du résistor.} En utilisant le critère $5\tau \leq \SI{3}{s}$, déterminer la valeur maximale $R_{\max}$ de la résistance permettant de respecter la contrainte de rapidité. Choisir ensuite une valeur normalisée courante (par exemple $\SI{2,2}{k\Omega}$ ou $\SI{2,7}{k\Omega}$) et vérifier qu'elle convient.
\item \textbf{Association de condensateurs.} Montrer qu'en associant deux condensateurs de $C_0 = \SI{100}{\micro F}$ en parallèle, on obtient une capacité $C' = \SI{200}{\micro F}$. Recalculer l'énergie emmagasinée sous $\SI{300}{V}$ avec cette association : la contrainte de $\SI{8}{J}$ reste-t-elle satisfaite ? Justifier pourquoi l'association en parallèle (et non en série) est ici appropriée, sachant que chaque condensateur supporte au maximum $\SI{350}{V}$.
\item \textbf{Étude graphique.} Tracer l'allure de la courbe $u_C(t)$ pendant la charge sous $E = \SI{300}{V}$ avec la valeur de $R$ retenue. Placer sur le graphique la constante de temps $\tau$ et la valeur de $u_C$ à l'instant $t = \tau$.
\item \textbf{Synthèse.} Rédiger en cinq à dix lignes un conseil technique à Monsieur Etoundi : le projet est-il viable ? Quelles précautions doit-il prendre ?
\end{enumerate}

\textbf{Ressources à mobiliser}\par

\begin{itemize}
\item Relation charge--tension : $q = C\,u$, avec $q$ en coulombs (C), $C$ en farads (F), $u$ en volts (V).
\item Énergie emmagasinée : $E_C = \dfrac{1}{2}C\,u^2 = \dfrac{q^2}{2C} = \dfrac{1}{2}q\,u$, en joules (J).
\item Constante de temps du dipôle $RC$ : $\tau = RC$ (en secondes si $R$ est en ohms et $C$ en farads).
\item Charge du condensateur : équation différentielle $E = RC\,\dfrac{du_C}{dt} + u_C$ ; solution $u_C(t) = E\left(1 - e^{-t/\tau}\right)$ ; à $t = \tau$, $u_C \approx 0{,}63\,E$ ; à $t = 5\tau$, $u_C \approx 0{,}99\,E$ (charge quasi complète).
\item Associations : en parallèle, $C_{\text{éq}} = C_1 + C_2$ (la tension est commune) ; en série, $\dfrac{1}{C_{\text{éq}}} = \dfrac{1}{C_1} + \dfrac{1}{C_2}$ (la charge est commune).
\end{itemize}

\textbf{Démarche et corrigé commenté}\par

\begin{exoresolu}[Le flash du photographe de quartier]
On donne : $C = \SI{220}{\micro F} = \SI{2,20e-4}{F}$ ; $E = \SI{300}{V}$ ; $E_{\min} = \SI{8}{J}$ ; durée maximale de recharge à $99\,\%$ : $\SI{3}{s}$ ; condensateurs de stock : $C_0 = \SI{100}{\micro F}$ ($\SI{350}{V}$ max).
\tcblower
\textbf{1. Vérification énergétique.}

L'énergie emmagasinée par le condensateur chargé sous $E = \SI{300}{V}$ est :
\[
E_C = \frac{1}{2}C\,E^2 = \frac{1}{2} \times \num{2,20e-4} \times 300^2 = \frac{1}{2} \times \num{2,20e-4} \times \num{9,0e4} = \SI{9,9}{J}.
\]
Comparaison au besoin : $E_C = \SI{9,9}{J} > E_{\min} = \SI{8}{J}$. La contrainte de puissance lumineuse est \textbf{satisfaite}, avec une marge d'environ $24\,\%$.

\medskip
\textbf{2. Charge maximale.}

En fin de charge, $u_C = E$, donc :
\[
Q_{\max} = C\,E = \num{2,20e-4} \times 300 = \SI{6,6e-2}{C} = \SI{66}{mC}.
\]

\medskip
\textbf{3. Dimensionnement du résistor.}

\emph{Équation différentielle de la charge.} D'après la loi des mailles : $E = u_R + u_C$. Comme $u_R = R\,i$ et $i = C\,\dfrac{du_C}{dt}$, on obtient :
\[
E = RC\,\frac{du_C}{dt} + u_C \quad\Longleftrightarrow\quad \frac{du_C}{dt} + \frac{1}{RC}u_C = \frac{E}{RC}.
\]
La solution avec la condition initiale $u_C(0)=0$ est $u_C(t) = E\left(1 - e^{-t/\tau}\right)$ avec $\tau = RC$.

La charge est quasi complète ($99\,\%$) au bout de $5\tau$. La contrainte de rapidité s'écrit :
\[
5\tau \leq \SI{3}{s} \quad\Longleftrightarrow\quad 5RC \leq 3 \quad\Longleftrightarrow\quad R \leq \frac{3}{5C}.
\]
Application numérique :
\[
R_{\max} = \frac{3}{5 \times \num{2,20e-4}} = \frac{3}{\num{1,10e-3}} \approx \SI{2,73e3}{\Omega} = \SI{2,73}{k\Omega}.
\]
On choisit la valeur normalisée $R = \SI{2,2}{k\Omega}$ (inférieure à $R_{\max}$). Vérification :
\[
\tau = RC = \num{2,2e3} \times \num{2,20e-4} = \SI{0,484}{s} \quad\Rightarrow\quad 5\tau = \SI{2,42}{s} < \SI{3}{s}.
\]
La contrainte de rapidité est respectée : le flash est prêt en environ $\SI{2,4}{s}$.

\medskip
\textbf{4. Association de condensateurs.}

Deux condensateurs de $C_0 = \SI{100}{\micro F}$ montés \textbf{en parallèle} donnent :
\[
C' = C_0 + C_0 = \SI{200}{\micro F}.
\]
En parallèle, chaque condensateur est soumis à la même tension $\SI{300}{V}$, inférieure à la tension maximale supportée ($\SI{350}{V}$) : le montage est sûr. En série, la capacité équivalente ne vaudrait que $\SI{50}{\micro F}$, ce qui stockerait quatre fois moins d'énergie ($\SI{2,25}{J} < \SI{8}{J}$) : le choix série est donc à rejeter ici.

Nouvelle énergie emmagasinée sous $\SI{300}{V}$ :
\[
E'_C = \frac{1}{2}C'E^2 = \frac{1}{2} \times \num{2,00e-4} \times 300^2 = \SI{9,0}{J}.
\]
$E'_C = \SI{9,0}{J} > \SI{8}{J}$ : la contrainte énergétique reste satisfaite. Le nouveau temps de recharge avec $R = \SI{2,2}{k\Omega}$ : $\tau' = \num{2,2e3} \times \num{2,00e-4} = \SI{0,44}{s}$, soit $5\tau' = \SI{2,2}{s} < \SI{3}{s}$. Toutes les contraintes demeurent respectées.

\medskip
\textbf{5. Allure de la courbe $u_C(t)$.}

La tension suit la loi $u_C(t) = 300\left(1 - e^{-t/0{,}484}\right)$ (en volts). À $t = \tau = \SI{0,484}{s}$ : $u_C(\tau) = 300(1 - e^{-1}) \approx 0{,}63 \times 300 \approx \SI{190}{V}$.

\begin{popfigure}
\begin{center}
\begin{tikzpicture}
\begin{axis}[
  width=0.85\linewidth, height=6cm,
  xlabel={$t$ (s)}, ylabel={$u_C$ (V)},
  xmin=0, xmax=3, ymin=0, ymax=340,
  xtick={0,0.484,1,1.5,2,2.42,3},
  xticklabels={$0$,$\tau$,$1$,$1{,}5$,$2$,$5\tau$,$3$},
  ytick={190,300},
  yticklabels={$0{,}63E$,$E=300$},
  grid=both, grid style={popline!40},
]
\addplot[popcurve,domain=0:3,samples=200]{300*(1-exp(-x/0.484))};
\addplot[dashed,popmuted,domain=0:3]{300};
\draw[dashed,poporange] (axis cs:0.484,0) -- (axis cs:0.484,190) -- (axis cs:0,190);
\draw[dashed,popgreen] (axis cs:2.42,0) -- (axis cs:2.42,297);
\addplot[only marks,mark=*,poporange] coordinates {(0.484,190)};
\end{axis}
\end{tikzpicture}
\end{center}
\legende{Charge du condensateur sous $E = \SI{300}{V}$ ($R = \SI{2,2}{k\Omega}$, $C = \SI{220}{\micro F}$). À $t = \tau$, la tension atteint $63\,\%$ de sa valeur finale ; à $t = 5\tau$, elle dépasse $99\,\%$.}
\end{popfigure}

\medskip
\textbf{6. Synthèse : conseil technique à Monsieur Etoundi.}

« Monsieur Etoundi, le projet du technicien est viable. Le condensateur de $\SI{220}{\micro F}$ chargé sous $\SI{300}{V}$ stocke $\SI{9,9}{J}$, ce qui dépasse les $\SI{8}{J}$ nécessaires à votre éclair. Avec un résistor de $\SI{2,2}{k\Omega}$, la recharge atteint $99\,\%$ en $\SI{2,4}{s}$ environ : vous pourrez enchaîner les prises de vue sans faire attendre vos clients. Si le condensateur de $\SI{220}{\micro F}$ est introuvable, deux condensateurs de $\SI{100}{\micro F}$ branchés en parallèle conviennent ($\SI{9,0}{J}$ stockés, recharge en $\SI{2,2}{s}$). Attention toutefois : un condensateur chargé conserve son énergie même appareil éteint ; ne touchez jamais ses bornes et déchargez-le à travers un résistor avant toute manipulation. »
\end{exoresolu}

\textbf{Bilan de l'activité}\par

Cette activité t'a permis de mobiliser l'ensemble des savoir-faire de la leçon dans une situation professionnelle authentique :

\begin{itemize}
\item la relation $q = C\,u$ a donné la charge maximale $Q_{\max} = \SI{66}{mC}$ ;
\item la relation $E_C = \dfrac{1}{2}C\,u^2$ a permis de valider le critère énergétique ($\SI{9,9}{J} \geq \SI{8}{J}$) : on remarque au passage que l'énergie croît avec le \emph{carré} de la tension, ce qui explique pourquoi les flashs utilisent des tensions élevées malgré des capacités modestes ;
\item l'équation différentielle $E = RC\,\dfrac{du_C}{dt} + u_C$ et sa solution $u_C(t) = E(1-e^{-t/\tau})$ ont modélisé la charge ; la constante de temps $\tau = RC$ et le critère $5\tau$ ont servi d'outil de \cle{dimensionnement} : c'est le physicien qui \emph{calcule} la résistance, et non le hasard du bricolage ;
\item l'association en parallèle a montré comment reconstituer une capacité avec les composants disponibles, tout en respectant la tension maximale supportée par chaque condensateur ;
\item le tracé de $u_C(t)$ a donné une image concrète de la charge exponentielle : rapide au début, de plus en plus lente ensuite.
\end{itemize}

Tu as également développé des savoir-être essentiels : rigueur dans les conversions d'unités ($\SI{}{\micro F}$ en F), esprit critique dans le choix d'une valeur normalisée, et sens de la communication technique dans la rédaction du conseil final. Ces compétences sont exactement celles attendues au Baccalauréat C dans les situations d'intégration portant sur les circuits électriques à évolution temporelle.

\newpage

\coursec{Méthodes et exercices résolus}

\coursec{Le condensateur : constitution, charge et capacité}

\begin{definition}[Condensateur -- Constitution et symbole]
Un \textbf{condensateur} est un dipôle passif constitué de deux conducteurs (appelés \emph{armatures}) séparés par un isolant (le \emph{diélectrique}). Il se charge sous l'effet d'une tension : l'armature reliée au pôle $+$ du générateur porte une charge $+q$, l'autre une charge $-q$. La charge $q$ du condensateur est, par définition, la charge de l'armature positive.
Son symbole normalisé (norme CEI) est deux traits parallèles ; pour un condensateur polarisé (chimique), l'un des traits est remplacé par un trait plus épais ou un $+$ (armature positive).
\end{definition}

\begin{propriete}[Relation fondamentale $q = C\,u_C$]
La charge $q$ (en coulombs, C) emmagasinée sur l'armature positive est proportionnelle à la tension $u_C$ (en volts, V) aux bornes du condensateur :
\[ q = C\,u_C \]
Le coefficient de proportionnalité $C$ est la \textbf{capacité} (en farads, F). $1~\text{F} = 1~\text{C}\cdot\text{V}^{-1}$.
\emph{Convention} : $u_C$ est fléchée en convention récepteur (de l'armature $+$ vers l'armature $-$). Le courant $i$ entrant par l'armature positive vérifie $i = \dfrac{dq}{dt} = C\,\dfrac{du_C}{dt}$.
\end{propriete}

\courssub{Fiche méthode 1 : Réaliser le schéma d'un circuit comportant un condensateur}

\begin{methode}[Schématiser un circuit avec condensateur]
\begin{enumerate}
\item \textbf{Identifier le rôle du condensateur} dans l'énoncé : charge (circuit avec générateur), décharge (circuit avec résistor seul), ou les deux (interrupteur à deux positions).
\item \textbf{Placer les dipôles en série} : générateur de fem $E$, interrupteur $K$, résistor $R$, condensateur $C$. Le condensateur se représente par deux traits parallèles ; l'armature reliée au pôle $+$ du générateur porte la charge $+q$.
\item \textbf{Choisir et flécher les grandeurs} : flèche du courant $i$ dans le sens réel pendant la charge (du $+$ du générateur vers l'armature $A$) ; flèches de tension en convention récepteur pour $u_C$ et $u_R$ (flèche opposée à $i$), en convention générateur pour $E$.
\item \textbf{Indiquer les charges} $+q$ et $-q$ sur les armatures, avec $q = q_A$ la charge de l'armature vers laquelle arrive $i$, de sorte que $i = \dfrac{dq}{dt}$.
\item \textbf{Vérifier la cohérence} : la loi des mailles doit pouvoir s'écrire directement : $E = u_R + u_C = Ri + u_C$.
\end{enumerate}
\begin{attention}[Réflexe]
Un schéma sans flèches de tension ni sens du courant est un schéma inutilisable : toute l'étude du dipôle RC repose sur les conventions choisies. Au Bac, le correcteur vérifie systématiquement la cohérence entre le schéma et l'équation différentielle.
\end{attention}
\end{methode}

\begin{exoresolu}[Circuit de charge et de décharge d'un flash d'appareil photo]
Le flash d'un appareil photo fonctionne grâce à un condensateur de capacité $C$ que l'on charge sous une tension $E = \num{6,0}$ V à travers un conducteur ohmique de résistance $R$ (interrupteur en position 1), puis que l'on décharge dans la lampe du flash (position 2).\par
\textbf{1.} Réaliser le schéma du circuit complet avec l'interrupteur à deux positions, en fléchant $i$, $u_C$ et $u_R$ pour la phase de charge.\par
\textbf{2.} Écrire la loi des mailles pendant la charge.
\tcblower
\textbf{1.} Schéma du circuit :
\begin{center}
\begin{circuitikz}[european]
\draw (0,0) to[battery1, v=$E$] (0,3) -- (1.5,3) to[nos, l=$K$] (3,3);
\draw (3,3.4) -- (5,3) ;
\draw (5,3) to[R=$R$, v=$u_R$] (8,3) to[C=$C$, v=$u_C$] (8,0) -- (0,0);
\draw (3,2.6) -- (3,1.2) to[R=$R'$, l_={flash}] (8,1.2);
\draw (2.2,3.6) node[above]{$i$};
\draw[->, popblue, thick] (1.6,3.5) -- (2.8,3.5);
\end{circuitikz}
\end{center}
\textbf{Conventions choisies} (position 1, charge) : le courant $i$ circule du pôle $+$ du générateur vers l'armature supérieure $A$ du condensateur ; $u_R$ et $u_C$ sont fléchées en convention récepteur (flèches opposées au sens de $i$). L'armature $A$ porte la charge $+q$, l'armature $B$ la charge $-q$, et on a bien $i = \dfrac{dq}{dt} > 0$ pendant la charge.\par
\textbf{2.} Loi des mailles (parcours dans le sens de $i$) : la somme des tensions aux bornes des dipôles récepteurs est égale à la fem du générateur :
\[ E = u_R + u_C \quad \text{soit, avec } u_R = Ri \text{ (loi d'Ohm)} : \quad E = Ri + u_C. \]
\textbf{Conclusion} : le circuit correctement fléché conduit à la relation fondamentale $E = Ri + u_C$, point de départ de l'équation différentielle de la charge.
\end{exoresolu}

\courssub{Fiche méthode 2 : Exploiter la relation $q = Cu$}

\begin{methode}[Exploiter la relation entre charge et tension]
\begin{enumerate}
\item \textbf{Écrire la relation fondamentale} $q = C\,u_C$, où $q$ est la charge de l'armature \emph{positive} (celle vers laquelle arrive le courant) et $u_C$ la tension fléchée en convention récepteur.
\item \textbf{Convertir les unités} avant tout calcul : capacités en farads ($1~\mu$F $= 10^{-6}$ F ; $1$ nF $= 10^{-9}$ F ; $1$ pF $= 10^{-12}$ F), charges en coulombs, tensions en volts.
\item \textbf{Exploiter la continuité} : la charge $q$ (donc la tension $u_C$) ne peut pas varier brutalement ; juste après une commutation, $u_C$ garde la valeur qu'elle avait juste avant.
\item \textbf{Relier courant et charge} : $i = \dfrac{dq}{dt} = C\,\dfrac{du_C}{dt}$. C'est la clé pour passer d'une équation en $q$ à une équation en $u_C$.
\item \textbf{Conclure avec unité et chiffres significatifs} : vérifier l'homogénéité ($[C]\times[u] = $ F$\cdot$V $=$ C).
\end{enumerate}
\end{methode}

\begin{exoresolu}[Sauvetage des données d'un compteur ENEO]
Certains compteurs électriques électroniques utilisés par ENEO comportent un condensateur qui maintient l'alimentation de la mémoire pendant les microcoupures. Ce condensateur a une capacité $C = \num{470}~\mu$F et est chargé sous la tension $U = \num{5,0}$ V.\par
\textbf{1.} Calculer la charge $q$ emmagasinée sur l'armature positive.\par
\textbf{2.} Pendant une coupure, le condensateur débite un courant supposé constant d'intensité $I = \num{2,0}~\mu$A dans le circuit mémoire. Pendant combien de temps $\Delta t$ la tension aux bornes du condensateur reste-t-elle supérieure à $U_{\min} = \num{2,5}$ V, seuil de fonctionnement de la mémoire ?
\tcblower
\textbf{1.} Conversion : $C = \num{470}~\mu\text{F} = \num{4,70e-4}$ F.\par
Relation fondamentale : $q = C\,U$.\par
Application numérique :
\[ q = \num{4,70e-4} \times \num{5,0} = \num{2,35e-3}~\text{C} \approx \num{2,4}~\text{mC}. \]
Vérification dimensionnelle : F$\times$V $=$ (C/V)$\times$V $=$ C. La charge de l'armature positive vaut $q \approx \num{2,4}$ mC.\par
\textbf{2.} La tension atteint le seuil $U_{\min}$ lorsque la charge a diminué de :
\[ \Delta q = C\,(U - U_{\min}) = \num{4,70e-4} \times (\num{5,0} - \num{2,5}) = \num{4,70e-4} \times \num{2,5} = \num{1,175e-3}~\text{C}. \]
Le courant étant supposé constant, $I = \dfrac{\Delta q}{\Delta t}$, donc :
\[ \Delta t = \dfrac{\Delta q}{I} = \dfrac{\num{1,175e-3}}{\num{2,0e-6}} = \num{5,9e2}~\text{s} \approx \num{9,8}~\text{min}. \]
\textbf{Conclusion} : le condensateur maintient la mémoire du compteur pendant environ $10$ minutes, largement suffisant pour traverser les microcoupures du réseau. La relation $q = Cu$ permet de relier directement la chute de tension à la charge fournie.
\end{exoresolu}

\coursec{Capacité du condensateur plan}

\begin{definition}[Condensateur plan -- Permittivité]
Un \textbf{condensateur plan} est constitué de deux armatures planes, parallèles, de surface $S$ (en m$^2$), séparées par une distance $e$ (en m) faible devant $\sqrt{S}$. L'espace entre les armatures est rempli d'un diélectrique de permittivité $\varepsilon$ (en F/m). La permittivité du vide est $\varepsilon_0 = \num{8,85e-12}$ F/m. Pour un diélectrique, $\varepsilon = \varepsilon_r \varepsilon_0$ où $\varepsilon_r$ (sans unité) est la permittivité relative (constante diélectrique).
\end{definition}

\begin{propriete}[Capacité du condensateur plan]
\[ C = \varepsilon\,\frac{S}{e} \]
\emph{Démonstration (exigible au Bac)} : Champ uniforme $E = u_C/e$ entre les armatures. Loi de Gauss : $\Phi = E S = q/\varepsilon \Rightarrow q = \varepsilon S E = \varepsilon S u_C/e$. Identification avec $q = C u_C$ donne la formule.
\emph{Analyse dimensionnelle} : $[\varepsilon S/e] = $ (F/m)$\cdot$m$^2$/m $=$ F.
\end{propriete}

\begin{methode}[Calculer la capacité d'un condensateur plan]
\begin{enumerate}
\item Identifier $S$ (en m$^2$), $e$ (en m) et le diélectrique (donc $\varepsilon_r$ ou $\varepsilon$).
\item Convertir toutes les grandeurs en SI.
\item Appliquer $C = \varepsilon_0 \varepsilon_r \dfrac{S}{e}$.
\item Vérifier l'ordre de grandeur : pour $S \sim 1$ cm$^2$, $e \sim 1$ mm, air ($\varepsilon_r=1$), $C \sim 1$ pF.
\end{enumerate}
\end{methode}

\begin{exemplebox}[Condensateur plan à diélectrique variable]
Un condensateur plan d'armatures carrées de côté $L = \num{10}$ cm est rempli d'huile de transformateur ($\varepsilon_r = 2,2$). L'épaisseur de la couche d'huile est $e = \num{1,0}$ mm. Calculer sa capacité.
\tcblower
$S = L^2 = \num{0,01}$ m$^2$, $e = \num{1,0e-3}$ m, $\varepsilon = 2,2 \times \num{8,85e-12} = \num{1,947e-11}$ F/m.
\[ C = \num{1,947e-11} \times \frac{\num{0,01}}{\num{1,0e-3}} = \num{1,947e-10}~\text{F} \approx \num{195}~\text{pF}. \]
\end{exemplebox}

\begin{savaistu}[Condensateurs au Cameroun : correction du facteur de puissance]
Sur le réseau de transport SONATREL (225 kV, 90 kV) et de distribution ENEO (30 kV, 0,4 kV), les lignes et les moteurs à induction (pompes agricoles, ventilateurs d'usines, compresseurs frigorifiques) absorbent du courant réactif, ce qui échauffe les câbles et fait chuter la tension. On installe des \textbf{banques de condensateurs} (souvent des condensateurs planaires à diélectrique polypropylène, immergés dans l'huile) en parallèle sur le réseau pour fournir le réactif localement. Cela relève le facteur de puissance $\cos\varphi$ vers 1, réduit les pertes Joule ($P = RI^2$) et évite les pénalités de facturation pour $\cos\varphi < 0,9$. Au barrage de Nachtigal (420 MW), les groupes de production sont équipés de condensateurs de synchronisation et de filtrage harmonique pour garantir la qualité de l'onde injectée sur le réseau interconnecté.
\end{savaistu}

\coursec{Associations de condensateurs}

\begin{propriete}[Associations série et parallèle]
\begin{itemize}
\item \textbf{Parallèle} : même tension $u$ aux bornes de chaque condensateur. Charges additives : $q = q_1 + q_2$. Capacité équivalente : $C_{\text{éq}} = C_1 + C_2 + \cdots$ ($C_{\text{éq}} > \max(C_i)$).
\item \textbf{Série} : même charge $q$ sur chaque condensateur (conservation de la charge sur l'îlot conducteur). Tensions additives : $u = u_1 + u_2$. Capacité équivalente : $\dfrac{1}{C_{\text{éq}}} = \dfrac{1}{C_1} + \dfrac{1}{C_2} + \cdots$ ($C_{\text{éq}} < \min(C_i)$). Pour deux : $C_{\text{éq}} = \dfrac{C_1 C_2}{C_1 + C_2}$.
\end{itemize}
\end{propriete}

\courssub{Fiche méthode 3 : Résoudre un problème d'association de condensateurs}

\begin{methode}[Associations série et parallèle]
\begin{enumerate}
\item \textbf{Identifier le type d'association} :
\begin{itemize}
\item \emph{Parallèle} : les condensateurs sont soumis à la \textbf{même tension} $u$ ; les charges s'ajoutent : $q = q_1 + q_2$.
\item \emph{Série} : les condensateurs portent la \textbf{même charge} $q$ (conservation de la charge sur l'îlot d'armatures) ; les tensions s'ajoutent : $u = u_1 + u_2$.
\end{itemize}
\item \textbf{Appliquer la formule de la capacité équivalente} :
\begin{itemize}
\item Parallèle : $C_{\text{éq}} = C_1 + C_2 + \cdots$ (la capacité équivalente est \emph{plus grande} que chaque capacité).
\item Série : $\dfrac{1}{C_{\text{éq}}} = \dfrac{1}{C_1} + \dfrac{1}{C_2} + \cdots$ (la capacité équivalente est \emph{plus petite} que la plus petite des capacités). Pour deux condensateurs : $C_{\text{éq}} = \dfrac{C_1 C_2}{C_1 + C_2}$.
\end{itemize}
\item \textbf{Traiter les associations mixtes} par blocs : réduire d'abord chaque bloc (série ou parallèle), puis recommencer jusqu'à un seul condensateur équivalent.
\item \textbf{Remonter aux grandeurs demandées} : une fois $C_{\text{éq}}$ connue, calculer $q = C_{\text{éq}}\,u$, puis redescendre dans le schéma en appliquant les conservations (même $u$ en parallèle, même $q$ en série).
\item \textbf{Vérifier} : en série, la tension la plus grande est aux bornes de la \emph{plus petite} capacité ($u = q/C$).
\end{enumerate}
\end{methode}

\begin{exoresolu}[Association mixte dans un poste radio]
Dans un circuit de poste radio, on associe trois condensateurs : $C_1 = \num{4,0}~\mu$F, $C_2 = \num{12}~\mu$F montés en série, l'ensemble étant branché en parallèle avec $C_3 = \num{2,0}~\mu$F. Le tout est soumis à une tension $U = \num{24}$ V.\par
\textbf{1.} Calculer la capacité équivalente $C_{\text{éq}}$ de l'association.\par
\textbf{2.} Calculer la charge $q_3$ de $C_3$ et la charge $q$ portée par chacun des condensateurs $C_1$ et $C_2$.\par
\textbf{3.} Calculer les tensions $U_1$ et $U_2$ aux bornes de $C_1$ et $C_2$. Vérifier la loi d'additivité des tensions.
\tcblower
\textbf{1.} \emph{Étape 1 : bloc série $C_1$--$C_2$.}
\[ \dfrac{1}{C_{12}} = \dfrac{1}{C_1} + \dfrac{1}{C_2} \quad\Rightarrow\quad C_{12} = \dfrac{C_1 C_2}{C_1 + C_2} = \dfrac{\num{4,0}\times\num{12}}{\num{4,0}+\num{12}} = \dfrac{48}{16} = \num{3,0}~\mu\text{F}. \]
\emph{Étape 2 : parallèle de $C_{12}$ et $C_3$.}
\[ C_{\text{éq}} = C_{12} + C_3 = \num{3,0} + \num{2,0} = \num{5,0}~\mu\text{F}. \]
\textbf{2.} $C_3$ est directement soumis à la tension $U$ (association parallèle) :
\[ q_3 = C_3\,U = \num{2,0e-6}\times 24 = \num{4,8e-5}~\text{C} = \num{48}~\mu\text{C}. \]
En série, $C_1$ et $C_2$ portent la même charge, égale à celle du bloc équivalent :
\[ q = C_{12}\,U = \num{3,0e-6}\times 24 = \num{7,2e-5}~\text{C} = \num{72}~\mu\text{C}. \]
\textbf{3.} Tensions aux bornes des condensateurs en série :
\[ U_1 = \dfrac{q}{C_1} = \dfrac{\num{7,2e-5}}{\num{4,0e-6}} = \num{18}~\text{V} \qquad ; \qquad U_2 = \dfrac{q}{C_2} = \dfrac{\num{7,2e-5}}{\num{12e-6}} = \num{6,0}~\text{V}. \]
Vérification : $U_1 + U_2 = 18 + \num{6,0} = 24$ V $= U$. La loi d'additivité des tensions est bien vérifiée.\par
\textbf{Conclusion} : $C_{\text{éq}} = \num{5,0}~\mu$F ; on remarque que la plus petite capacité ($C_1$) supporte la plus grande tension ($18$ V), point de contrôle classique au Bac.
\end{exoresolu}

\coursec{Charge du condensateur dans un dipôle RC}

\begin{experience}[Observation de la charge d'un condensateur]
\textbf{Matériel :} Générateur de tension continue $E$ (pile 9 V ou alim. stabilisée), résistance $R = \num{100}~$k$\Omega$, condensateur $C = \num{100}~\mu$F (électrochimique, respecter la polarité), interrupteur, multimètre (voltmètre) ou acquisition numérique (Arduino/ExAO).\par
\textbf{Protocole :} Monter le circuit série $E$--$K$--$R$--$C$. Fermer $K$ à $t=0$. Relever $u_C(t)$ à intervalles réguliers (ex. toutes les 5 s).\par
\textbf{Observations :} $u_C$ part de 0, croît rapidement au début, puis ralentit pour tendre vers $E$. Le courant $i$ (mesuré par la chute de tension sur $R$) part de $E/R$ et décroît vers 0. La courbe $u_C(t)$ est une exponentielle saturante.
\end{experience}

\courssub{Fiche méthode 4 : Établir et résoudre l'équation différentielle de la charge (dipôle RC)}

\begin{methode}[Charge d'un condensateur à travers un résistor]
\begin{enumerate}
\item \textbf{Écrire la loi des mailles} avec les conventions du schéma : $E = u_R + u_C = Ri + u_C$.
\item \textbf{Injecter les relations du condensateur} : $i = \dfrac{dq}{dt}$ et $q = C\,u_C$, d'où $i = C\,\dfrac{du_C}{dt}$.
\item \textbf{Obtenir l'équation différentielle} sous forme canonique :
\[ RC\,\dfrac{du_C}{dt} + u_C = E \qquad \text{ou} \qquad \tau\,\dfrac{du_C}{dt} + u_C = E \quad \text{avec } \tau = RC. \]
\item \textbf{Écrire la solution} en tenant compte de la condition initiale $u_C(0) = 0$ (condensateur initialement déchargé) :
\[ u_C(t) = E\left(1 - e^{-t/\tau}\right) \qquad ; \qquad i(t) = \dfrac{E}{R}\,e^{-t/\tau}. \]
\item \textbf{Exploiter la constante de temps} : $\tau = RC$ s'exprime en secondes (vérifier : $\Omega \times $ F $=$ s). À $t = \tau$, $u_C = 0{,}63\,E$ ; la charge est pratiquement terminée à $t = 5\tau$ ($u_C \approx 0{,}99\,E$). Graphiquement, $\tau$ est l'abscisse de l'intersection de la tangente à l'origine avec l'asymptote $u_C = E$.
\item \textbf{Vérifier les limites} : $u_C(0)=0$, $u_C(\infty) = E$ ; $i(0) = E/R$, $i(\infty) = 0$.
\end{enumerate}
\end{methode}

\begin{exoresolu}[Charge du condensateur d'un défibrillateur]
Un défibrillateur d'un hôpital de district charge un condensateur de capacité $C = \num{100}~\mu$F à travers une résistance $R = \num{2,0}~$k$\Omega$ sous une tension $E = \num{2000}$ V. Le condensateur est initialement déchargé.\par
\textbf{1.} Établir l'équation différentielle vérifiée par $u_C(t)$ pendant la charge.\par
\textbf{2.} Vérifier que $u_C(t) = E\left(1 - e^{-t/\tau}\right)$ est solution et donner l'expression de $\tau$ ; calculer sa valeur.\par
\textbf{3.} Calculer le temps nécessaire pour que la tension atteigne $u_C = \num{1900}$ V.\par
\textbf{4.} Calculer l'intensité du courant à $t = 0$.
\tcblower
\textbf{1.} Loi des mailles : $E = u_R + u_C = Ri + u_C$. Avec $i = \dfrac{dq}{dt} = C\dfrac{du_C}{dt}$ :
\[ RC\,\dfrac{du_C}{dt} + u_C = E. \]
\textbf{2.} Dérivons la solution proposée : $\dfrac{du_C}{dt} = \dfrac{E}{\tau}\,e^{-t/\tau}$. Injectons dans le membre de gauche :
\[ RC\cdot\dfrac{E}{\tau}\,e^{-t/\tau} + E\left(1 - e^{-t/\tau}\right) = E\left(\dfrac{RC}{\tau} - 1\right)e^{-t/\tau} + E. \]
En posant $\tau = RC$, le terme exponentiel s'annule et le membre de gauche vaut $E$ : l'équation est vérifiée. De plus $u_C(0) = E(1-1) = 0$ : la condition initiale est respectée.\par
Valeur de la constante de temps :
\[ \tau = RC = \num{2,0e3}\times\num{100e-6} = \num{0,20}~\text{s}. \]
Analyse dimensionnelle : $[RC] = \Omega\cdot\text{F} = \dfrac{\text{V}}{\text{A}}\times\dfrac{\text{C}}{\text{V}} = \dfrac{\text{C}}{\text{A}} = $ s. \par
\textbf{3.} On résout $E\left(1 - e^{-t/\tau}\right) = \num{1900}$ :
\[ e^{-t/\tau} = 1 - \dfrac{1900}{2000} = \num{0,050} \quad\Rightarrow\quad t = -\tau\ln(\num{0,050}) = \num{0,20}\times\ln 20 = \num{0,20}\times\num{3,00} \approx \num{0,60}~\text{s}. \]
\textbf{4.} À $t = 0$, le condensateur déchargé se comporte comme un fil ($u_C = 0$) : toute la tension $E$ est aux bornes de $R$ :
\[ i(0) = \dfrac{E}{R} = \dfrac{2000}{\num{2,0e3}} = \num{1,0}~\text{A}. \]
\textbf{Conclusion} : la charge à $95~\%$ ne prend que $\num{0,60}$ s (soit $3\tau$) ; le défibrillateur est opérationnel presque instantanément, ce qui est vital en réanimation.
\end{exoresolu}

\coursec{Décharge du condensateur dans un résistor}

\begin{experience}[Observation de la décharge d'un condensateur]
\textbf{Matériel :} Même montage que pour la charge, mais l'interrupteur $K$ est un inverseur à deux positions : Pos. 1 (charge sur $E$), Pos. 2 (décharge dans $R$ seul).\par
\textbf{Protocole :} Charger le condensateur (Pos. 1, attendre $t > 5\tau$). Basculer brutalement sur Pos. 2 à $t=0$. Relever $u_C(t)$.\par
\textbf{Observations :} $u_C$ part de $U_0 = E$ et décroît exponentiellement vers 0. Le courant $i$ inverse son sens par rapport à la charge (il sort de l'armature positive) : $i(t) = -\dfrac{U_0}{R}e^{-t/\tau}$. La tangente à l'origine coupe l'axe des temps en $t = \tau$.
\end{experience}

\courssub{Fiche méthode 5 : Étudier la décharge d'un condensateur dans un résistor}

\begin{methode}[Décharge dans un résistor]
\begin{enumerate}
\item \textbf{Écrire la loi des mailles} du circuit condensateur--résistor seul : $u_C + u_R = 0$, soit $u_C + Ri = 0$ (le condensateur joue le rôle de générateur).
\item \textbf{Obtenir l'équation différentielle} avec $i = C\dfrac{du_C}{dt}$ :
\[ RC\,\dfrac{du_C}{dt} + u_C = 0. \]
\item \textbf{Écrire la solution} avec la condition initiale $u_C(0) = U_0$ (tension de charge initiale) :
\[ u_C(t) = U_0\,e^{-t/\tau} \qquad ; \qquad i(t) = -\dfrac{U_0}{R}\,e^{-t/\tau} \]
(le signe de $i$ traduit l'inversion du courant par rapport à la charge).
\item \textbf{Exploiter $\tau = RC$} : à $t = \tau$, $u_C = 0{,}37\,U_0$ ; décharge quasi totale à $t = 5\tau$. La tangente à l'origine coupe l'axe des temps en $t = \tau$.
\item \textbf{Déterminer $\tau$ graphiquement} si une courbe expérimentale est donnée : méthode de la tangente à l'origine, ou lecture de l'abscisse où $u_C = 0{,}37\,U_0$.
\item \textbf{En déduire une inconnue} : $R = \tau/C$ ou $C = \tau/R$.
\end{enumerate}
\end{methode}

\begin{exoresolu}[Décharge et détermination expérimentale d'une capacité]
Au laboratoire du lycée, on charge un condensateur de capacité $C$ inconnue sous la tension $U_0 = \num{12}$ V, puis on le décharge dans un résistor de résistance $R = \num{1,0}~$k$\Omega$. L'enregistrement de $u_C(t)$ montre que la tangente à la courbe à l'origine coupe l'axe des abscisses en $t = \num{47}$ ms.\par
\textbf{1.} Établir l'équation différentielle de la décharge et donner sa solution.\par
\textbf{2.} Montrer que la tangente à l'origine coupe l'axe des temps en $t = \tau$.\par
\textbf{3.} En déduire la valeur de $C$.\par
\textbf{4.} Quelle est la valeur de $u_C$ à $t = \tau$ ? À quel instant la décharge est-elle terminée à $99~\%$ ?
\tcblower
\textbf{1.} Le circuit ne contient que $C$ et $R$ : la loi des mailles donne $u_C + Ri = 0$. Avec $i = C\dfrac{du_C}{dt}$ :
\[ RC\,\dfrac{du_C}{dt} + u_C = 0, \qquad \text{de solution } u_C(t) = U_0\,e^{-t/\tau} \text{ avec } \tau = RC, \]
car $u_C(0) = U_0$ et $\dfrac{du_C}{dt} = -\dfrac{U_0}{\tau}e^{-t/\tau}$ vérifie bien $RC\dfrac{du_C}{dt} = -u_C$.\par
\textbf{2.} Équation de la tangente à l'origine : $y = u_C(0) + u_C'(0)\,t = U_0 - \dfrac{U_0}{\tau}\,t$.\par
Elle coupe l'axe des temps ($y = 0$) pour $U_0\left(1 - \dfrac{t}{\tau}\right) = 0$, soit $t = \tau$. CQFD.\par
\textbf{3.} On lit $\tau = \num{47}$ ms $= \num{4,7e-2}$ s. Donc :
\[ C = \dfrac{\tau}{R} = \dfrac{\num{4,7e-2}}{\num{1,0e3}} = \num{4,7e-5}~\text{F} = \num{47}~\mu\text{F}. \]
\textbf{4.} À $t = \tau$ : $u_C(\tau) = U_0\,e^{-1} = 12\times\num{0,37} \approx \num{4,4}$ V.\par
Décharge à $99~\%$ : il reste $1~\%$ de $U_0$, soit $e^{-t/\tau} = \num{0,01}$, d'où :
\[ t = \tau\ln 100 = \num{4,7e-2}\times\num{4,6} \approx \num{0,22}~\text{s} \quad(\approx 5\tau). \]
\textbf{Conclusion} : la mesure de $\tau$ par la méthode de la tangente est un moyen simple et précis de déterminer une capacité inconnue ; ici $C = \num{47}~\mu$F.
\end{exoresolu}

\coursec{Énergie emmagasinée par un condensateur}

\begin{propriete}[Énergie électrostatique d'un condensateur]
L'énergie $\mathcal{E}$ (en joules, J) emmagasinée dans le champ électrique d'un condensateur de capacité $C$ chargé à la tension $u_C$ (charge $q$) est :
\[ \mathcal{E} = \dfrac{1}{2}\,C\,u_C^{\,2} = \dfrac{1}{2}\,q\,u_C = \dfrac{q^2}{2C}. \]
\emph{Démonstration (exigible au Bac)} : Puissance reçue $P = u_C\,i = u_C\,C\,\dfrac{du_C}{dt} = \dfrac{d}{dt}\left(\dfrac{1}{2}Cu_C^2\right)$. Intégration de $0$ à $t$ : $\mathcal{E}(t) = \dfrac{1}{2}Cu_C(t)^2$ (énergie nulle pour $u_C=0$).
\emph{Analyse dimensionnelle} : $[\tfrac12 Cu^2] = $ F$\cdot$V$^2 = $ C$\cdot$V $=$ J.
\end{propriete}

\courssub{Fiche méthode 6 : Écrire et exploiter l'énergie emmagasinée}

\begin{methode}[Énergie d'un condensateur]
\begin{enumerate}
\item \textbf{Choisir la formule adaptée aux données} :
\[ \mathcal{E} = \dfrac{1}{2}\,C\,u_C^{\,2} = \dfrac{1}{2}\,q\,u_C = \dfrac{q^2}{2C}. \]
\item \textbf{Convertir en unités SI} : $C$ en farads, $u_C$ en volts, $q$ en coulombs ; l'énergie s'obtient en joules.
\item \textbf{Pour une association de condensateurs} : calculer d'abord $C_{\text{éq}}$, puis $\mathcal{E} = \dfrac{1}{2}C_{\text{éq}}u^2$ ; ou sommer les énergies de chaque condensateur (l'énergie est additive).
\item \textbf{Exploiter un transfert d'énergie} : lors de la décharge dans un résistor, l'énergie emmagasinée est intégralement dissipée par effet Joule : $\mathcal{E}_{\text{Joule}} = \dfrac{1}{2}CU_0^2$.
\item \textbf{Vérifier l'homogénéité} : $[\tfrac12 Cu^2] = $ F$\cdot$V$^2 = $ J.
\end{enumerate}
\end{methode}

\begin{exoresolu}[Énergie du flash et du défibrillateur]
\textbf{A.} Le condensateur du flash d'un appareil photo ($C = \num{330}~\mu$F) est chargé sous $U = \num{300}$ V.\par
\textbf{1.} Calculer l'énergie emmagasinée $\mathcal{E}$ et la charge $q$ du condensateur.\par
\textbf{2.} La décharge dans le tube du flash dure $\Delta t = \num{1,0}$ ms. Calculer la puissance moyenne du flash.\par
\textbf{B.} Le condensateur d'un défibrillateur ($C = \num{100}~\mu$F) est chargé sous $U' = \num{2000}$ V. Lors du choc, $70~\%$ de l'énergie stockée est délivrée au patient. Calculer l'énergie reçue par le patient.
\tcblower
\textbf{A.1.} Énergie emmagasinée :
\[ \mathcal{E} = \dfrac{1}{2}CU^2 = \dfrac{1}{2}\times\num{330e-6}\times 300^2 = \dfrac{1}{2}\times\num{330e-6}\times\num{9,0e4} = \num{14,85}~\text{J} \approx \num{15}~\text{J}. \]
Charge : $q = CU = \num{330e-6}\times 300 = \num{9,9e-2}$ C.\par
Vérification croisée : $\dfrac{q^2}{2C} = \dfrac{(\num{9,9e-2})^2}{2\times\num{330e-6}} = \dfrac{\num{9,8e-3}}{\num{6,6e-4}} \approx \num{14,9}$ J. Cohérent.\par
\textbf{A.2.} Puissance moyenne :
\[ \mathcal{P}_{\text{moy}} = \dfrac{\mathcal{E}}{\Delta t} = \dfrac{\num{14,85}}{\num{1,0e-3}} \approx \num{1,5e4}~\text{W} = \num{15}~\text{kW}. \]
\textbf{B.} Énergie stockée :
\[ \mathcal{E}' = \dfrac{1}{2}C{U'}^2 = \dfrac{1}{2}\times\num{100e-6}\times (2000)^2 = \dfrac{1}{2}\times\num{100e-6}\times\num{4,0e6} = \num{200}~\text{J}. \]
Énergie délivrée au patient :
\[ \mathcal{E}_{\text{patient}} = \num{0,70}\times 200 = \num{140}~\text{J}. \]
\textbf{Conclusion} : le condensateur est un réservoir d'énergie capable de la restituer quasi instantanément : $15$ J en $1$ ms pour le flash (puissance de $15$ kW !), $140$ J pour le choc du défibrillateur, dose typique en réanimation cardiaque.
\end{exoresolu}

\begin{savaistu}[Le condensateur dans la vie quotidienne au Cameroun]
\begin{itemize}
\item \textbf{Démarrage des moteurs monophasés} (pompes à eau, ventilateurs, compresseurs de frigidaire) : un condensateur de démarrage (souvent $20$--$60~\mu$F, $450$ V) crée un déphasage pour générer un couple de démarrage. S'il est HS, le moteur « ronronne » sans tourner (surconsommation, risque de brûlure).
\item \textbf{Alimentations à découpage (chargeurs téléphone MTN/Orange, PC, TV)} : condensateurs chimiques $400$ V ($10$--$100~\mu$F) pour lisser la tension redressée ; condensateurs CMS céramiques (nF--µF) pour le filtrage haute fréquence.
\item \textbf{Correction du facteur de puissance} : dans les usines (aluminium, ciment, brasseries) et les bâtiments administratifs, des batteries de condensateurs (souvent $25$--$50$ kvar par palier) sont commutées par contacteurs pour maintenir $\cos\varphi > 0,92$ et éviter les pénalités ENEO.
\item \textbf{Paratonnerres et parafoudres} : les varistances (comportement non-linéaire) et les condensateurs à gaz protègent les entrées HT/BT des postes de transformation contre les surtensions de foudre (fréquentes en saison des pluies à Douala/Yaoundé).
\end{itemize}
\end{savaistu}

\begin{attention}[Les 5 erreurs qui coûtent des points au Bac]
\begin{enumerate}
\item \textbf{Oublier les conversions d'unités.} Les capacités sont presque toujours données en $\mu$F, nF ou pF : un calcul mené avec $C = 100$ au lieu de $\num{100e-6}$ F donne un résultat faux d'un facteur $10^6$. Écris systématiquement la conversion \emph{avant} l'application numérique.
\item \textbf{Inverser les formules des associations.} C'est le piège classique : pour les condensateurs, la formule « en inverse » ($\dfrac{1}{C_{\text{éq}}} = \sum\dfrac{1}{C_i}$) concerne l'association \textbf{série} (contrairement aux résistors !). Contrôle rapide : en série, $C_{\text{éq}}$ doit être plus petite que la plus petite capacité ; en parallèle, plus grande que la plus grande.
\item \textbf{Mal établir l'équation différentielle à cause des conventions.} Si $i = \dfrac{dq}{dt}$ avec $q$ charge de l'armature vers laquelle \emph{arrive} le courant, alors la charge donne $RC\dfrac{du_C}{dt} + u_C = E$ et la décharge $RC\dfrac{du_C}{dt} + u_C = 0$. Un schéma non fléché ou une convention récepteur non respectée conduit à un signe faux, et toute la suite est fausse.
\item \textbf{Confondre les valeurs remarquables de $\tau$.} Pendant la \emph{charge}, $u_C(\tau) = 0{,}63\,E$ ; pendant la \emph{décharge}, $u_C(\tau) = 0{,}37\,U_0$. Et dans les deux cas, le régime permanent est atteint à $t \approx 5\tau$, pas à $t = \tau$. Précise toujours de quelle phase tu parles.
\item \textbf{Écrire l'énergie sans le facteur $\dfrac{1}{2}$ ou avec une mauvaise variable.} $\mathcal{E} = \dfrac{1}{2}Cu^2 = \dfrac{q^2}{2C}$, et non $Cu^2$ ni $\dfrac{q^2}{C}$. Autre piège : utiliser $E$ (fem du générateur) au lieu de la tension $u_C$ réelle aux bornes du condensateur à l'instant considéré. Termine toujours par une vérification dimensionnelle : F$\cdot$V$^2 = $ J.
\end{enumerate}
\end{attention}

\newpage

\coursec{Exercices}

\courssub{Je vérifie mes connaissances}

\begin{exercice}[QCM : réflexes sur le condensateur]
Pour chaque question, une seule réponse est exacte. Recopie la lettre de la bonne réponse.
\begin{enumerate}
\item La relation entre la charge $q$ portée par une armature d'un condensateur et la tension $u$ à ses bornes est :
\begin{enumerate}
\item $q = \dfrac{u}{C}$ \qquad \item $q = C\,u$ \qquad \item $q = \dfrac{C}{u}$ \qquad \item $q = C\,u^2$
\end{enumerate}
\item L'unité de capacité, le farad, équivaut à :
\begin{enumerate}
\item $\SI{1}{C.V}$ \qquad \item $\SI{1}{V.C^{-1}}$ \qquad \item $\SI{1}{C.V^{-1}}$ \qquad \item $\SI{1}{A.V^{-1}}$
\end{enumerate}
\item Deux condensateurs de capacités $C_1$ et $C_2$ associés en parallèle équivalent à un condensateur de capacité :
\begin{enumerate}
\item $\dfrac{C_1 C_2}{C_1 + C_2}$ \qquad \item $C_1 + C_2$ \qquad \item $\dfrac{1}{C_1} + \dfrac{1}{C_2}$ \qquad \item $\sqrt{C_1 C_2}$
\end{enumerate}
\item La constante de temps d'un dipôle RC vaut :
\begin{enumerate}
\item $\tau = \dfrac{R}{C}$ \qquad \item $\tau = \dfrac{C}{R}$ \qquad \item $\tau = R\,C$ \qquad \item $\tau = \dfrac{1}{RC}$
\end{enumerate}
\item L'énergie emmagasinée par un condensateur chargé s'écrit :
\begin{enumerate}
\item $E = \dfrac{1}{2} C u$ \qquad \item $E = \dfrac{1}{2} C u^2$ \qquad \item $E = C u^2$ \qquad \item $E = \dfrac{q^2}{C}$
\end{enumerate}
\end{enumerate}
\end{exercice}

\begin{exercice}[Vrai ou faux, avec justification]
Réponds par vrai ou faux en justifiant chaque réponse par une phrase ou un calcul.
\begin{enumerate}
\item La capacité d'un condensateur plan augmente lorsqu'on écarte ses armatures.
\item En régime permanent continu, un condensateur complètement chargé se comporte comme un interrupteur ouvert.
\item Deux condensateurs identiques montés en série ont une capacité équivalente double de celle de chacun.
\item À l'instant $t = \tau$ lors de la charge d'un condensateur, la tension $u_C$ atteint environ $63\,\%$ de sa valeur finale.
\item L'énergie dissipée par effet Joule dans le résistor pendant la décharge complète d'un condensateur est égale à l'énergie initialement stockée dans le condensateur.
\end{enumerate}
\end{exercice}

\begin{exercice}[Définitions et schéma]
\begin{enumerate}
\item Définis un condensateur et donne son symbole normalisé.
\item Réalise le schéma d'un circuit permettant, grâce à un interrupteur à deux positions, de charger un condensateur avec un générateur de tension continue $E$ à travers un résistor $R$, puis de le décharger dans ce même résistor. Flèche les tensions $u_C$ et $u_R$ en convention récepteur.
\item Donne l'expression de la capacité $C$ d'un condensateur plan en précisant le nom et l'unité SI de chaque grandeur.
\end{enumerate}
\end{exercice}

\begin{exercice}[Calculs directs]
\begin{enumerate}
\item Un condensateur de capacité $C = \SI{100}{\micro F}$ est soumis à une tension $u = \SI{12}{V}$. Calcule la charge $q$ portée par son armature positive et l'énergie emmagasinée.
\item Un condensateur porte une charge $q = \SI{4,8e-4}{C}$ lorsque la tension à ses bornes vaut $u = \SI{24}{V}$. Calcule sa capacité.
\item Un dipôle RC est constitué d'un résistor de résistance $R = \SI{2,2}{k\Omega}$ et d'un condensateur de capacité $C = \SI{47}{\micro F}$. Calcule la constante de temps $\tau$ et la durée approximative de la charge quasi totale (on prendra $5\tau$).
\end{enumerate}
\end{exercice}

\courssub{Je m'entraîne}

\begin{exercice}[Condensateur plan au mica]
Un condensateur plan est formé de deux armatures circulaires de rayon $r = \SI{10}{cm}$, séparées par une feuille de mica d'épaisseur $e = \SI{0,20}{mm}$ et de permittivité relative $\varepsilon_r = 7$.
On donne : $\varepsilon_0 = \SI{8,85e-12}{F.m^{-1}}$.
\begin{enumerate}
\item Calcule l'aire $S$ d'une armature.
\item Calcule la capacité $C$ de ce condensateur.
\item On le charge sous une tension $U = \SI{24}{V}$. Calcule la charge $q$ emmagasinée et l'énergie $E$ stockée.
\end{enumerate}
\end{exercice}

\begin{exercice}[Association mixte de condensateurs]
Deux condensateurs de capacités $C_1 = \SI{4}{\micro F}$ et $C_2 = \SI{12}{\micro F}$ sont montés en parallèle ; l'ensemble est branché en série avec un condensateur $C_3 = \SI{8}{\micro F}$. Le tout est alimenté sous une tension $U = \SI{12}{V}$.
\begin{enumerate}
\item Réalise le schéma du montage.
\item Calcule la capacité $C'$ du groupement parallèle, puis la capacité équivalente $C_{\text{éq}}$ de l'association complète.
\item Calcule la charge totale $Q$ fournie par le générateur.
\item En déduire la tension $U_3$ aux bornes de $C_3$ et la tension $U'$ aux bornes du groupement parallèle, puis la charge de $C_1$ et celle de $C_2$.
\end{enumerate}
\end{exercice}

\begin{exercice}[Décharge d'un condensateur]
Un condensateur de capacité $C = \SI{470}{\micro F}$, chargé sous la tension $U_0 = \SI{15}{V}$, se décharge à travers un résistor de résistance $R = \SI{4,7}{k\Omega}$.
\begin{enumerate}
\item Établis l'équation différentielle vérifiée par $u_C(t)$ pendant la décharge et donne sa solution.
\item Calcule la constante de temps $\tau$.
\item Calcule la tension $u_C$ aux instants $t = \tau$ et $t = 5\tau$. Conclus.
\item Calcule l'énergie initialement stockée dans le condensateur. Que devient cette énergie au bout d'une décharge complète ?
\end{enumerate}
\end{exercice}

\begin{exercice}[Série ou parallèle : le match énergétique]
On dispose de deux condensateurs identiques de capacité $C = \SI{220}{\micro F}$ chacun, et d'un générateur de tension $E = \SI{12}{V}$.
\begin{enumerate}
\item Les deux condensateurs sont montés en série puis chargés sous $E$. Calcule la capacité équivalente, la charge totale et l'énergie totale stockée $E_{\text{série}}$.
\item Mêmes questions avec les deux condensateurs montés en parallèle : calcule $E_{\text{par}}$.
\item Compare $E_{\text{par}}$ et $E_{\text{série}}$ et conclus sur le montage le plus « énergétique ». Ce résultat dépend-il de la valeur de $C$ ?
\end{enumerate}
\end{exercice}

\courssub{Je me prépare au Bac}

\begin{exercice}[Type Bac : étude complète de la charge d'un dipôle RC]
Au laboratoire du lycée, un groupe d'élèves de Terminale C étudie la charge d'un condensateur de capacité $C$ inconnue à travers un résistor de résistance $R = \SI{1,0}{k\Omega}$, alimenté par un générateur de force électromotrice $E = \SI{6,0}{V}$. Un système d'acquisition enregistre la tension $u_C(t)$ aux bornes du condensateur. La courbe obtenue part de l'origine, croît et tend vers une asymptote horizontale $u_C = E$ ; la tangente à l'origine coupe cette asymptote au point d'abscisse $t_0 = \SI{0,22}{s}$.
\begin{enumerate}
\item Réalise le schéma du circuit de charge en fléchant $u_C$ et $u_R$ en convention récepteur.
\item En appliquant la loi d'additivité des tensions et la loi d'Ohm, établis l'équation différentielle vérifiée par $u_C(t)$ :
\[ R C \dfrac{d u_C}{dt} + u_C = E. \]
\item Vérifie que la fonction $u_C(t) = E\left(1 - e^{-t/\tau}\right)$, avec $\tau = RC$, est bien solution de cette équation différentielle et qu'elle satisfait la condition initiale $u_C(0) = 0$.
\item À l'aide de la construction de la tangente à l'origine décrite ci-dessus, détermine graphiquement la constante de temps $\tau$. Justifie la méthode.
\item En déduire la valeur de la capacité $C$ du condensateur.
\item Calcule l'intensité du courant à l'instant $t = 0$, puis l'énergie maximale $E_{\max}$ emmagasinée par le condensateur en fin de charge.
\item Quelle est, en fin de charge, l'énergie totale fournie par le générateur ? Compare-la à $E_{\max}$ et interprète physiquement la différence.
\end{enumerate}
\end{exercice}

\begin{exercice}[Type Bac : le flash de l'appareil photo]
Le flash électronique d'un appareil photo utilise un condensateur de capacité $C = \SI{1,0}{mF}$ chargé sous une tension $U_0 = \SI{300}{V}$. Au moment de la prise de vue, le condensateur se décharge à travers la lampe au xénon, assimilable à un résistor de résistance $R = \SI{2,0}{\Omega}$.
\begin{enumerate}
\item Calcule la charge $Q_0$ et l'énergie $E_0$ emmagasinées par le condensateur avant le déclenchement du flash.
\item Établis l'équation différentielle de la décharge et donne l'expression de $u_C(t)$.
\item Calcule la constante de temps $\tau$ du circuit de décharge. Justifie pourquoi le flash produit un éclair très bref.
\item Calcule l'intensité $i(0)$ du courant au début de la décharge. Commenter cette valeur.
\item On considère que l'éclair est terminé lorsque la tension ne vaut plus que $1\,\%$ de $U_0$. Détermine l'instant $t_1$ correspondant en fonction de $\tau$, puis sa valeur numérique.
\item Quelle énergie a alors été convertie dans la lampe ? Sous quelles formes cette énergie est-elle transférée ?
\end{enumerate}
\end{exercice}

\begin{exercice}[Type Bac : supercondensateur pour l'éclairage de secours]
Dans un quartier de Douala où les coupures d'électricité sont fréquentes, un technicien étudie un système d'éclairage de secours utilisant un supercondensateur de capacité $C = \SI{10}{F}$, chargé sous $U = \SI{12}{V}$ par un panneau solaire. En cas de coupure, le supercondensateur alimente une lampe à LED assimilée à un résistor de résistance $R = \SI{24}{\Omega}$.
\begin{enumerate}
\item Calcule l'énergie $E$ emmagasinée par le supercondensateur chargé. Compare cette valeur à l'énergie stockée par un condensateur classique de $\SI{470}{\micro F}$ chargé sous la même tension.
\item Le technicien hésite entre un seul supercondensateur de $\SI{10}{F}$ et deux supercondensateurs de $\SI{5,0}{F}$ montés en parallèle. Montre que les deux solutions stockent la même énergie sous $\SI{12}{V}$.
\item Pourquoi ne peut-il pas monter les deux supercondensateurs de $\SI{5,0}{F}$ en série pour alimenter la lampe sous $\SI{12}{V}$ ? Calcule la capacité équivalente de ce montage série et l'énergie qu'il stockerait sous $\SI{12}{V}$.
\item Le supercondensateur de $\SI{10}{F}$ se décharge dans la lampe. Calcule la constante de temps $\tau$ du circuit.
\item Écris l'expression de $u_C(t)$ pendant la décharge et calcule la tension aux bornes de la lampe après une durée $t = \SI{4,0}{min}$.
\item La lampe cesse d'éclairer correctement lorsque la tension descend sous $U_{\min} = \SI{6,0}{V}$. Détermine la durée d'autonomie $t_a$ de l'éclairage de secours.
\end{enumerate}
\end{exercice}



\newpage

\coursec{Corrigés des exercices}

\begin{corrige}[Exercice 1 --- QCM : réflexes sur le condensateur]
\begin{enumerate}
\item \textbf{b.} $q = C\,u$ : c'est la relation fondamentale du condensateur, la charge portée par une armature étant proportionnelle à la tension appliquée entre les armatures.
\item \textbf{c.} De $C = \dfrac{q}{u}$, on tire $\SI{1}{F} = \SI{1}{C.V^{-1}}$ : un condensateur de \SI{1}{F} porte une charge de \SI{1}{C} sous une tension de \SI{1}{V}.
\item \textbf{b.} En parallèle, la tension est commune aux deux condensateurs et les charges s'ajoutent : $q = q_1 + q_2 = (C_1 + C_2)\,u$, donc $C_{\text{éq}} = C_1 + C_2$.
\item \textbf{c.} $\tau = RC$ ; vérification dimensionnelle : $\Omega \cdot \mathrm{F} = \dfrac{V}{A}\times\dfrac{C}{V} = \dfrac{C}{A} = \dfrac{A\cdot s}{A} = \mathrm{s}$, homogène à un temps.
\item \textbf{b.} $E = \dfrac{1}{2}Cu^2$, équivalente à $\dfrac{q^2}{2C}$ puisque $q = Cu$ ; la réponse d vaut le double de l'énergie emmagasinée.
\end{enumerate}
\end{corrige}

\begin{corrige}[Exercice 2 --- Vrai ou faux, avec justification]
\begin{enumerate}
\item \textbf{Faux.} $C = \dfrac{\varepsilon S}{e}$ : la capacité est inversement proportionnelle à l'écartement $e$ ; écarter les armatures fait \emph{diminuer} $C$.
\item \textbf{Vrai.} En fin de charge, $u_C = E$ est constante, donc $i = C\dfrac{du_C}{dt} = 0$ : plus aucun courant ne circule, le condensateur se comporte comme un interrupteur ouvert.
\item \textbf{Faux.} En série, $\dfrac{1}{C_{\text{éq}}} = \dfrac{1}{C} + \dfrac{1}{C} = \dfrac{2}{C}$, donc $C_{\text{éq}} = \dfrac{C}{2}$ : la capacité est \emph{divisée par deux}.
\item \textbf{Vrai.} $u_C(\tau) = E(1 - e^{-1}) \approx 0,632\,E$, soit environ $63\,\%$ de la valeur finale.
\item \textbf{Vrai.} Par conservation de l'énergie, l'énergie $\dfrac{1}{2}CU_0^2$ initialement stockée est intégralement dissipée par effet Joule dans le résistor au cours de la décharge complète.
\end{enumerate}
\end{corrige}

\begin{corrige}[Exercice 3 --- Définitions et schéma]
\begin{enumerate}
\item Un \cle{condensateur} est un dipôle constitué de deux conducteurs (les \cle{armatures}) séparés par un isolant (le \cle{diélectrique}) ; il peut emmagasiner des charges électriques opposées $+q$ et $-q$ sur ses armatures. Symbole normalisé : deux traits parallèles de même longueur, perpendiculaires aux fils de connexion.
\item Schéma du circuit charge--décharge :
\begin{popfigure}
\begin{tikzpicture}[european, scale=1.0]
\draw (0,0) to[battery1, l=$E$] (0,2.5);
\draw (0,2.5) -- (1.2,2.5);
\draw (1.2,2.5) node[above left]{1} -- (2.2,2.7);
\draw (2.2,2.5) node[above right]{2} (2.4,2.5) -- (3.4,2.5);
\draw (1.2,2.5) node[circ]{} (2.4,2.5) node[circ]{};
\draw (3.4,2.5) to[R, R=$R$, v=$u_R$] (6,2.5);
\draw (6,2.5) to[C, C=$C$, v=$u_C$] (6,0);
\draw (6,0) -- (0,0);
\draw (2.2,2.7) node[above]{$K$};
\end{tikzpicture}
\legende{Interrupteur $K$ en position 1 : charge du condensateur ; en position 2 : décharge à travers $R$.}
\end{popfigure}
\item Pour un condensateur plan : $C = \dfrac{\varepsilon\, S}{e}$, avec $\varepsilon$ la permittivité du diélectrique ($\si{F.m^{-1}}$), $S$ l'aire en regard des armatures ($\si{m^2}$) et $e$ l'écartement entre armatures ($\si{m}$). Pour un diélectrique de permittivité relative $\varepsilon_r$ : $\varepsilon = \varepsilon_r\,\varepsilon_0$ avec $\varepsilon_0 = \SI{8,85e-12}{F.m^{-1}}$.
\end{enumerate}
\end{corrige}

\begin{corrige}[Exercice 4 --- Calculs directs]
\begin{enumerate}
\item $q = Cu = \num{100e-6} \times 12 = \SI{1,2e-3}{C} = \SI{1,2}{mC}$.\\
$E = \dfrac{1}{2}Cu^2 = \dfrac{1}{2}\times \num{100e-6} \times 12^2 = \SI{7,2e-3}{J} = \SI{7,2}{mJ}$.
\item $C = \dfrac{q}{u} = \dfrac{\num{4,8e-4}}{24} = \SI{2,0e-5}{F} = \SI{20}{\micro F}$.
\item $\tau = RC = \num{2,2e3} \times \num{47e-6} = \SI{0,103}{s} \approx \SI{0,10}{s}$.\\
Durée de charge quasi totale : $5\tau \approx \SI{0,52}{s}$.
\end{enumerate}
\end{corrige}

\begin{corrige}[Exercice 5 --- Condensateur plan au mica]
\begin{enumerate}
\item $S = \pi r^2 = \pi \times (0,10)^2 = \SI{3,14e-2}{m^2}$.
\item $C = \dfrac{\varepsilon_0 \varepsilon_r S}{e} = \dfrac{\num{8,85e-12}\times 7 \times \num{3,14e-2}}{\num{2,0e-4}} = \SI{9,7e-9}{F} \approx \SI{9,7}{nF}$.
\item $q = CU = \num{9,7e-9}\times 24 = \SI{2,3e-7}{C}$.\\
$E = \dfrac{1}{2}CU^2 = \dfrac{1}{2}\times \num{9,7e-9}\times 24^2 = \SI{2,8e-6}{J} = \SI{2,8}{\micro J}$.
\end{enumerate}
\end{corrige}

\begin{corrige}[Exercice 6 --- Association mixte de condensateurs]
\begin{enumerate}
\item Schéma du montage :
\begin{popfigure}
\begin{tikzpicture}[european, scale=1.0]
\draw (0,0) to[battery1, l=$U$] (0,3);
\draw (0,3) -- (1,3) to[C, l=$C_3$] (3,3) -- (4,3);
\draw (4,3) -- (4,4) to[C, l=$C_1$] (6.5,4) -- (6.5,3);
\draw (4,3) -- (4,2) to[C, l_=$C_2$] (6.5,2) -- (6.5,3);
\draw (6.5,3) -- (7.5,3) -- (7.5,0) -- (0,0);
\end{tikzpicture}
\legende{$C_1$ et $C_2$ en parallèle, le groupement en série avec $C_3$.}
\end{popfigure}
\item Groupement parallèle : $C' = C_1 + C_2 = 4 + 12 = \SI{16}{\micro F}$.\\
Association série de $C'$ et $C_3$ : $C_{\text{éq}} = \dfrac{C'\,C_3}{C' + C_3} = \dfrac{16\times 8}{16+8} = \SI{5,3}{\micro F}$.
\item $Q = C_{\text{éq}}\,U = \num{5,33e-6}\times 12 = \SI{6,4e-5}{C} = \SI{64}{\micro C}$.
\item En série, $C_3$ porte la charge totale : $U_3 = \dfrac{Q}{C_3} = \dfrac{64}{8} = \SI{8,0}{V}$.\\
Par additivité : $U' = U - U_3 = 12 - 8 = \SI{4,0}{V}$.\\
$q_1 = C_1 U' = 4\times 4 = \SI{16}{\micro C}$ ; \quad $q_2 = C_2 U' = 12\times 4 = \SI{48}{\micro C}$.\\
Vérification : $q_1 + q_2 = 64\ \mu\mathrm{C} = Q$.
\end{enumerate}
\end{corrige}

\begin{corrige}[Exercice 7 --- Décharge d'un condensateur]
\begin{enumerate}
\item Pendant la décharge, le générateur est débranché : la loi d'additivité des tensions donne $u_C + u_R = 0$, avec $u_R = Ri$ et $i = \dfrac{dq}{dt} = C\dfrac{du_C}{dt}$ (convention récepteur). D'où :
\[ RC\,\dfrac{du_C}{dt} + u_C = 0. \]
Solution satisfaisant la condition initiale $u_C(0) = U_0$ : $u_C(t) = U_0\,e^{-t/\tau}$ avec $\tau = RC$.
\item $\tau = RC = \num{4,7e3}\times \num{470e-6} = \SI{2,2}{s}$.
\item $u_C(\tau) = 15\,e^{-1} = \SI{5,5}{V}$ (environ $37\,\%$ de $U_0$).\\
$u_C(5\tau) = 15\,e^{-5} = \SI{0,10}{V}$ (moins de $1\,\%$ de $U_0$) : la décharge est quasi totale au bout de $5\tau \approx \SI{11}{s}$.
\item $E_0 = \dfrac{1}{2}CU_0^2 = \dfrac{1}{2}\times \num{470e-6}\times 15^2 = \SI{5,3e-2}{J}$.\\
Au bout de la décharge complète, cette énergie a été intégralement dissipée par effet Joule dans le résistor (échauffement).
\end{enumerate}
\end{corrige}

\begin{corrige}[Exercice 8 --- Série ou parallèle : le match énergétique]
\begin{enumerate}
\item En série : $C_{\text{éq}} = \dfrac{C}{2} = \SI{110}{\micro F}$.\\
$Q = C_{\text{éq}}E = \num{110e-6}\times 12 = \SI{1,32e-3}{C}$.\\
$E_{\text{série}} = \dfrac{1}{2}C_{\text{éq}}E^2 = \dfrac{1}{2}\times \num{110e-6}\times 12^2 = \SI{7,9e-3}{J}$.
\item En parallèle : $C'_{\text{éq}} = 2C = \SI{440}{\micro F}$ ; $Q' = 2CE = \SI{5,28e-3}{C}$.\\
$E_{\text{par}} = \dfrac{1}{2}(2C)E^2 = \dfrac{1}{2}\times \num{440e-6}\times 144 = \SI{3,2e-2}{J}$.
\item $\dfrac{E_{\text{par}}}{E_{\text{série}}} = \dfrac{2C}{C/2} = 4$ : le montage parallèle stocke \textbf{quatre fois plus} d'énergie. Ce rapport ne dépend pas de $C$ : pour deux condensateurs identiques chargés sous la même tension, le parallèle est toujours quatre fois plus « énergétique » que la série.
\end{enumerate}
\end{corrige}

\begin{corrige}[Exercice 9 --- Type Bac : étude complète de la charge d'un dipôle RC]
\begin{enumerate}
\item Schéma du circuit de charge :
\begin{popfigure}
\begin{tikzpicture}[european, scale=1.0]
\draw (0,0) to[battery1, l=$E$] (0,3);
\draw (0,3) to[nos, l=$K$] (2,3);
\draw (2,3) to[R, l=$R$, v=$u_R$] (5,3);
\draw (5,3) to[C, l=$C$, v=$u_C$] (5,0);
\draw (5,0) -- (0,0);
\end{tikzpicture}
\legende{Circuit de charge du dipôle RC série.}
\end{popfigure}
\item Loi d'additivité des tensions : $E = u_R + u_C$. Loi d'Ohm : $u_R = Ri$, avec $i = \dfrac{dq}{dt} = C\dfrac{du_C}{dt}$. Donc :
\[ E = RC\,\dfrac{du_C}{dt} + u_C. \]
\item Dérivons la solution proposée $u_C(t) = E\left(1 - e^{-t/\tau}\right)$ : $\dfrac{du_C}{dt} = \dfrac{E}{\tau}e^{-t/\tau}$. Alors, en tenant compte de $\tau = RC$ :
\[ RC\,\dfrac{du_C}{dt} + u_C = \tau\cdot\dfrac{E}{\tau}e^{-t/\tau} + E - Ee^{-t/\tau} = E. \]
L'équation différentielle est vérifiée. De plus $u_C(0) = E(1 - e^{0}) = 0$ : la condition initiale (condensateur initialement déchargé) est satisfaite.
\item La tangente à l'origine a pour équation $u = \left(\dfrac{du_C}{dt}\right)_{t=0} t = \dfrac{E}{\tau}\,t$. Elle coupe l'asymptote $u = E$ en $t = \tau$. Donc $\tau = t_0 = \SI{0,22}{s}$.
\item $C = \dfrac{\tau}{R} = \dfrac{0,22}{\num{1,0e3}} = \SI{2,2e-4}{F} = \SI{220}{\micro F}$.
\item À $t = 0$, $u_C = 0$ donc $u_R = E$ : $i(0) = \dfrac{E}{R} = \dfrac{6,0}{\num{1,0e3}} = \SI{6,0e-3}{A} = \SI{6,0}{mA}$.\\
$E_{\max} = \dfrac{1}{2}CE^2 = \dfrac{1}{2}\times \num{2,2e-4}\times 6,0^2 = \SI{4,0e-3}{J} = \SI{4,0}{mJ}$.
\item La charge totale ayant traversé le générateur est $Q = CE$ ; l'énergie fournie est $E_{\text{gén}} = QE = CE^2 = 2E_{\max} = \SI{7,9}{mJ}$. La moitié de l'énergie fournie est stockée dans le condensateur, l'autre moitié est dissipée par effet Joule dans le résistor pendant la charge.
\end{enumerate}
\textbf{Barème indicatif (sur 8 points) :} schéma 1 pt ; équation différentielle 1,5 pt ; vérification de la solution 1,5 pt ; méthode de la tangente 1 pt ; $C$ 1 pt ; $i(0)$ et $E_{\max}$ 1 pt ; bilan énergétique 1 pt.\\
\textbf{Points de vigilance :} citer explicitement la convention récepteur et la relation $i = C\dfrac{du_C}{dt}$ ; pour la question 3, dériver puis \emph{substituer} dans l'équation (ne pas se contenter d'affirmer) ; justifier la méthode graphique par le calcul de la pente à l'origine ; ne pas oublier les unités (farad, seconde, joule).
\end{corrige}

\begin{corrige}[Exercice 10 --- Type Bac : le flash de l'appareil photo]
\begin{enumerate}
\item $Q_0 = CU_0 = \num{1,0e-3}\times 300 = \SI{0,30}{C}$.\\
$E_0 = \dfrac{1}{2}CU_0^2 = \dfrac{1}{2}\times \num{1,0e-3}\times 300^2 = \SI{45}{J}$.
\item Loi d'additivité : $u_C + u_R = 0$ avec $u_R = Ri = RC\dfrac{du_C}{dt}$, d'où :
\[ RC\,\dfrac{du_C}{dt} + u_C = 0, \qquad u_C(t) = U_0\,e^{-t/\tau}. \]
\item $\tau = RC = 2,0\times \num{1,0e-3} = \SI{2,0e-3}{s} = \SI{2,0}{ms}$. La constante de temps est de l'ordre de la milliseconde : la décharge est quasi totale en $5\tau \approx \SI{10}{ms}$, d'où un éclair extrêmement bref.
\item $i(0) = \dfrac{U_0}{R} = \dfrac{300}{2,0} = \SI{150}{A}$. Cette intensité est énorme (comparée aux quelques ampères d'un appareil domestique), mais elle ne dure que quelques millisecondes : c'est ce pic de courant qui porte le gaz xénon à incandescence.
\item $u_C(t_1) = 0,01\,U_0 = U_0 e^{-t_1/\tau}$ donne $e^{-t_1/\tau} = 0,01$, soit $t_1 = \tau\ln 100 \approx 4,6\,\tau$. Numériquement : $t_1 = 4,6\times \num{2,0e-3} = \SI{9,2e-3}{s} \approx \SI{9,2}{ms}$.
\item Énergie restante dans le condensateur : $E(t_1) = \dfrac{1}{2}C(0,01\,U_0)^2 = 10^{-4}\,E_0$, négligeable. L'énergie convertie dans la lampe vaut donc pratiquement $E_0 = \SI{45}{J}$. Elle est transférée sous forme de \cle{rayonnement lumineux} (l'éclair) et de \cle{chaleur} (échauffement de la lampe).
\end{enumerate}
\textbf{Barème indicatif (sur 7 points) :} $Q_0$ et $E_0$ 1 pt ; équation différentielle et solution 1,5 pt ; $\tau$ et justification de la brièveté 1 pt ; $i(0)$ et commentaire 1 pt ; $t_1$ (littéral puis numérique) 1,5 pt ; bilan énergétique 1 pt.\\
\textbf{Points de vigilance :} pour la décharge, la loi des mailles s'écrit $u_C + u_R = 0$ (le générateur est débranché) ; l'extraction de $t_1$ exige le passage au logarithme népérien, à rédiger proprement ; penser à convertir les millisecondes ; le commentaire physique sur le pic de courant est attendu.
\end{corrige}

\begin{corrige}[Exercice 11 --- Type Bac : supercondensateur pour l'éclairage de secours]
\begin{enumerate}
\item $E = \dfrac{1}{2}CU^2 = \dfrac{1}{2}\times 10\times 12^2 = \SI{720}{J}$.\\
Condensateur classique : $E' = \dfrac{1}{2}\times \num{470e-6}\times 144 = \SI{3,4e-2}{J}$.\\
Rapport : $\dfrac{E}{E'} \approx \num{2,1e4}$ : le supercondensateur stocke environ \num{21000} fois plus d'énergie.
\item Deux condensateurs de $\SI{5,0}{F}$ en parallèle : $C_{\text{éq}} = 5,0 + 5,0 = \SI{10}{F}$, identique au supercondensateur unique. Sous la même tension $U = \SI{12}{V}$, $E = \dfrac{1}{2}C_{\text{éq}}U^2 = \SI{720}{J}$ : même énergie stockée.
\item En série : $C_{\text{éq}} = \dfrac{5,0}{2} = \SI{2,5}{F}$. L'énergie stockée sous $\SI{12}{V}$ serait $E = \dfrac{1}{2}\times 2,5\times 144 = \SI{180}{J}$, soit quatre fois moins. De plus, chaque supercondensateur supporterait $\SI{6}{V}$, ce qui peut dépasser sa tension maximale d'usage (souvent $\SI{2,7}{V}$ pour un supercondensateur élémentaire) : le montage série est à proscrire ici.
\item $\tau = RC = 24\times 10 = \SI{240}{s} = \SI{4,0}{min}$.
\item $u_C(t) = U\,e^{-t/\tau} = 12\,e^{-t/240}$ (avec $t$ en secondes).\\
À $t = \SI{4,0}{min} = \SI{240}{s} = \tau$ : $u_C = 12\,e^{-1} = \SI{4,4}{V}$.
\item $U_{\min} = U\,e^{-t_a/\tau}$ donne $e^{-t_a/\tau} = \dfrac{6,0}{12} = 0,5$, soit $t_a = \tau\ln 2 = 240\times 0,693 = \SI{166}{s} \approx \SI{2,8}{min}$. L'éclairage de secours fonctionne correctement pendant environ $2$ minutes $46$ secondes.
\end{enumerate}
\textbf{Barème indicatif (sur 7 points) :} énergies et comparaison 1,5 pt ; équivalence parallèle 1 pt ; analyse du montage série 1,5 pt ; $\tau$ 0,5 pt ; $u_C(t)$ et valeur à $t = \SI{4}{min}$ 1 pt ; autonomie $t_a$ 1,5 pt.\\
\textbf{Points de vigilance :} attention à la cohérence des unités de temps (convertir les minutes en secondes avant d'utiliser l'exponentielle) ; le calcul de $t_a$ exige $\ln 2 \approx 0,693$ ; la justification du refus du montage série doit invoquer à la fois l'énergie divisée par quatre et la tension maximale admissible par composant.
\end{corrige}

\newpage

\coursec{Fiche bilan}

\courssub{Carte mentale de la leçon}

\begin{popfigure}
\begin{center}\tcbox[colback=popdark,colframe=popdark,coltext=white,arc=9pt,boxsep=4pt,left=6pt,right=6pt]{\bfseries\small Le condensateur}\end{center}
\vspace{-2pt}
\begin{tcbraster}[raster columns=3,raster equal height=rows,raster column skip=6pt,raster row skip=6pt,enhanced,blankest]
\begin{tcolorbox}[enhanced,colback=popblue,colframe=popblue,coltext=white,boxrule=0.9pt,arc=3pt,left=4pt,right=4pt,top=3pt,bottom=3pt,boxsep=1pt,halign=left]\bfseries\scriptsize Constitution, charge $q = C\,u$\end{tcolorbox}
\begin{tcolorbox}[enhanced,colback=poporange,colframe=poporange,coltext=white,boxrule=0.9pt,arc=3pt,left=4pt,right=4pt,top=3pt,bottom=3pt,boxsep=1pt,halign=left]\bfseries\scriptsize Condensateur plan $C=\dfrac{\varepsilon S}{e}$\end{tcolorbox}
\begin{tcolorbox}[enhanced,colback=popgreen,colframe=popgreen,coltext=white,boxrule=0.9pt,arc=3pt,left=4pt,right=4pt,top=3pt,bottom=3pt,boxsep=1pt,halign=left]\bfseries\scriptsize Associations série / parallèle\end{tcolorbox}
\begin{tcolorbox}[enhanced,colback=poppurple,colframe=poppurple,coltext=white,boxrule=0.9pt,arc=3pt,left=4pt,right=4pt,top=3pt,bottom=3pt,boxsep=1pt,halign=left]\bfseries\scriptsize Charge (dipôle RC) $u_C=E(1-e^{-t/\tau})$\end{tcolorbox}
\begin{tcolorbox}[enhanced,colback=poppink,colframe=poppink,coltext=white,boxrule=0.9pt,arc=3pt,left=4pt,right=4pt,top=3pt,bottom=3pt,boxsep=1pt,halign=left]\bfseries\scriptsize Décharge $u_C=E\,e^{-t/\tau}$\end{tcolorbox}
\begin{tcolorbox}[enhanced,colback=popgold,colframe=popgold,coltext=white,boxrule=0.9pt,arc=3pt,left=4pt,right=4pt,top=3pt,bottom=3pt,boxsep=1pt,halign=left]\bfseries\scriptsize Énergie $E_e=\tfrac12 C u^2$\end{tcolorbox}
\end{tcbraster}

\legende{Carte mentale : les six idées-forces de la leçon sur le condensateur.}
\end{popfigure}

\begin{definition}[Définitions clés]
\begin{itemize}
  \item \cle{Condensateur} : dipôle formé de deux \cle{armatures} conductrices séparées par un \cle{isolant} (diélectrique). Il emmagasine des charges $+q$ et $-q$ sur ses armatures.
  \item \cle{Charge} $q$ portée par l'armature reliée à l'entrée du courant (convention récepteur) : $q = C\,u$, avec $C$ la \cle{capacité} en farads (F).
  \item \cle{Constante de temps} du dipôle RC : $\tau = RC$ ; elle mesure la rapidité de la charge ou de la décharge (homogénéité : $\Omega \cdot \text{F} = \text{s}$).
\end{itemize}
\end{definition}

\begin{propriete}[Formules à connaître par cœur]
\begin{center}
\begin{tabularx}{\linewidth}{|l|X|l|}
\hline
\rowcolor{popblueL} \textbf{Grandeur} & \textbf{Relation} & \textbf{Unités SI} \\
\hline
Charge--tension & $q = C\,u$ & C, F, V \\
\hline
Intensité (conv. récepteur) & $i = \dfrac{dq}{dt} = C\,\dfrac{du}{dt}$ & A \\
\hline
Condensateur plan & $C = \dfrac{\varepsilon\, S}{e}$ & F \\
\hline
Association parallèle & $C_{\text{éq}} = C_1 + C_2 + \cdots$ & F \\
\hline
Association série & $\dfrac{1}{C_{\text{éq}}} = \dfrac{1}{C_1} + \dfrac{1}{C_2} + \cdots$ & F \\
\hline
Constante de temps & $\tau = RC$ & s \\
\hline
Charge (dipôle RC série) & $u_C(t) = E\left(1 - e^{-t/\tau}\right)$ & V \\
\hline
Décharge (dipôle RC série) & $u_C(t) = E\,e^{-t/\tau}$ & V \\
\hline
Énergie emmagasinée & $E_e = \dfrac{1}{2}C u^2 = \dfrac{q^2}{2C} = \dfrac{1}{2}q\,u$ & J \\
\hline
\end{tabularx}
\end{center}
\end{propriete}

\begin{propriete}[Équations différentielles du dipôle RC série]
\begin{itemize}
  \item \textbf{Charge} (générateur de tension $E$, interrupteur fermé sur $E$ à $t=0$) : 
  \[ RC\,\dfrac{du_C}{dt} + u_C = E \qquad\text{avec } u_C(0)=0. \]
  \item \textbf{Décharge} (générateur court-circuité ou déconnecté, condensateur initialement chargé à $E$) : 
  \[ RC\,\dfrac{du_C}{dt} + u_C = 0 \qquad\text{avec } u_C(0)=E. \]
  \item \textbf{Vérification d'une solution proposée} : on la dérive, on la reporte dans l'équation différentielle, on identifie les termes constants et les termes en $e^{-t/\tau}$.
\end{itemize}
\end{propriete}

\begin{methode}[Méthodes express]
\begin{enumerate}
  \item \cle{Déterminer $\tau$ graphiquement} : tracer la tangente à la courbe $u_C(t)$ à l'origine ($t=0$) ; elle coupe l'asymptote ($u_C = E$ en charge, $u_C = 0$ en décharge) à l'abscisse $t = \tau$. Sinon, lire $t$ tel que $u_C = 0{,}63\,E$ (charge) ou $u_C = 0{,}37\,E$ (décharge).
  \item \cle{Régime permanent} : en charge terminée ($t \gg 5\tau$), le condensateur se comporte comme un interrupteur ouvert : $i = 0$ et $u_C = E$.
  \item \cle{Associations} : parallèle $\Rightarrow$ les capacités s'ajoutent ($C_{\text{éq}} = \sum C_k$) ; série $\Rightarrow$ les inverses s'ajoutent ($1/C_{\text{éq}} = \sum 1/C_k$) — l'inverse des règles pour les résistances !
  \item \cle{Continuité de la tension} : la tension $u_C$ aux bornes d'un condensateur ne subit jamais de discontinuité ; c'est l'intensité $i$ qui peut sauter (ex. : fermeture d'un interrupteur).
\end{enumerate}
\end{methode}

\begin{aretenir}[Checklist avant le Bac]
\begin{itemize}
  \item[$\square$] Je sais dessiner le symbole normalisé du condensateur et flécher $u_C$ et $i$ en convention récepteur.
  \item[$\square$] Je connais la relation $q = C\,u$ et je sais que $i = C\,\dfrac{du_C}{dt}$ (convention récepteur).
  \item[$\square$] Je sais que la capacité d'un condensateur plan est $C = \dfrac{\varepsilon S}{e}$ : $C$ augmente si $S$ augmente ou si $e$ diminue.
  \item[$\square$] Je ne confonds pas les règles d'association : parallèle $\Rightarrow$ somme des capacités ; série $\Rightarrow$ somme des inverses.
  \item[$\square$] Je sais établir l'équation différentielle de la charge ou de la décharge à partir de la loi d'additivité des tensions (loi des mailles).
  \item[$\square$] Je connais les solutions $u_C = E(1 - e^{-t/\tau})$ (charge) et $u_C = E\,e^{-t/\tau}$ (décharge) avec $\tau = RC$.
  \item[$\square$] Je sais déterminer $\tau$ graphiquement (tangente à l'origine, ou points à $63\,\%$ et $37\,\%$ de la variation totale).
  \item[$\square$] Je sais qu'un condensateur chargé isolé conserve sa charge, et qu'en régime permanent il bloque le courant continu ($i=0$).
  \item[$\square$] Je maîtrise les trois expressions de l'énergie : $E_e = \tfrac12 C u^2 = \dfrac{q^2}{2C} = \tfrac12 q u$.
  \item[$\square$] Je vérifie toujours l'homogénéité : $RC$ en secondes, $Cu^2$ en joules.
  \item[$\square$] Je pense que la tension $u_C$ est continue à $t = 0$ : $u_C(0^+) = u_C(0^-)$.
  \item[$\square$] Je donne chaque résultat numérique avec son unité SI et un nombre cohérent de chiffres significatifs.
\end{itemize}
\end{aretenir}

\findecours

\end{document}
